% !TeX program = pdflatex
% !TeX encoding = UTF-8
%==============================================================================
%  REFERAT
%  Semantyka denotacyjna i równania stałopunktowe
%  (na podstawie: M. Gordon, "Denotacyjny opis języków programowania";
%   Harvard CS152 Lecture 6; Cornell CS4110 Lecture 8)
%
%  Kompilacja:  pdflatex referat.tex   (dwa przebiegi --- spis treści i odsyłacze)
%  Styl i środowiska analogiczne do pliku rownania_stalopunktowe.tex.
%==============================================================================
\documentclass[11pt,a4paper]{article}

%----------------------------- Kodowanie i język ------------------------------
\usepackage[T1]{fontenc}
\usepackage[utf8]{inputenc}
\usepackage[polish]{babel}
\usepackage{lmodern}
\usepackage{enumitem}
\usepackage{float}

\usepackage{tikz}
\usetikzlibrary{positioning}
\usetikzlibrary{decorations.pathreplacing}
\usetikzlibrary{calc,shapes.geometric}

%----------------------------- Matematyka -------------------------------------
\usepackage{amsmath,amssymb,amsthm,mathtools}
\usepackage{stmaryrd}   % \llbracket \rrbracket  (nawiasy semantyczne)
\usepackage{mathrsfs}   % \mathscr  (kaligrafia funkcji semantycznych)

%----------------------------- Układ i odsyłacze ------------------------------
\usepackage[margin=2.5cm]{geometry}
\usepackage{xcolor}
\usepackage{hyperref}
\hypersetup{hypertexnames=false, colorlinks=true, linkcolor=blue!50!black, citecolor=blue!50!black}
\usepackage{comment}


%----------------------------- Środowiska twierdzeń ---------------------------
% Każdy typ ma własny licznik numerowany w obrębie sekcji (np. Definicja 2.1).
\theoremstyle{definition}
\newtheorem{definicja}{Definicja}[section]
\newtheorem{uwaga}[definicja]{Uwaga}
\newtheorem{przyklad}[definicja]{Przykład}
\newtheorem{oznaczenie}[definicja]{Oznaczenie}

\theoremstyle{plain}
\newtheorem{twierdzenie}[definicja]{Twierdzenie}
\newtheorem{lemat}[definicja]{Lemat}
\newtheorem{wniosek}[definicja]{Wniosek}
\newtheorem{fakt}[definicja]{Fakt}

% Nazwa dowodu po polsku
\addto\captionspolish{\renewcommand{\proofname}{Dowód}}

%----------------------------- Makra pomocnicze -------------------------------
\newcommand{\bt}{\bot}                            % element najmniejszy
\newcommand{\Fix}{\operatorname{Fix}}            % zbiór wszystkich punktów stałych
\newcommand{\lfp}{\operatorname{lfp}}
\newcommand{\dom}{\operatorname{dom}}
% Nawiasy semantyczne: \sem{...}
\DeclarePairedDelimiter{\semdelim}{\llbracket}{\rrbracket}
\newcommand{\sem}[1]{\semdelim{#1}}
% Funkcje semantyczne
\newcommand{\Cc}{\mathscr{C}}   % instrukcje (Commands)
\newcommand{\Aa}{\mathscr{A}}   % wyrażenia arytmetyczne (Arithmetic)
\newcommand{\Bb}{\mathscr{B}}   % wyrażenia logiczne (Boolean)
% Zbiory dziedzinowe
\newcommand{\Store}{\textbf{Store}}
\newcommand{\Int}{\textbf{Int}}
\newcommand{\Var}{\textbf{Var}}
\newcommand{\Bool}{\{\kw{true},\kw{false}\}}
\newcommand{\kw}[1]{\ensuremath{\mathsf{#1}}} % slowa kluczowe / literaly IMP

%==============================================================================
\title{Twierdzenia o punktach stałych w programowaniu}
\author{Adrian Madajewski}
\date{\today}

\begin{document}
\maketitle
\tableofcontents
\newpage

%==============================================================================
\section{Wstęp}
%==============================================================================

Semantyka denotacyjna przypisuje konstrukcjom języka programowania obiekty matematyczne opisujące ich znaczenie, abstrahując od sposobu wykonania programu. Dla większości konstrukcji przypisanie to jest bezpośrednie, jednak pętla \kw{while} prowadzi do \emph{równania rekurencyjnego}, w którym znaczenie pętli występuje po obu stronach. Rozwiązania takiego równania poszukujemy jako \emph{najmniejszego punktu stałego} odpowiedniego odwzorowania. Celem niniejszego referatu jest zbudowanie aparatu nadającego temu pojęciu ścisły sens --- teorii zbiorów łańcuchowo zupełnych i funkcji ciągłych --- udowodnienie twierdzenia Kleenego o najmniejszym punkcie stałym oraz zastosowanie go do semantyki pętli \kw{while} w modelowym języku \texttt{IMP}. Prezentacja opiera się na monografii Gordona~\cite{gordon1979} oraz na notatkach do wykładów z~semantyki denotacyjnej~\cite{cs152denot,cs4110denot}.

W dalszej części referatu przedstawione zostanie również alternatywne podejście do problemu punktów stałych, oparte na \emph{metrykach częściowych}. W tym ujęciu pojęcia przybliżenia i granicy nie są opisywane za pomocą porządku, lecz za pomocą odległości, która może również mierzyć stopień określenia elementu. Pozwala to sformułować odpowiednik klasycznego twierdzenia Banacha o punkcie stałym dla zupełnych przestrzeni metrycznych częściowych. Podejście to, związane z pracą Matthewsa~\cite{matthews1994}, zostanie przedstawione jako uzupełnienie podejścia porządkowego. Szczególną uwagę poświęcimy relacji pomiędzy porządkiem indukowanym przez metrykę częściową, ciągami Cauchy'ego oraz kresami górnymi łańcuchów.

Oba podejścia prowadzą do konstrukcji punktów stałych za pomocą iteracji, lecz opierają się na innych założeniach. W przypadku twierdzenia Kleenego kluczową rolę odgrywają zupełność porządkowa oraz ciągłość odwzorowania, natomiast w przypadku twierdzenia Matthewsa --- zupełność metryczna i warunek kontrakcji. Pozwala to spojrzeć na problem punktów stałych z dwóch różnych perspektyw matematycznych oraz wskazać, które elementy konstrukcji są specyficzne dla semantyki denotacyjnej, a które mają charakter bardziej ogólny.


%==============================================================================
\section{Podstawowe pojęcia}
%==============================================================================

Przypomnijmy najpierw klasyczne pojęcie funkcji jako relacji.

\begin{definicja}\label{def:funkcja}
\emph{Funkcją} ze zbioru $A$ w zbiór $B$, oznaczaną $f:A\to B$, nazywamy relację
$f\subseteq A\times B$ spełniającą dwa warunki:
\begin{enumerate}
  \item[(1)] \emph{całkowitość}: dla każdego $a\in A$ istnieje $b\in B$
        takie, że $(a,b)\in f$;
  \item[(2)] \emph{jednoznaczność}: jeśli $(a,b)\in f$ oraz $(a,b')\in f$,
        to $b=b'$.
\end{enumerate}
Ten jedyny element $b$ oznaczamy $f(a)$, tzn.\ $(a,b)\in f \iff b=f(a)$.
Ze względu na warunek~(1) (całkowitość) taką funkcję nazywamy tu \emph{funkcją
całkowitą} --- w odróżnieniu od funkcji częściowej wprowadzonej niżej.
\end{definicja}

\begin{definicja}\label{def:funkcja-czesciowa}
\emph{Funkcją częściową} ze zbioru $A$ w zbiór $B$, oznaczaną $f:A\rightharpoonup B$,
nazywamy relację $f\subseteq A\times B$ spełniającą \emph{tylko} warunek
jednoznaczności:
\begin{enumerate}
  \item[(2)] \emph{jednoznaczność}: jeśli $(a,b)\in f$ oraz $(a,b')\in f$, to $b=b'$.
\end{enumerate}
Warunku~(1) z Definicji~\ref{def:funkcja} (całkowitości) \emph{nie} wymagamy: dla
niektórych $a\in A$ może nie istnieć żadna para $(a,b)\in f$; mówimy wtedy, że
$f(a)$ jest \emph{nieokreślone}, a gdy $(a,b)\in f$, piszemy $b=f(a)$. Funkcję nigdzie
nieokreśloną utożsamiamy z relacją pustą~$\varnothing$.
\end{definicja}

\begin{definicja}\label{def:dziedzina_okreslonosci}
	\emph{Dziedziną określoności} funkcji częściowej $f:A\rightharpoonup B$ nazywamy zbiór tych argumentów, dla których $f$ jest określona, tzn.
	\[
	\dom(f)=\{a\in A:\ \exists b\in B\ (a,b)\in f\}\subseteq A.
	\]
	Funkcja całkowita (Definicja~\ref{def:funkcja}) jest szczególnym przypadkiem, w którym $\dom(f)=A$.
\end{definicja}

\begin{przyklad}\label{prz:poprzednik}
Niech $\mathbb{N}=\{1,2,3,\dots\}$ oraz $\mathbb{N}_0=\mathbb{N}\cup\{0\}$. Funkcja
poprzednika $p:\mathbb{N}_0\rightharpoonup\mathbb{N}_0$, $p(n)=n-1$, jest częściowa:
$\dom(p)=\mathbb{N}=\{1,2,3,\dots\}$, a $p(0)$ jest nieokreślone (argument $0$ należy
do $\mathbb{N}_0$, lecz $0-1\notin\mathbb{N}_0$).
\end{przyklad}

\begin{comment}
	\begin{uwaga}
		Częściowość jest kluczowa. Program może się \emph{nie zatrzymać}. Dla stanu
		wejściowego, z którego obliczenie się zapętla, nie ma stanu wyjściowego --- funkcja
		jest tam nieokreślona.
	\end{uwaga}
\end{comment}

%==============================================================================
\section{Język \texttt{IMP}: składnia i semantyka denotacyjna}
%==============================================================================

Modelowym językiem rozważanym w tym referacie jest \texttt{IMP}; jego składnię
oraz semantykę denotacyjną przedstawiamy za notatkami do wykładów~\cite{cs152denot,cs4110denot}.

\subsection{Składnia języka programowania}

Przed zdefiniowaniem semantyki języka programowania należy określić jego składnię,
czyli zbiór poprawnych programów, wyrażeń oraz innych konstrukcji językowych.
W tym celu składnię można opisać za pomocą gramatyki formalnej, która określa
sposób konstruowania poprawnych wyrażeń danego języka.

W dalszej części będziemy posługiwać się zapisem reguł gramatycznych w notacji Backusa--Naura (BNF, ang. \emph{Backus--Naur Form}), wprowadzonej przez Johna Backusa w kontekście formalnego opisu składni języków programowania i rozwiniętej w opisie języka ALGOL 60 przez Petera Naura \cite{naur1960}. Reguły tej postaci zapisywane są jako
\[
\langle\text{wyrażenie}\rangle ::= \ldots
\]
gdzie symbol $\langle\text{wyrażenie}\rangle$ oznacza kategorię składniową,
natomiast prawa strona reguły określa dopuszczalne sposoby konstruowania
elementów należących do tej kategorii. Symbol $::=$ należy rozumieć jako
„jest definiowane przez” lub „może być postaci”.

Przykładowo, prostą klasę wyrażeń arytmetycznych można zdefiniować za pomocą
reguł
\[
\begin{aligned}
	\langle\text{wyrażenie}\rangle
	&::=
	\langle\text{liczba}\rangle
	\mid
	\langle\text{wyrażenie}\rangle
	+
	\langle\text{wyrażenie}\rangle,
	\\
	\langle\text{liczba}\rangle
	&::=
	0 \mid 1 \mid 2 \mid \ldots
\end{aligned}
\]
Reguły te nie określają jeszcze znaczenia wyrażeń, lecz jedynie wskazują,
które ciągi symboli są poprawnymi wyrażeniami języka.

Takie podejście jest ogólne: poprzez odpowiednio dobraną gramatykę można
formalnie określić składnię języka programowania. W szczególności reguły
składniowe mogą opisywać zarówno proste języki rachunkowe, jak i języki
zawierające zmienne, funkcje, instrukcje warunkowe czy konstrukcje iteracyjne.
Dopiero po określeniu składni można przypisać poszczególnym konstrukcjom języka
ich znaczenie za pomocą semantyki denotacyjnej.

W dalszej części referatu skupimy się na prostym, modelowym języku
programowania \texttt{IMP}. Jest to niewielki język imperatywny, zawierający
podstawowe konstrukcje występujące w językach programowania, takie jak
zmienne, przypisania, operacje arytmetyczne i logiczne, instrukcje warunkowe
oraz pętle. Jego niewielki rozmiar pozwala w przejrzysty sposób przedstawić
zasady konstruowania semantyki denotacyjnej.

\subsection{Gramatyka języka \texttt{IMP}}

Składnię języka \texttt{IMP} określamy za pomocą następującej gramatyki:
\[
\begin{aligned}
	\langle a \rangle
	&::=
	x
	\mid n
	\mid \langle a \rangle + \langle a \rangle,
	\\[2mm]
	\langle b \rangle
	&::=
	\kw{true}
	\mid \kw{false}
	\mid \langle a \rangle < \langle a \rangle,
	\\[2mm]
	\langle c \rangle
	&::=
	\kw{skip}
	\mid
	x := \langle a \rangle
	\mid
	\langle c \rangle ; \langle c \rangle
	\\
	&\qquad\mid
	\kw{if}\ \langle b\rangle\
	\kw{then}\ \langle c\rangle\
	\kw{else}\ \langle c\rangle
	\mid
	\kw{while}\ \langle b\rangle\
	\kw{do}\ \langle c\rangle.
\end{aligned}
\]

W powyższej gramatyce $\langle a\rangle$ oznacza wyrażenie arytmetyczne,
$\langle b\rangle$ oznacza wyrażenie logiczne, natomiast $\langle c\rangle$
oznacza instrukcję. Symbol $n$ oznacza stałą liczbową, a $x$ zmienną programu.

Konstrukcja
\[
x := \langle a\rangle
\]
reprezentuje przypisanie wartości wyrażenia arytmetycznego do zmiennej.
Średnik oznacza sekwencyjne wykonanie dwóch instrukcji. Konstrukcja
\[
\kw{if}\ \langle b\rangle\
\kw{then}\ \langle c\rangle\
\kw{else}\ \langle c\rangle
\]
reprezentuje instrukcję warunkową, natomiast
\[
\kw{while}\ \langle b\rangle\
\kw{do}\ \langle c\rangle
\]
reprezentuje instrukcję iteracyjną.

Na tym etapie gramatyka określa jedynie, jakie konstrukcje należą do języka
\texttt{IMP}. Nie określa natomiast, jaki efekt wywołuje ich wykonanie. To
właśnie jest przedmiotem semantyki denotacyjnej, która pozwoli każdej
konstrukcji składniowej przypisać odpowiedni obiekt matematyczny opisujący
jej znaczenie.


Przez denotacyjny opis języka programowania rozumiemy w tym kontekście jego
matematyczną interpretację, która abstrahuje od konkretnej składni języka
i przypisuje poszczególnym konstrukcjom programistycznym ich formalne znaczenie.

Skupimy się teraz na denotacyjnym opisie języka \texttt{IMP}. W tym celu
wprowadzimy kilka podstawowych definicji.

\begin{definicja}
	Ustalamy dwa \emph{zbiory podstawowe} języka. Zbiór
	$\Var$ jest przeliczalnym zbiorem nazw zmiennych (identyfikatorów),
	które mogą wystąpić w programie, natomiast
	\[
	\Int=\mathbb{Z}
	\]
	jest zbiorem liczb całkowitych, stanowiących wartości wyrażeń
	arytmetycznych oraz wartości przechowywane przez zmienne. Zbiory
	$\Var$ i $\Int$ traktujemy jako dane i nie definiujemy ich
	wewnętrznej struktury.
\end{definicja}

\begin{definicja}\label{def:store}
	\emph{Stanem} nazywamy funkcję
	$\sigma:\Var \to \Int$ przypisującą zmiennym ich wartości.
	Zbiór wszystkich stanów oznaczamy przez $\Store$. Wartość zmiennej $x$
	w stanie $\sigma$ oznaczamy przez $\sigma(x)$. Ponadto
	\[
	\big(\sigma[x\mapsto n]\big)(y)=
	\begin{cases}
		n & \text{gdy } y=x,\\
		\sigma(y) & \text{gdy } y\neq x,
	\end{cases}
	\]
	oznacza stan powstały z $\sigma$ przez przypisanie zmiennej $x$ wartości
	$n$.
\end{definicja}

\begin{figure}[H]
	\centering
	\begin{tikzpicture}[
		>=stealth,
		every node/.style={font=\small},
		entry/.style={text height=1.5ex, text depth=0.25ex,
			minimum width=1.2em, inner sep=2pt}
		]
		
		% Store
		\draw[rounded corners=8pt, thick] (-6.0,-2.7) rectangle (6.0,2.7);
		\node[fill=white, inner sep=4pt] at (0,2.7)
		{$\Store=\{\sigma\mid\sigma\colon\Var\to\Int\}$};
		
		% =========================================================
		% sigma_1
		% =========================================================
		\begin{scope}[shift={(-4.0,0)}]
			\draw[rounded corners=5pt, thick] (-1.6,-2.0) rectangle (1.6,1.9);
			\node at (0,1.50) {$\sigma_1$};
			
			\node[entry] (x1) at (-0.85, 0.95) {$x$};
			\node[entry] (y1) at (-0.85, 0.35) {$y$};
			\node[entry] (z1) at (-0.85,-0.25) {$z$};
			
			\node[entry] (x1v) at (0.85, 0.95) {$3$};
			\node[entry] (y1v) at (0.85, 0.35) {$5$};
			\node[entry] (z1v) at (0.85,-0.25) {$1$};
			
			\draw[->] (x1) -- (x1v);
			\draw[->] (y1) -- (y1v);
			\draw[->] (z1) -- (z1v);
			
			\node at (-0.85,-0.80) {$\vdots$};
			\node at ( 0.85,-0.80) {$\vdots$};
			
			\node at (0,-1.55) {$\sigma_1\colon\Var\to\Int$};
		\end{scope}
		
		% =========================================================
		% sigma_2
		% =========================================================
		\begin{scope}[shift={(-0.4,0)}]
			\draw[rounded corners=5pt, thick] (-1.6,-2.0) rectangle (1.6,1.9);
			\node at (0,1.50) {$\sigma_2$};
			
			\node[entry] (x2) at (-0.85, 0.95) {$x$};
			\node[entry] (y2) at (-0.85, 0.35) {$y$};
			\node[entry] (z2) at (-0.85,-0.25) {$z$};
			
			\node[entry] (x2v) at (0.85, 0.95) {$1$};
			\node[entry] (y2v) at (0.85, 0.35) {$2$};
			\node[entry] (z2v) at (0.85,-0.25) {$9$};
			
			\draw[->] (x2) -- (x2v);
			\draw[->] (y2) -- (y2v);
			\draw[->] (z2) -- (z2v);
			
			\node at (-0.85,-0.80) {$\vdots$};
			\node at ( 0.85,-0.80) {$\vdots$};
			
			\node at (0,-1.55) {$\sigma_2\colon\Var\to\Int$};
		\end{scope}
		
		% =========================================================
		% sigma_3
		% =========================================================
		\begin{scope}[shift={(3.2,0)}]
			\draw[rounded corners=5pt, thick] (-1.6,-2.0) rectangle (1.6,1.9);
			\node at (0,1.50) {$\sigma_3$};
			
			\node[entry] (x3) at (-0.85, 0.95) {$x$};
			\node[entry] (y3) at (-0.85, 0.35) {$y$};
			\node[entry] (z3) at (-0.85,-0.25) {$z$};
			
			\node[entry] (x3v) at (0.85, 0.95) {$7$};
			\node[entry] (y3v) at (0.85, 0.35) {$2$};
			\node[entry] (z3v) at (0.85,-0.25) {$4$};
			
			\draw[->] (x3) -- (x3v);
			\draw[->] (y3) -- (y3v);
			\draw[->] (z3) -- (z3v);
			
			\node at (-0.85,-0.80) {$\vdots$};
			\node at ( 0.85,-0.80) {$\vdots$};
			
			\node at (0,-1.55) {$\sigma_3\colon\Var\to\Int$};
		\end{scope}
		
		% more states
		\node at (5.45,0) {\large$\cdots$};
		
	\end{tikzpicture}
	
	\caption{Zbiór stanów $\Store$ jako zbiór funkcji $\sigma\colon\Var\to\Int$.}
	\label{fig:store}
\end{figure}
\begin{przyklad}\label{prz:store}
	Niech $\sigma$ będzie stanem takim, że
	$\sigma(\texttt{foo})=3$ oraz $\sigma(\texttt{bar})=5$.
	Wtedy
	$\sigma[\texttt{foo}\mapsto 7]$ jest nowym stanem, w którym
	$\texttt{foo}$ ma wartość $7$, a $\texttt{bar}$ nadal ma wartość $5$:
	\[
	\big(\sigma[\texttt{foo}\mapsto 7]\big)(\texttt{bar})=5.
	\]
\end{przyklad}

Naszą uwagę powinniśmy przede wszystkim skupić na poniższym pojęciu
funkcji semantycznych.

\begin{definicja}\label{def:funkcje_semantyczne}
	Definiujemy trzy podstawowe \emph{funkcje semantyczne}:
	\[
	\Aa\sem{a}:\Store\to\Int,\qquad
	\Bb\sem{b}:\Store\to\Bool,\qquad
	\Cc\sem{c}:\Store\rightharpoonup\Store,
	\]
	przy czym zapis $\Aa\sem{a}$ oznacza \emph{denotację wyrażenia arytmetycznego}
	$a$, zapis $\Bb\sem{b}$ oznacza \emph{denotację wyrażenia logicznego} $b$, natomiast
	zapis $\Cc\sem{c}$ oznacza \emph{denotację instrukcji} $c$.
\end{definicja}

Funkcje semantyczne różnią się przeciwdziedzinami (zob.\
Rysunek~\ref{fig:funkcje_semantyczne}). Denotacja wyrażenia
arytmetycznego przypisuje każdemu stanowi \emph{liczbę} --- wartość
wyrażenia w tym stanie, a denotacja wyrażenia logicznego --- \emph{wartość
logiczną}. Obie są funkcjami całkowitymi, bo wyrażenie ma wartość
w każdym stanie. Jedynie denotacja instrukcji przenosi nas \emph{ze stanu
w stan}: stanowi początkowemu przypisuje stan końcowy. Jest to przy tym
funkcja częściowa --- jeżeli wykonanie instrukcji rozpoczęte w pewnym
stanie się nie kończy, to temu stanowi nie odpowiada żaden stan końcowy
(na rysunku stan $\sigma_3$).

\begin{figure}[H]
	\centering
	\begin{tikzpicture}[>=stealth, font=\small,
		pt/.style={inner sep=1.5pt},
		set/.style={draw, rounded corners=10pt, thick, minimum width=1.5cm, minimum height=2.6cm}]
		
		% ---------- A[[a]] : Store -> Int ----------
		\node[set] (S1) at (0,0) {};
		\node[above] at (S1.north) {$\Store$};
		\node[set] (T1) at (2.6,0) {};
		\node[above] at (T1.north) {$\Int$};
		\node[pt] (a1) at (0, 0.8) {$\sigma_1$};
		\node[pt] (a2) at (0, 0.0) {$\sigma_2$};
		\node[pt] (a3) at (0,-0.8) {$\sigma_3$};
		\node[pt] (n1) at (2.6, 0.5) {$5$};
		\node[pt] (n2) at (2.6,-0.5) {$7$};
		\draw[->] (a1) -- (n1);
		\draw[->] (a2) -- (n1);
		\draw[->] (a3) -- (n2);
		\node at (1.3,-1.9) {$\Aa\sem{a}:\Store\to\Int$};
		
		% ---------- B[[b]] : Store -> Bool ----------
		\begin{scope}[shift={(5.2,0)}]
			\node[set] (S2) at (0,0) {};
			\node[above] at (S2.north) {$\Store$};
			\node[set] (T2) at (2.6,0) {};
			\node[above] at (T2.north) {$\Bool$};
			\node[pt] (b1) at (0, 0.8) {$\sigma_1$};
			\node[pt] (b2) at (0, 0.0) {$\sigma_2$};
			\node[pt] (b3) at (0,-0.8) {$\sigma_3$};
			\node[pt] (t) at (2.6, 0.5) {$\kw{true}$};
			\node[pt] (f) at (2.6,-0.5) {$\kw{false}$};
			\draw[->] (b1) -- (t);
			\draw[->] (b2) -- (t);
			\draw[->] (b3) -- (f);
			\node at (1.3,-1.9) {$\Bb\sem{b}:\Store\to\Bool$};
		\end{scope}
		
		% ---------- C[[c]] : Store -/-> Store ----------
		\begin{scope}[shift={(10.4,0)}]
			\node[set] (S3) at (0,0) {};
			\node[above] at (S3.north) {$\Store$};
			\node[set] (T3) at (2.6,0) {};
			\node[above] at (T3.north) {$\Store$};
			\node[pt] (c1) at (0, 0.8) {$\sigma_1$};
			\node[pt] (c2) at (0, 0.0) {$\sigma_2$};
			\node[pt] (c3) at (0,-0.8) {$\sigma_3$};
			\node[pt] (d1) at (2.6, 0.8) {$\sigma_1'$};
			\node[pt] (d2) at (2.6, 0.0) {$\sigma_2'$};
			\draw[->] (c1) -- (d1);
			\draw[->] (c2) -- (d2);
			\draw[->, dashed, gray] (c3) -- ++(1.0,0) node[right, text=black, font=\footnotesize] {$\times$};
			\node at (1.3,-1.9) {$\Cc\sem{c}:\Store\rightharpoonup\Store$};
		\end{scope}
	\end{tikzpicture}
	\caption{Trzy funkcje semantyczne. Wyrażenia przyporządkowują stanom
		liczby (lewa część) lub wartości logiczne (środkowa część), natomiast
		instrukcje --- stany końcowe (prawa część). Dla stanu $\sigma_3$
		wykonanie instrukcji $c$ się nie kończy, więc $\Cc\sem{c}(\sigma_3)$
		jest nieokreślone.}
	\label{fig:funkcje_semantyczne}
\end{figure}

Definicję tę szczegółowo zobrazujemy dla każdego z tych trzech podstawowych
typów poniżej.

\subsection{Denotacje wyrażeń}

Znaczeniem wyrażenia jest funkcja przypisująca każdemu stanowi jego
wartość. Zgodnie z Definicją~\ref{def:funkcja}, dla wygody zapisu
funkcję tę przedstawiamy jako \emph{zbiór par} (relację). Para
$(\sigma,n)\in\Aa\sem{a}$ oznacza, że w stanie $\sigma$ wyrażenie $a$
przyjmuje wartość $n$.

\subsubsection{Stała liczbowa}

Denotację stałej liczbowej $n$ definiujemy jako
\[
\Aa\sem{n}=\{(\sigma,n)\mid \sigma\in\Store\}.
\]
Tutaj $n$ oznacza \emph{stałą liczbową} (numerał), czyli literał zapisany
w programie, np.\ \texttt{0}, \texttt{7} lub \texttt{42}. Symbol $n$
występujący w wyrażeniu oznacza zatem składniową postać stałej, natomiast
$n$ występujący w parze $(\sigma,n)$ oznacza odpowiadającą jej liczbę
całkowitą. Dla uproszczenia oba obiekty oznaczamy tą samą literą.
Wartość stałej liczbowej nie zależy od stanu.

\begin{przyklad}
	\[
	\Aa\sem{7}=\{(\sigma,7)\mid \sigma\in\Store\}.
	\]
	Zatem w każdym stanie $\sigma$ literał \texttt{7} przyjmuje wartość $7$.
\end{przyklad}

\subsubsection{Zmienna}

Denotację zmiennej $x$ definiujemy jako
\[
\Aa\sem{x}=\{(\sigma,\sigma(x))\mid \sigma\in\Store\}.
\]
Wartość zmiennej $x$ jest więc określona przez stan $\sigma$, a zatem może
zależeć od stanu.

\begin{przyklad}
	Jeżeli $\sigma(x)=5$, to
	\[
	(\sigma,5)\in\Aa\sem{x}.
	\]
	Oznacza to, że w stanie $\sigma$ zmienna $x$ przyjmuje wartość $5$.
\end{przyklad}

\subsubsection{Dodawanie}

Denotację sumy dwóch wyrażeń arytmetycznych definiujemy jako
\[
\Aa\sem{a_1+a_2}
=
\left\{
(\sigma,n)
\;\middle|\;
\begin{array}{c}
	(\sigma,n_1)\in\Aa\sem{a_1}
	\;\wedge\;
	(\sigma,n_2)\in\Aa\sem{a_2}\\
	{}\wedge\; n=n_1+n_2
\end{array}
\right\}.
\]
Definicja ta ma charakter kompozycyjny, ponieważ denotacja wyrażenia
$a_1+a_2$ jest określona na podstawie denotacji jego podwyrażeń $a_1$
i $a_2$. W stanie $\sigma$ najpierw wyznaczamy wartości $n_1$ i $n_2$
podwyrażeń $a_1$ i $a_2$, a następnie wartość całego wyrażenia jest równa
ich sumie:
\[
n=n_1+n_2.
\]

\begin{przyklad}\label{prz:wyr}
	Niech $\sigma(x)=5$. Wówczas
	\[
	(\sigma,5)\in\Aa\sem{x}
	\qquad\text{oraz}\qquad
	(\sigma,2)\in\Aa\sem{2}.
	\]
	Zatem
	\[
	n=5+2=7,
	\]
	więc
	\[
	(\sigma,7)\in\Aa\sem{x+2}.
	\]
	Oznacza to, że w stanie, w którym $x$ przyjmuje wartość $5$, wyrażenie
	$x+2$ przyjmuje wartość $7$.
\end{przyklad}

\subsubsection{Stała logiczna}

Denotację stałej logicznej $\kw{true}$ definiujemy jako
\[
\Bb\sem{\kw{true}}
=
\{(\sigma,\kw{true})\mid \sigma\in\Store\}.
\]
Analogicznie
\[
\Bb\sem{\kw{false}}
=
\{(\sigma,\kw{false})\mid \sigma\in\Store\}.
\]
Wartość stałej logicznej jest niezależna od stanu.

\subsubsection{Porównanie}

Denotację porównania $a_1<a_2$ definiujemy jako
\[
\begin{aligned}
	\Bb\sem{a_1<a_2}
	={}&
	\left\{
	(\sigma,\kw{true})
	\;\middle|\;
	\begin{array}{c}
		(\sigma,n_1)\in\Aa\sem{a_1},\quad
		(\sigma,n_2)\in\Aa\sem{a_2},\\
		n_1<n_2
	\end{array}
	\right\}\\
	&\cup
	\left\{
	(\sigma,\kw{false})
	\;\middle|\;
	\begin{array}{c}
		(\sigma,n_1)\in\Aa\sem{a_1},\quad
		(\sigma,n_2)\in\Aa\sem{a_2},\\
		n_1\geq n_2
	\end{array}
	\right\}.
\end{aligned}
\]
Oznacza to, że w stanie $\sigma$ najpierw wyznaczamy wartości $n_1$ i $n_2$
wyrażeń $a_1$ i $a_2$. Jeżeli
\[
n_1<n_2,
\]
to wyrażenie $a_1<a_2$ przyjmuje wartość $\kw{true}$, natomiast jeżeli
\[
n_1\geq n_2,
\]
to przyjmuje wartość $\kw{false}$.

\begin{przyklad}
	Niech $\sigma(x)=5$. Wówczas
	\[
	(\sigma,5)\in\Aa\sem{x}
	\qquad\text{oraz}\qquad
	(\sigma,1)\in\Aa\sem{1}.
	\]
	Ponieważ $5\not<1$, otrzymujemy
	\[
	(\sigma,\kw{false})\in\Bb\sem{x<1}.
	\]
	Jeżeli natomiast $\sigma(x)=0$, to $0<1$, a zatem
	\[
	(\sigma,\kw{true})\in\Bb\sem{x<1}.
	\]
\end{przyklad}

\subsection{Denotacje instrukcji}

Denotacją instrukcji jest funkcja \emph{częściowa}, która każdemu stanowi
wejściowemu przypisuje stan końcowy, o ile wykonanie instrukcji kończy się.
Zgodnie z Definicją~\ref{def:funkcja-czesciowa} funkcję tę przedstawiamy
jako zbiór par uporządkowanych. Para
$(\sigma,\sigma')\in\Cc\sem{c}$ oznacza, że wykonanie instrukcji $c$
rozpoczęte w stanie $\sigma$ kończy się w stanie $\sigma'$.

Funkcja ta jest częściowa, ponieważ wykonanie instrukcji rozpoczęte
w pewnym stanie może nie zakończyć się. W takim przypadku dla danego
stanu $\sigma$ funkcja $\Cc\sem{c}$ nie jest określona, czyli nie istnieje
żaden stan $\sigma'$ taki, że
\[
(\sigma,\sigma')\in\Cc\sem{c}.
\]
Brak takiej pary oznacza zatem niezakończenie wykonania instrukcji
rozpoczętego w stanie $\sigma$. 

\subsubsection{Instrukcja pusta}

Denotację instrukcji \kw{skip} definiujemy jako
\[
\Cc\sem{\kw{skip}}=\{(\sigma,\sigma)\mid \sigma\in\Store\}.
\]
Instrukcja \kw{skip} nie zmienia stanu, a więc stan początkowy i końcowy
są identyczne. Oznacza to, że wykonanie instrukcji \kw{skip} w stanie $\sigma$
kończy się w tym samym stanie $\sigma$ (zob.\ Rysunek~\ref{fig:skip}).

\begin{figure}[H]
	\centering
	\begin{tikzpicture}[>=stealth]
		\node (s) {$\sigma$};
		\draw[->] (s) edge[loop right, min distance=14mm, in=-30, out=30] node[right] {$\Cc\sem{\kw{skip}}$} (s);
	\end{tikzpicture}
	\caption{Denotacja instrukcji \kw{skip}: stan końcowy jest równy początkowemu.}
	\label{fig:skip}
\end{figure}

\subsubsection{Instrukcja przypisania}

Denotację instrukcji przypisania definiujemy jako
\[
\Cc\sem{x:=a}
=
\left\{
(\sigma,\sigma[x\mapsto n])
\;\middle|\;
(\sigma,n)\in\Aa\sem{a}
\right\}.
\]
Oznacza to, że w stanie $\sigma$ najpierw wyznaczamy wartość $n$ wyrażenia
$a$, a następnie otrzymujemy nowy stan $\sigma[x\mapsto n]$, w którym
zmienna $x$ przyjmuje wartość $n$, natomiast wartości pozostałych zmiennych
pozostają bez zmian (zob.\ Rysunek~\ref{fig:przypisanie}).

\begin{figure}[H]
	\centering
	\begin{tikzpicture}[>=stealth, node distance=4.2cm]
		\node (s) {$\sigma$};
		\node (s1) [right=of s] {$\sigma[x\mapsto n]$};
		\node (n) at ($(s)!0.5!(s1)+(0,1.2)$) {$n$};
		\draw[->] (s) -- node[below] {$\Cc\sem{x:=a}$} (s1);
		\draw[->, dashed] (s) -- node[above left] {$\Aa\sem{a}$} (n);
		\draw[->, dashed] (n) -- node[above right] {nowa wartość $x$} (s1);
	\end{tikzpicture}
	\caption{Denotacja przypisania: najpierw obliczamy wartość $n$ wyrażenia $a$
		w stanie $\sigma$, a następnie zmieniamy wartość zmiennej $x$ na $n$.}
	\label{fig:przypisanie}
\end{figure}

\begin{przyklad}\label{prz:instr}
	Niech stan $\sigma$ spełnia
	\[
	\sigma(x)=5
	\qquad\text{oraz}\qquad
	\sigma(y)=9.
	\]
	Wówczas
	\[
	(\sigma,7)\in\Aa\sem{x+2},
	\]
	a zatem
	\[
	(\sigma,\sigma[x\mapsto7])\in\Cc\sem{x:=x+2}.
	\]
	W stanie $\sigma[x\mapsto7]$ zmienna $x$ przyjmuje wartość $7$,
	natomiast zmienna $y$ nadal przyjmuje wartość $9$.
\end{przyklad}

\subsubsection{Złożenie instrukcji}

Instrukcja $c_1;c_2$ oznacza wykonanie dwóch instrukcji kolejno:
najpierw instrukcji $c_1$, a następnie instrukcji $c_2$. Denotację takiej
instrukcji definiujemy przez złożenie denotacji obu instrukcji:
\[
\Cc\sem{c_1;c_2}
=
\Cc\sem{c_2}\circ\Cc\sem{c_1}.
\]
Oznacza to, że stan początkowy $\sigma$ najpierw poddajemy działaniu
instrukcji $c_1$. Jeżeli jej wykonanie prowadzi do stanu $\sigma'$, a
następnie wykonanie $c_2$ ze stanu $\sigma'$ prowadzi do stanu $\sigma''$,
to całe złożenie $c_1;c_2$ prowadzi ze stanu $\sigma$ do stanu $\sigma''$
(zob.\ Rysunek~\ref{fig:zlozenie}):
\[
(\sigma,\sigma')\in\Cc\sem{c_1}
\quad\land\quad
(\sigma',\sigma'')\in\Cc\sem{c_2}
\quad\Longrightarrow\quad
(\sigma,\sigma'')\in\Cc\sem{c_1;c_2}.
\]

\begin{figure}[H]
	\centering
	\begin{tikzpicture}[>=stealth, node distance=3cm]
		\node (s) {$\sigma$};
		\node (s1) [right=of s] {$\sigma'$};
		\node (s2) [right=of s1] {$\sigma''$};
		\draw[->] (s) -- node[above] {$\Cc\sem{c_1}$} (s1);
		\draw[->] (s1) -- node[above] {$\Cc\sem{c_2}$} (s2);
		\draw[->, bend right=20, dashed] (s) to node[below] {$\Cc\sem{c_1;c_2}=\Cc\sem{c_2}\circ\Cc\sem{c_1}$} (s2);
	\end{tikzpicture}
	\caption{Denotacja złożenia instrukcji: przejście ze stanu $\sigma$ do
		$\sigma''$ przez stan pośredni $\sigma'$.}
	\label{fig:zlozenie}
\end{figure}

Ponieważ denotacje instrukcji są funkcjami częściowymi, formalnie
$\circ$ oznacza tutaj składanie relacji:
\[
R_2\circ R_1
=
\left\{
(\sigma,\sigma'')
\;\middle|\;
\exists\sigma'\,
\big(
(\sigma,\sigma')\in R_1
\land
(\sigma',\sigma'')\in R_2
\big)
\right\}.
\]

\begin{przyklad}\label{prz:zloz}
	Rozważmy instrukcję
	\[
	x:=1;\ x:=x+2
	\]
	oraz stan $\sigma$ taki, że $\sigma(x)=5$.
	
	Z definicji denotacji instrukcji przypisania mamy
	\[
	(\sigma,\sigma[x\mapsto1])
	\in
	\Cc\sem{x:=1}.
	\]
	Oznaczmy stan pośredni przez
	\[
	\sigma'=\sigma[x\mapsto1].
	\]
	Wówczas $\sigma'(x)=1$, a zatem
	\[
	(\sigma',3)
	\in
	\Aa\sem{x+2}.
	\]
	Stąd, ponownie z definicji denotacji przypisania,
	\[
	(\sigma',\sigma'[x\mapsto3])
	\in
	\Cc\sem{x:=x+2}.
	\]
	Z definicji złożenia denotacji otrzymujemy więc
	\[
	(\sigma,\sigma'[x\mapsto3])
	\in
	\Cc\sem{x:=1;\ x:=x+2}.
	\]
	Ponieważ
	\[
	\sigma'=\sigma[x\mapsto1],
	\]
	a zmiana wartości zmiennej $x$ z $1$ na $3$ daje
	\[
	\sigma'[x\mapsto3]=\sigma[x\mapsto3],
	\]
	ostatecznie
	\[
	(\sigma,\sigma[x\mapsto3])
	\in
	\Cc\sem{x:=1;\ x:=x+2}.
	\]
	Całe obliczenie można przedstawić schematycznie:
	\begin{center}
		\begin{tikzpicture}[>=stealth, node distance=2.6cm]
			\node (s) {$\sigma$};
			\node (s1) [right=of s] {$\sigma'=\sigma[x\mapsto1]$};
			\node (s2) [right=of s1] {$\sigma[x\mapsto3]$};
			\draw[->] (s) -- node[above] {$x:=1$} (s1);
			\draw[->] (s1) -- node[above] {$x:=x+2$} (s2);
			\draw[->, bend right=15, dashed] (s) to node[below] {$x:=1;\ x:=x+2$} (s2);
		\end{tikzpicture}
	\end{center}
\end{przyklad}

\subsubsection{Instrukcja warunkowa \kw{if}}

Denotację instrukcji warunkowej definiujemy jako
\begin{align*}
	\Cc\sem{\kw{if}\ b\ \kw{then}\ c_1\ \kw{else}\ c_2}
	={}&
	\left\{
	(\sigma,\sigma')
	\;\middle|\;
	\begin{array}{c}
		(\sigma,\kw{true})\in\Bb\sem{b},\\
		(\sigma,\sigma')\in\Cc\sem{c_1}
	\end{array}
	\right\}\\
	&{}\cup
	\left\{
	(\sigma,\sigma')
	\;\middle|\;
	\begin{array}{c}
		(\sigma,\kw{false})\in\Bb\sem{b},\\
		(\sigma,\sigma')\in\Cc\sem{c_2}
	\end{array}
	\right\}.
\end{align*}

Definicja ta rozpatruje dwa możliwe przypadki. Jeżeli w stanie $\sigma$
wyrażenie logiczne $b$ przyjmuje wartość $\kw{true}$, wykonywana jest
instrukcja $c_1$. Jeżeli natomiast $b$ przyjmuje wartość
$\kw{false}$, wykonywana jest instrukcja $c_2$ (zob.\ Rysunek~\ref{fig:if}).

\begin{figure}[H]
	\centering
	\begin{tikzpicture}[>=stealth]
		\node (s) at (0,0) {$\sigma$};
		\node[draw, diamond, aspect=1.6, inner sep=1pt] (b) at (2.2,0) {$\Bb\sem{b}(\sigma)$};
		\node (s1) at (7.0,1.1) {$\sigma'$};
		\node (s2) at (7.0,-1.1) {$\sigma'$};
		\draw[->] (s) -- (b);
		\draw[->] (b.north) |- node[pos=0.75, above] {$\kw{true}$:\ \ $\Cc\sem{c_1}$} (s1);
		\draw[->] (b.south) |- node[pos=0.75, below] {$\kw{false}$:\ \ $\Cc\sem{c_2}$} (s2);
	\end{tikzpicture}
	\caption{Denotacja instrukcji warunkowej: wartość warunku $b$ w stanie
		$\sigma$ decyduje, czy stan końcowy wyznacza denotacja $c_1$, czy $c_2$.}
	\label{fig:if}
\end{figure}

\begin{przyklad}\label{prz:if}
	Rozważmy instrukcję
	\[
	\kw{if}\ x<1\ \kw{then}\ y:=0\ \kw{else}\ y:=1
	\]
	oraz stan $\sigma$ taki, że
	\[
	\sigma(x)=5.
	\]
	
	Z definicji denotacji porównania otrzymujemy
	\[
	(\sigma,\kw{false})\in\Bb\sem{x<1},
	\]
	ponieważ
	\[
	5\not<1.
	\]
	Zatem w definicji denotacji instrukcji warunkowej spełniony jest drugi
	przypadek. Z definicji przypisania mamy
	\[
	(\sigma,\sigma[y\mapsto1])
	\in
	\Cc\sem{y:=1}.
	\]
	Wobec tego
	\[
	(\sigma,\sigma[y\mapsto1])
	\in
	\Cc\sem{
		\kw{if}\ x<1\ \kw{then}\ y:=0\ \kw{else}\ y:=1
	}.
	\]
	
	Rozważmy teraz stan $\sigma'$ taki, że
	\[
	\sigma'(x)=0.
	\]
	Wówczas
	\[
	(\sigma',\kw{true})\in\Bb\sem{x<1},
	\]
	ponieważ
	\[
	0<1.
	\]
	Tym razem spełniony jest pierwszy przypadek definicji. Zatem
	\[
	(\sigma',\sigma'[y\mapsto0])
	\in
	\Cc\sem{y:=0},
	\]
	a w konsekwencji
	\[
	(\sigma',\sigma'[y\mapsto0])
	\in
	\Cc\sem{
		\kw{if}\ x<1\ \kw{then}\ y:=0\ \kw{else}\ y:=1
	}.
	\]
\end{przyklad}

\subsubsection{Problem pętli \kw{while}}\label{subsec:problem_while}

Semantykę instrukcji \kw{while} możemy najpierw opisać za pomocą
poznanej wcześniej instrukcji warunkowej. Wykonanie
\[
\kw{while}\ b\ \kw{do}\ c
\]
możemy rozpatrywać w dwóch przypadkach. Jeżeli warunek $b$ jest fałszywy,
pętla kończy działanie bez wykonania instrukcji $c$, co jest równoważne
wykonaniu instrukcji \kw{skip}. Jeżeli natomiast warunek $b$ jest
prawdziwy, wykonujemy instrukcję $c$, a następnie ponownie wykonujemy
całą pętlę. Możemy zatem zapisać
\[
\boxed{
	\kw{while}\ b\ \kw{do}\ c
	\quad\equiv\quad
	\kw{if}\ b\ \kw{then}\ (c;\kw{while}\ b\ \kw{do}\ c)
	\ \kw{else}\ \kw{skip}.
}
\]
Działanie pętli przedstawia Rysunek~\ref{fig:while}: z każdego stanu
sprawdzamy warunek $b$; jeżeli jest fałszywy, pętla kończy się bez zmiany
stanu, a jeżeli jest prawdziwy, wykonujemy ciało $c$ i \emph{wracamy} do
sprawdzenia warunku w nowym stanie. Właśnie ten powrót sprawia, że
denotacji pętli nie da się zapisać tak bezpośrednio jak w przypadku
pozostałych instrukcji.

\begin{figure}[H]
	\centering
	\begin{tikzpicture}[>=stealth]
		\node (s) at (0,0) {$\sigma$};
		\node[draw, diamond, aspect=1.6, inner sep=1pt] (b) at (2.4,0) {$\Bb\sem{b}$};
		\node[draw, rounded corners, minimum width=1.4cm, minimum height=0.7cm] (c) at (2.4,-2.0) {$\Cc\sem{c}$};
		\node (end) at (6.6,0) {stan końcowy};
		\draw[->] (s) -- (b);
		\draw[->] (b) -- node[above] {$\kw{false}$} (end);
		\draw[->] (b) -- node[right] {$\kw{true}$} (c);
		\draw[->] (c.west) -- ++(-1.3,0) |- node[pos=0.25, left] {$\sigma''$} ($(b.west)+(-0.6,0)$);
	\end{tikzpicture}
	\caption{Wykonanie pętli $\kw{while}\ b\ \kw{do}\ c$: po każdym wykonaniu
		ciała $c$ wracamy do sprawdzenia warunku $b$ w nowym stanie $\sigma''$.}
	\label{fig:while}
\end{figure}

Powyższa równoważność pozwala opisać denotację pętli za pomocą jej
własnej denotacji. Oznaczmy przez
\[
W_{b,c}
=
\Cc\sem{\kw{while}\ b\ \kw{do}\ c}
\]
denotację instrukcji pętli
$\kw{while}\ b\ \kw{do}\ c$. Korzystając z definicji denotacji
instrukcji warunkowej oraz złożenia instrukcji, otrzymujemy następujące
\emph{równanie denotacyjne}:
\begin{equation}
	\label{eq:while_denotacja}
	\begin{aligned}
		W_{b,c}
		={}&
		\left\{
		(\sigma,\sigma)
		\;\middle|\;
		(\sigma,\kw{false})\in\Bb\sem{b}
		\right\}\\
		&{}\cup
		\left\{
		(\sigma,\sigma')
		\;\middle|\;
		\begin{array}{c}
			(\sigma,\kw{true})\in\Bb\sem{b},\\
			\exists\sigma''\in\Store\;
			\big(
			(\sigma,\sigma'')\in\Cc\sem{c}
			\land
			(\sigma'',\sigma')\in W_{b,c}
			\big)
		\end{array}
		\right\}.
	\end{aligned}
	\tag{W}
\end{equation}

Pierwszy składnik odpowiada sytuacji, w której warunek $b$ jest fałszywy,
a więc pętla kończy działanie bez zmiany stanu. Drugi składnik odpowiada
sytuacji, w której warunek $b$ jest prawdziwy: najpierw wykonywana jest
instrukcja $c$, prowadząc ze stanu $\sigma$ do stanu pośredniego
$\sigma''$, a następnie pętla jest wykonywana ponownie, prowadząc ze
stanu $\sigma''$ do stanu końcowego $\sigma'$.

Otrzymaliśmy w ten sposób \emph{równanie rekurencyjne} opisujące obiekt,
który chcemy zdefiniować. Nie jest to jednak jeszcze właściwa definicja
denotacji instrukcji \kw{while}, ponieważ $W_{b,c}$ występuje zarówno
po lewej, jak i po prawej stronie równania. Powstaje zatem problem
istnienia oraz wyboru odpowiedniego rozwiązania równania~\eqref{eq:while_denotacja}.

Do problemu tego powrócimy w dalszej części, po wprowadzeniu teorii punktów
stałych, która pozwoli formalnie wyznaczyć właściwe rozwiązanie równania~\eqref{eq:while_denotacja}.

%==============================================================================
\section{Równania rekurencyjne i punkty stałe}
%==============================================================================

Zaczniemy od jednej z najważniejszych definicji tego referatu.

\begin{definicja}\label{def:rownanie_stalopunktowe}
	\emph{Równaniem stałopunktowym} nazywamy każde równanie w postaci
	\begin{equation}
		\label{eq:rownanie_stalopunktowe}
		x=f(x)
		\tag{FP}
	\end{equation}
	gdzie $f:D\to D$ dla pewnego zbioru $D$. Element $x\in D$ spełniający
	równanie~\eqref{eq:rownanie_stalopunktowe} nazywamy \emph{punktem stałym} odwzorowania $f$. Zbiór wszystkich punktów
	stałych odwzorowania $f$ oznaczamy przez
	\[
	\Fix(f)=\{x\in D:f(x)=x\}.
	\]
	Jeżeli w zbiorze $\Fix(f)$ istnieje element najmniejszy względem
	ustalonego porządku na $D$, to nazywamy go \emph{najmniejszym punktem
		stałym} odwzorowania $f$ i oznaczamy przez $\lfp(f)$.
\end{definicja}

Rozważmy na początek następujący przykład.

\begin{przyklad}
	Niech $f:\mathbb{N}_0\to\mathbb{N}_0$ będzie funkcją określoną rekurencyjnie
	przez równanie
	\begin{equation}\label{prz:kwadrat}
		f(x)=
		\begin{cases}
			0 & x=0,\\
			f(x-1)+2x-1 & \text{wpp.}
		\end{cases}
	\end{equation}
	Chcemy zatem znaleźć funkcję $f$, która spełnia powyższą zależność
	rekurencyjną. Powstaje pytanie, czy taka funkcja istnieje oraz czy jest
	wyznaczona jednoznacznie.
	
	Spróbujmy skonstruować rozwiązanie poprzez kolejne przybliżenia. Załóżmy,
	że $f$ jest funkcją częściową, tzn.
	$f:\mathbb{N}_0\rightharpoonup\mathbb{N}_0$, i zdefiniujmy ciąg
	$(f_k)_{k\in\mathbb{N}_0}$ kolejnych przybliżeń. Rozpoczynamy od funkcji
	częściowej, która nie określa wartości dla żadnego argumentu:
	\[
	f_0=\varnothing.
	\]
	Następnie, korzystając z równania~\eqref{prz:kwadrat}, otrzymujemy
	\begin{align*}
		f_1
		&=
		\begin{cases}
			0 & \text{gdy }x=0,\\
			f_0(x-1)+2x-1 & \text{wpp.}
		\end{cases}
		=\{(0,0)\},\\[2pt]
		f_2
		&=
		\begin{cases}
			0 & \text{gdy }x=0,\\
			f_1(x-1)+2x-1 & \text{wpp.}
		\end{cases}
		=\{(0,0),(1,1)\},\\[2pt]
		f_3
		&=
		\begin{cases}
			0 & \text{gdy }x=0,\\
			f_2(x-1)+2x-1 & \text{wpp.}
		\end{cases}
		=\{(0,0),(1,1),(2,4)\},\\[2pt]
		&\;\;\vdots
	\end{align*}
	
	Widzimy zatem, że kolejne przybliżenia prowadzą do funkcji
	$f(x)=x^2$, która spełnia równanie~\eqref{prz:kwadrat}, ponieważ
	\[
	(x-1)^2+2x-1=x^2
	\]
	dla każdego $x\in\mathbb{N}$.
	
	Spójrzmy na ten przykład z nieco innej strony. Poszukujemy odwzorowania,
	które na podstawie danego przybliżenia rozwiązania $f_k$ wyznacza kolejne
	przybliżenie $f_{k+1}$. W istocie poszukujemy więc funkcjonału wyższego rzędu
	\[
	F:
	(\mathbb{N}_0\rightharpoonup\mathbb{N}_0)
	\to
	(\mathbb{N}_0\rightharpoonup\mathbb{N}_0)
	\]
	takiego, że
	\[
	(F(f))(x)=
	\begin{cases}
		0 & \text{gdy }x=0,\\
		f(x-1)+2x-1 & \text{wpp.}
	\end{cases}
	\]
	Wówczas kolejne przybliżenia możemy zapisać w postaci
	\[
	f_0=\varnothing,
	\qquad
	f_{k+1}=F(f_k).
	\]
	
	Rozwiązanie równania~\eqref{prz:kwadrat} jest więc funkcją $f$ spełniającą
	równanie
	\[
	F(f)=f.
	\]
	Innymi słowy, rozwiązanie jest punktem stałym odwzorowania $F$.
	
	Przeszliśmy zatem od rekurencyjnego równania określającego funkcję $f$
	do równania stałopunktowego, w którym szukamy punktu stałego odwzorowania
	$F$. Pozostaje odpowiedzieć na szereg pytań związanych z powyższym
	równaniem: kiedy rozwiązanie istnieje oraz --- jeśli rozwiązań jest więcej
	--- które z nich powinniśmy wybrać, aby odpowiadało ono naszym
	oczekiwaniom. Na te pytania postaramy się odpowiedzieć w dalszej części
	niniejszego referatu. Zastanowimy się bowiem, w jakiej klasie obiektów,
	zbiorów i funkcji rozwiązanie równania~\eqref{eq:rownanie_stalopunktowe}
	istnieje i spełnia nasze oczekiwania.
\end{przyklad}

%==============================================================================
\section{Zbiory łańcuchowo zupełne}\label{sec:cpo}
%==============================================================================

Będziemy poruszać się w terminologii zbiorów łańcuchowo zupełnych~\cite{gordon1979}; aby jednak
wprowadzić to pojęcie, potrzebujemy szeregu wcześniejszych definicji, które
zilustrujemy konkretnymi przykładami.

\subsection{Podstawowe pojęcia}

\begin{definicja}\label{def:czesciowy_porzadek}
\emph{Zbiorem częściowo uporządkowanym} nazywamy parę $(D,\sqsubseteq)$, gdzie $D$
jest dowolnym zbiorem, a $\sqsubseteq$ jest \emph{relacją binarną} na $D$ (tzn.\
${\sqsubseteq}\subseteq D\times D$), zwaną \emph{częściowym porządkiem}, spełniającą
dla dowolnych $a,b,c\in D$ następujące warunki:
\begin{itemize}
  \item[(1)] \emph{zwrotność}: $a\sqsubseteq a$;
  \item[(2)] \emph{przechodniość}: $a\sqsubseteq b \,\land\, b\sqsubseteq c \implies a\sqsubseteq c$;
  \item[(3)] \emph{antysymetryczność}: $a\sqsubseteq b \,\land\, b\sqsubseteq a \implies a=b$.
\end{itemize}
Zapis $x\sqsubseteq y$ czytamy „$x$ jest co najwyżej tak określone jak $y$”
(\emph{porządek informacyjny}). Jeżeli z kolei nie zachodzi ani $a \sqsubseteq b$, ani $b \sqsubseteq a$, to mówimy, że $a$ i $b$ są \emph{nieporównywalne}.
\end{definicja}

\begin{przyklad}\label{prz:alfabet}
Niech $A=\{a_1,a_2,\ldots,a_n\}$ dla pewnego $n\in\mathbb{N}$ będzie dowolnym
alfabetem i niech $A^*$ oznacza zbiór wszystkich słów skończonych nad $A$ wraz ze
słowem pustym $\varepsilon$. \emph{Porządkiem przedrostkowym} w $A^*$ nazywamy
relację $P\subseteq A^*\times A^*$ zdefiniowaną dla dowolnych $x,y\in A^*$ wzorem
\[
	x \, P \, y \iff y = xz \ \text{ dla pewnego } z\in A^*,
\]
tzn.\ $x$ jest \emph{przedrostkiem} słowa $y$. Wobec Definicji~\ref{def:czesciowy_porzadek}
relacja $P$ jest częściowym porządkiem w $A^*$. Dla ilustracji rozważmy ciąg
\[
	\texttt{3} \ P \ \texttt{3,1} \ P \ \texttt{3,14} \ P \ \cdots,
\]
stanowiący ciąg rozwinięć dziesiętnych kolejnych przybliżeń liczby $\pi$. Mówiąc, że
$\texttt{3,1}$ \emph{aproksymuje} $\texttt{3,14}$, mamy na myśli, że $\texttt{3,1}$
jest gorszym przybliżeniem liczby $\pi$ niż $\texttt{3,14}$. Zauważmy jeszcze, że w
tym przypadku słowa $\texttt{3,14}$ i $\texttt{3,15}$ są nieporównywalne (żadne nie
jest przedrostkiem drugiego).
\end{przyklad}

\begin{przyklad}\label{prz:potegowy}
	Niech $\mathcal{P}(D)$ oznacza zbiór wszystkich podzbiorów zbioru $D$. Zbiór ten jest częściowo uporządkowany ze względu na relację zawierania się zbiorów $\subseteq$. 
\end{przyklad}

Kluczowym w dalszych rozważaniach będzie pojęcie \emph{łańcucha}.

\begin{definicja}
	\emph{Łańcuchem} w zbiorze częściowo uporządkowanym $(D, \sqsubseteq)$ nazywamy ciąg $(d_n)_{n \in \mathbb{N}}$ elementów zbioru $D$ taki, że
	\[
	\forall n \in \mathbb{N} \quad d_n \sqsubseteq d_{n+1}.
	\]
\end{definicja}

Intuicyjnie elementy łańcucha stanowią coraz to lepsze \emph{aproksymacje} elementu granicznego, o ile taki element istnieje w rozważanym zbiorze. 

Wróćmy do Przykładu~\ref{prz:alfabet}. Łańcuch elementów zbioru $A^*$
\[
	a_1 \, P \, a_1a_1 \, P \, a_1a_1a_1 \, P \, \ldots
\]
nie ma w $A^*$ \emph{granicy}, ponieważ $A^*$ zawiera wyłącznie słowa skończone, podczas gdy \emph{granicą} tego łańcucha byłoby słowo nieskończone $a_1a_1a_1\ldots$. Gdybyśmy jednak rozszerzyli $A^*$ o wszystkie nieskończone ciągi nad $A$, to ten nieskończony ciąg można by uznać za \emph{granicę} naszego łańcucha.

Zdefiniujemy teraz dokładnie pojęcie granicy dla łańcucha.

\begin{definicja}
	\emph{Kresem górnym} lub \emph{granicą} łańcucha $(d_n)_{n \in \mathbb{N}}$  elementów zbioru częściowo uporządkowanego $(D, \sqsubseteq)$ nazywamy element $d \in D$ spełniający następujące dwa warunki:
	\begin{itemize}
		\item [(1)] $d_i \sqsubseteq d$ dla $i = 1,2,\ldots$, tzn. $d$ jest \emph{górnym ograniczeniem łańcucha},
		\item [(2)] dla każdego innego ograniczenia górnego $d'$ łańcucha zachodzi $d \sqsubseteq d'$.
	\end{itemize}
	przy czym kres górny łańcucha $(d_n)_{n \in \mathbb{N}}$ oznaczamy przez $\bigsqcup_{n=1}^{\infty} d_n.$
\end{definicja}

Finalnie możemy zdefiniować główny twór naszych rozważań. 

\begin{definicja}
	\emph{Zbiorem łańcuchowo zupełnym} nazywamy dowolny zbiór częściowo uporządkowany $(D, \sqsubseteq)$, w którym:
	\begin{itemize}
		\item [(1)] każdy łańcuch elementów z $D$ ma kres górny (granicę),
		\item [(2)] istnieje w $D$ element najmniejszy $\bt$, tzn. $\bt \sqsubseteq d$ dla każdego $d \in D$.
	\end{itemize}
\end{definicja}

Wprowadzimy dodatkowo następujące oznaczenie zbiorów łańcuchowo zupełnych.

\begin{oznaczenie}
	Niech $\boldsymbol{D} = (D, \sqsubseteq_D, \bt_D)$ jako trójka uporządkowana oznacza zbiór łańcuchowo zupełny, w którym $D$ nazwiemy \emph{nośnikiem}, $\sqsubseteq_D$ relacją porządku oraz $\bt_D$ elementem najmniejszym.
\end{oznaczenie}

\begin{comment}

\subsection{Operacje mnogościowe na zbiorach łańcuchowo zupełnych}

W semantyce denotacyjnej często spotykamy się z sytuacją, gdy mając dany zbiór łańcuchowo zupełny, chcemy utworzyć nowy zbiór łańcuchowo zupełny, przy czym zwykle operacje teoriomnogościowe na zbiorach są niewystarczające.

Ustalmy dwa zbiory łańcuchowo zupełne $\boldsymbol{A} = (A, \sqsubseteq_A, \bt_A)$ i $\boldsymbol{B} = (B, \sqsubseteq_B, \bt_B)$, aby nie powielać ich w kolejnych definicjach.

\begin{definicja}
	\emph{Sumą rozłączną} zbiorów łańcuchowo zupełnych $\boldsymbol{A}$ i $\boldsymbol{B}$ nazwiemy zbiór
	\[
	\boldsymbol{A} + \boldsymbol{B}
	=
	\left(
	\left(A - \{\bt_A\}\right)
	+
	\left(B - \{\bt_B\}\right)
	+
	\{\bt\},
	\sqsubseteq,
	\bt
	\right),
	\]
	gdzie:
	\begin{itemize}
		\item [(1)] Nośnik $\boldsymbol{A} + \boldsymbol{B}$ powstaje z nośników wyjściowych przez usunięcie z nich elementów najmniejszych, zsumowanie ich w sposób rozłączny, tzn. jeżeli $A$ i $B$ są rozłączne, to korzystamy ze zwykłej sumy teoriomnogościowej, w przeciwnym wypadku rozważamy sumę
		\[
		(A \times \{0\}) + (B \times \{1\}),
		\]
		a następnie uzupełniamy ją o nowy, dodatkowy, wyróżniony element $\bt$, który następnie uznajemy za element najmniejszy w zbiorze łańcuchowo zupełnym $\boldsymbol{A} + \boldsymbol{B}$.
			\item [(2)] Porządek w zbiorze $\boldsymbol{A} + \boldsymbol{B}$ definiujemy przez zachowanie w zbiorach $A$ i $B$ odpowiednich porządków, przyjmując, że elementy zbiorów $A$ i $B$ są ze sobą nieporównywalne, oraz uznając element $\bt$ za element najmniejszy dla całej sumy, tzn.
		\[
		x \sqsubseteq y \iff x = \bt \text{ lub } (x,y \in A \text{ oraz } x \sqsubseteq_A y) \text{ lub } (x,y \in B \text{ oraz } x \sqsubseteq_B y).
		\]
	\end{itemize}
	
	Nieformalnie widzimy, że suma rozłączna zbiorów łańcuchowo zupełnych zachowuje relacje wśród elementów składowych zbiorów, uznając elementy różnych zbiorów za nieporównywalne, oraz \emph{skleja} ze sobą elementy najmniejsze dwóch zbiorów w jeden.
\end{definicja}

\begin{definicja}
	Iloczynem kartezjańskim zbiorów łańcuchowo zupełnych $\boldsymbol{A}$ i $\boldsymbol{B}$ nazywamy zbiór
	\[
	\boldsymbol{A} \times \boldsymbol{B}
	=
	\left(A \times B, \sqsubseteq, (\bt_A,\bt_B)\right),
	\]
	gdzie:
	\begin{itemize}
			\item [(1)] $A \times B$ jest zwykłym iloczynem kartezjańskim zbiorów,
			\item [(2)] porządek w $\boldsymbol{A} \times \boldsymbol{B}$ definiujemy po współrzędnych, tzn.
			\[
			(a_1,b_1) \sqsubseteq (a_2,b_2)
			\iff
			a_1 \sqsubseteq_A a_2 \land b_1 \sqsubseteq_B b_2,
			\]
			\item [(3)] $(\bt_A,\bt_B)$ stanowi element najmniejszy zbioru $\boldsymbol{A} \times \boldsymbol{B}$.
		\end{itemize}
\end{definicja}

\end{comment}
%
\subsection{Porządek płaski}
%------------------------------------------------------------------------------

Najprostszy sposób nadania dowolnemu zbiorowi struktury zbioru łańcuchowo
zupełnego polega na dołożeniu jednego elementu najmniejszego, reprezentującego
brak informacji.

\begin{definicja}\label{def:zbior_plaski}
	Niech $D$ będzie dowolnym zbiorem. Rozszerzmy go o nowy element
	$\bt\notin D$, tworząc zbiór
	\[
	D_\bt = D \cup \{\bt\}.
	\]
	Parę $(D_\bt,\sqsubseteq)$ nazywamy \emph{płaską} (a porządek
	$\sqsubseteq$ --- \emph{płaskim}), jeżeli
	\begin{itemize}
		\item [(1)] $\bt$ jest elementem najmniejszym zbioru $D_\bt$,
		\item [(2)] $d_1 \sqsubseteq d_2 \iff d_1 = \bt \ \text{lub}\ d_1 = d_2$.
	\end{itemize}
\end{definicja}

Powyższą definicję można łatwo zilustrować na diagramie (Rysunek~\ref{fig:porzadek_plaski}), na którym odpowiednie strzałki oznaczają relację $\sqsubseteq$.

\begin{figure}[H]
	\centering
	\begin{tikzpicture}[node distance=2cm]
		\node (bottom) {$\bt$};
		\node (a) [above left=of bottom] {$d_1$};
		\node (b) [above=of bottom] {$d_2$};
		\node (c) [above right=of bottom] {$d_3$};
		\draw[->] (bottom) -- (a);
		\draw[->] (bottom) -- (b);
		\draw[->] (bottom) -- (c);
	\end{tikzpicture}
	\caption{Porządek płaski na zbiorze $D_\bt=\{\bt,d_1,d_2,d_3\}$.}
	\label{fig:porzadek_plaski}
\end{figure}

Elementy $d_1,d_2,d_3$ są parami nieporównywalne --- jedynym elementem leżącym
poniżej pozostałych jest $\bt$.

\begin{uwaga}\label{uw:lancuchy_plaskie}
	Każdy zbiór płaski $D_\bt$ jest łańcuchowo zupełny. W łańcuchu nie mogą
	bowiem wystąpić dwa różne elementy różne od $\bt$ (byłyby
	nieporównywalne), więc każdy łańcuch ma postać
	\[
	\bt\sqsubseteq\ldots\sqsubseteq\bt\sqsubseteq d\sqsubseteq d\sqsubseteq\ldots
	\]
	(co najwyżej jedna wartość $d\neq\bt$); jest zatem od pewnego miejsca stały,
	a jego kresem górnym jest $d$ (lub $\bt$, gdy wszystkie wyrazy są równe
	$\bt$). Elementem najmniejszym jest $\bt$.
\end{uwaga}

%------------------------------------------------------------------------------
\subsection{Porządki w zbiorach funkcji}
%------------------------------------------------------------------------------

Zbiory funkcji można porządkować na dwa naturalne sposoby, prowadzące do dwóch
różnych --- choć, jak się okaże, powiązanych --- konstrukcji zbiorów łańcuchowo
zupełnych.

\subsubsection{Porządek poziomy}

\begin{definicja}\label{def:zbior_funkcji_czesciowych}
	Niech $A \rightharpoonup B$ oznacza zbiór funkcji częściowych ze zbioru $A$ w zbiór $B$. 
\end{definicja}

Zbiór ten możemy naturalnie uporządkować relacją inkluzji:

\begin{twierdzenie}\label{tw:zbior_funkcji_czesciowych_czesciowo_uporzadkowany}
	Zbiór funkcji częściowych $A \rightharpoonup B$ można częściowo uporządkować relacją zawierania się wykresów
	\[
		f\sqsubseteq g \iff f\subseteq g.
	\]
\end{twierdzenie}

\begin{proof}
	Każdą funkcję częściową $f:A\rightharpoonup B$ utożsamiamy z jej
	wykresem, czyli relacją $f\subseteq A\times B$, spełniającą warunek
	jednoznaczności
	\[
	(a,b)\in f\land(a,b')\in f
	\implies b=b'.
	\]
	Na zbiorze funkcji częściowych rozważamy relację
	\[
	f\sqsubseteq g
	\iff
	f\subseteq g.
	\]
	Relacja ta jest częściowym porządkiem, ponieważ inkluzja zbiorów jest
	zwrotna, antysymetryczna i przechodnia.
\end{proof}

\begin{definicja}\label{def:porzadek_poziomy}
	Porządek w zbiorze funkcji częściowych $A \rightharpoonup B$ z relacją częściowego porządku daną wzorem 
	\[
		f\sqsubseteq g \iff f\subseteq g
	\]
	nazywamy \emph{porządkiem poziomym}. Intuicyjnie $f\sqsubseteq g$ znaczy, że $g$ jest \emph{co najmniej tak określona} jak $f$ --- ma
	dziedzinę nie mniejszą niż $f$ i zgadza się z $f$ tam, gdzie $f$ jest określona (zob.\ Rysunek~\ref{fig:porzadek_poziomy}).
\end{definicja}

\begin{figure}[H]
	\centering
	\begin{tikzpicture}[>=stealth, scale=1]
		\draw[->] (-1.6,0) -- (1.6,0) node[right] {$A$};
		\draw[->] (0,-0.3) -- (0,2.8) node[above] {$B$};
		\draw[thick] (-1.1,0.45) -- (-0.6,0.75) -- (-0.1,1.05) -- (0.35,1.35);
		\draw[thick,dashed] (0.35,1.35) -- (0.75,1.65) -- (1.15,1.9) -- (1.4,2.2);
		\node[above left] at (-0.6,0.75) {$f=g$};
		\node[above left] at (0.9,1.8) {$g$};
		\draw[dashed] (-1.1,0) -- (-1.1,-0.8);
		\draw[dashed] (0.35,0) -- (0.35,-0.8);
		\draw[dashed] (-1.1,-0.8) -- (0.35,-0.8);
		\node[below] at (-0.375,-0.8) {$\dom(f)$};
		\draw[dashed] (-1.1,-0.05) -- (-1.1,-1.35);
		\draw[dashed] (1.4,0) -- (1.4,-1.35);
		\draw[dashed] (-1.1,-1.35) -- (1.4,-1.35);
		\node[below] at (0.15,-1.35) {$\dom(g)$};
	\end{tikzpicture}
	\caption{Porządek poziomy: $f\sqsubseteq g \iff f\subseteq g$. Funkcja $g$ (linia ciągła i przerywana) rozszerza $f$ (linia ciągła): pokrywa się z $f$ na $\dom(f)\subseteq\dom(g)$ i jest dodatkowo określona poza $\dom(f)$.}
	\label{fig:porzadek_poziomy}
\end{figure}

\begin{twierdzenie}\label{tw:funkcje_czesciowe_cpo}
	Dla dowolnych zbiorów $A$ i $B$ zbiór funkcji częściowych
	\[
	\boldsymbol{A\rightharpoonup B} = (A \rightharpoonup B, \subseteq, \varnothing),
	\]
	uporządkowany relacją inkluzji wykresów, jest zbiorem łańcuchowo
	zupełnym. Elementem najmniejszym jest funkcja pusta $\varnothing$,
	a kres górny każdego łańcucha $(f_n)_{n\in\mathbb{N}}$ jest równy
	\[
	\bigsqcup_{n=1}^{\infty} f_n
	=
	\bigcup_{n=1}^{\infty} f_n.
	\]
\end{twierdzenie}

\begin{proof}
	\emph{Częściowy porządek}. Istotnie wobec Twierdzenia~\ref{tw:zbior_funkcji_czesciowych_czesciowo_uporzadkowany} zbiór funkcji częściowych $A \rightharpoonup B$ jest zbiorem częściowo uporządkowanym. 
	
	\emph{Element najmniejszy}. Funkcja pusta $\varnothing$ jest elementem najmniejszym, gdyż
	\[
	\varnothing\subseteq f
	\]
	dla każdej funkcji częściowej $f:A\rightharpoonup B$.
	
	\emph{Kres górny łańcucha}. Niech
	\[
	f_1\sqsubseteq f_2\sqsubseteq f_3\sqsubseteq\ldots
	\]
	będzie dowolnym łańcuchem funkcji częściowych $A\rightharpoonup B$.
	Zdefiniujmy relację
	\[
	f=\bigcup_{n=1}^{\infty}f_n.
	\]
	
	Najpierw pokażemy, że $f$ jest funkcją częściową. Ponieważ
	$f_n\subseteq A\times B$ dla każdego $n\in\mathbb{N}$, otrzymujemy
	\[
	f\subseteq A\times B.
	\]
	Niech
	\[
	(a,b)\in f
	\quad\land\quad
	(a,b')\in f.
	\]
	Z definicji sumy zbiorów istnieją $i,j\in\mathbb{N}$ takie, że
	\[
	(a,b)\in f_i
	\quad\land\quad
	(a,b')\in f_j.
	\]
	Ponieważ $(f_n)_{n\in\mathbb{N}}$ jest łańcuchem, elementy $f_i$
	i $f_j$ są porównywalne, zatem
	\[
	f_i\subseteq f_j
	\quad\lor\quad
	f_j\subseteq f_i.
	\]
	Bez straty ogólności załóżmy, że
	\[
	f_i\subseteq f_j.
	\]
	Wówczas
	\[
	(a,b)\in f_i
	\implies
	(a,b)\in f_j.
	\]
	Zatem
	\[
	(a,b)\in f_j
	\quad\land\quad
	(a,b')\in f_j.
	\]
	Ponieważ $f_j$ jest funkcją częściową, otrzymujemy
	\[
	b=b'.
	\]
	Stąd $f$ spełnia warunek jednoznaczności, a więc
	\[
	f\in A\rightharpoonup B.
	\]
	
	Następnie pokażemy, że $f$ jest ograniczeniem górnym łańcucha.
	Dla każdego $n\in\mathbb{N}$ z definicji sumy zbiorów mamy
	\[
	f_n\subseteq\bigcup_{k=1}^{\infty}f_k=f.
	\]
	Zatem
	\[
	\forall n\in\mathbb{N}\quad f_n\sqsubseteq f,
	\]
	co oznacza, że $f$ jest ograniczeniem górnym łańcucha.
	
	Pozostaje wykazać, że $f$ jest najmniejszym ograniczeniem górnym.
	Niech $g:A\rightharpoonup B$ będzie dowolnym ograniczeniem górnym
	łańcucha $(f_n)_{n\in\mathbb{N}}$. Wówczas
	\[
	\forall n\in\mathbb{N}\quad f_n\subseteq g.
	\]
	Weźmy dowolny element $(a,b)\in f$. Z definicji sumy zbiorów
	istnieje $n\in\mathbb{N}$ takie, że
	\[
	(a,b)\in f_n.
	\]
	Ponieważ $f_n\subseteq g$, otrzymujemy
	\[
	(a,b)\in g.
	\]
	Zatem
	\[
	f\subseteq g.
	\]
	Stąd
	\[
	f\sqsubseteq g
	\]
	dla każdego ograniczenia górnego $g$ tego łańcucha. W konsekwencji
	$f$ jest jego najmniejszym ograniczeniem górnym, czyli
	\[
	\bigsqcup_{n=1}^{\infty}f_n
	=
	\bigcup_{n=1}^{\infty}f_n.
	\]
	
	Ostatecznie $A\rightharpoonup B$ jest częściowo uporządkowany,
	posiada element najmniejszy $\varnothing$, a każdy jego łańcuch
	posiada kres górny. Zatem
	\[
	\boldsymbol{A\rightharpoonup B}
	=
	(A\rightharpoonup B,\subseteq,\varnothing)
	\]
	jest zbiorem łańcuchowo zupełnym.
\end{proof}

Analogicznie można pokazać, że zbiór $\mathcal{P}(D)$ z Przykładu \ref{prz:potegowy} jest zbiorem łańcuchowo zupełnym.

\subsubsection{Porządek pionowy}

\begin{definicja}\label{def:zbior_funkcji_calkowitych}
	Niech $A \to B$ oznacza zbiór wszystkich funkcji całkowitych ze zbioru $A$ w zbiór $B$.
\end{definicja}

Zdefiniujemy teraz porządek w zbiorze funkcji całkowitych.

\begin{twierdzenie}\label{tw:zbior_funkcji_calkowitych_czesciowo_uporzadkowany}
	Niech $(B, \sqsubseteq_B)$ będzie zbiorem częściowo uporządkowanym. Wówczas zbiór funkcji całkowitych $A \to B$ możemy uporządkować relacją częściowego porządku
	\[
		f\sqsubseteq g \iff \forall a\in A\quad f(a)\sqsubseteq_B g(a).
	\]
\end{twierdzenie}

\begin{proof}
	Relacja $\sqsubseteq$ jest zadana punktowo, więc dziedziczy własności po $\sqsubseteq_B$. Dla dowolnych $f,g,h:A\to B$:
	\begin{itemize}
		\item[(1)] \emph{zwrotność}: dla każdego $a\in A$ mamy
		$f(a)\sqsubseteq_B f(a)$, zatem $f\sqsubseteq f$;
		\item[(2)] \emph{przechodniość}: jeżeli $f\sqsubseteq g$ oraz
		$g\sqsubseteq h$, to dla każdego $a \in A$ zachodzi
		$f(a)\sqsubseteq_B g(a)\sqsubseteq_B h(a)$, skąd
		$f(a)\sqsubseteq_B h(a)$, czyli $f\sqsubseteq h$;
		\item[(3)] \emph{antysymetryczność}: jeżeli $f\sqsubseteq g$ oraz
		$g\sqsubseteq f$, to dla każdego $a \in A$ mamy $f(a)\sqsubseteq_B g(a)$
		i $g(a)\sqsubseteq_B f(a)$. Stąd z antysymetrii $\sqsubseteq_B$ wynika
		$f(a)=g(a)$, a więc $f=g$.
	\end{itemize}
\end{proof}

\begin{definicja}\label{def:porzadek_pionowy}
	Niech $(B, \sqsubseteq_B)$ będzie zbiorem częściowo uporządkowanym. Porządek w zbiorze funkcji całkowitych $A \to B$ zadany relacją 
	\[
	f\sqsubseteq g \iff \forall a\in A\quad f(a)\sqsubseteq_B g(a),
	\]
	nazywamy \emph{porządkiem pionowym}. Intuicyjnie $f\sqsubseteq g$ znaczy, że wykres $g$
	leży ,,nad'' wykresem $f$ w sensie porządku $\sqsubseteq_B$ na wartościach
	(zob.\ Rysunek~\ref{fig:porzadek_pionowy}).
\end{definicja}

\begin{figure}[H]
	\centering
	\begin{tikzpicture}[>=stealth, scale=1]
		
		% ============================================================
		% PORZĄDEK PIONOWY
		% ============================================================
		\draw[->] (4.4,0) -- (7.6,0) node[right] {$A$};
		\draw[->] (6,-0.3) -- (6,2.8) node[above] {$B$};
		
		% f
		\draw[thick]
		(4.8,0.55) -- (5.5,0.95) -- (6.3,1.25) -- (7.2,1.6);
		
		% g
		\draw[thick,dashed]
		(4.8,0.95) -- (5.5,1.35) -- (6.3,1.65) -- (7.2,2.0);
		
		\node[right] at (7.2,1.6) {$f$};
		\node[right] at (7.2,2.0) {$g$};
		
	\end{tikzpicture}
	\caption{Porządek pionowy: $f\sqsubseteq g \iff \forall a\in A\ f(a)\sqsubseteq_B g(a)$. Wykres funkcji $g$ (linia przerywana) leży nad wykresem funkcji $f$ (linia ciągła).}
	\label{fig:porzadek_pionowy}
\end{figure}

\begin{twierdzenie}\label{tw:funkcje_calkowite_cpo}
	Niech $A$ będzie dowolnym zbiorem, a $\boldsymbol{B}=(B,\sqsubseteq_B,\bt_B)$
	--- zbiorem łańcuchowo zupełnym. Wówczas zbiór funkcji całkowitych
	$\boldsymbol{A\to B}=(A\to B,\sqsubseteq,\varphi)$ z porządkiem pionowym jest zbiorem łańcuchowo
	zupełnym. Kres górny łańcucha wyznaczamy punktowo,
	\[
	\Big(\bigsqcup_{n=1}^{\infty} f_n\Big)(a)=\bigsqcup_{n=1}^{\infty} f_n(a)
	\qquad(a\in A),
	\]
	a elementem najmniejszym jest funkcja stała $\varphi(a)=\bt_B$ dla dowolnego $a \in A$.
\end{twierdzenie}

\begin{proof}
	\emph{Częściowy porządek.} Istotnie wobec Twierdzenia~\ref{tw:zbior_funkcji_calkowitych_czesciowo_uporzadkowany} zbiór funkcji całkowitych jest zbiorem częściowo uporządkowanym względem relacji $\sqsubseteq$. 
	
	\emph{Element najmniejszy.}
	Funkcja $\varphi:A\to B$ jest określona wzorem
	\[
	\varphi(a)=\bt_B
	\qquad(a\in A).
	\]
	Niech $f:A\to B$ będzie dowolną funkcją. Ponieważ $\bt_B$ jest
	elementem najmniejszym zbioru $\boldsymbol{B}$, zachodzi
	\[
	\bt_B\sqsubseteq_B f(a)
	\qquad(a\in A).
	\]
	Stąd
	\[
	\varphi(a)\sqsubseteq_B f(a)
	\qquad(a\in A),
	\]
	a zatem, z definicji porządku pionowego,
	\[
	\varphi\sqsubseteq f.
	\]
	Ponieważ $f\in A\to B$ było dowolne, $\varphi$ jest elementem
	najmniejszym zbioru $A\to B$.
	
	\emph{Kres górny łańcucha.}
	Niech
	\[
	f_1\sqsubseteq f_2\sqsubseteq f_3\sqsubseteq\ldots
	\]
	będzie dowolnym łańcuchem w $A\to B$.
	
	Pokażemy najpierw, że dla każdego $a\in A$ ciąg
	\[
	f_1(a)\sqsubseteq_B f_2(a)\sqsubseteq_B f_3(a)\sqsubseteq_B\ldots
	\]
	jest łańcuchem w $\boldsymbol{B}$. Istotnie, z
	\[
	f_n\sqsubseteq f_{n+1}
	\]
	oraz z definicji porządku punktowego wynika
	\[
	f_n(a)\sqsubseteq_B f_{n+1}(a)
	\qquad(a\in A).
	\]
	
	Ponieważ $\boldsymbol{B}$ jest zbiorem łańcuchowo zupełnym, dla każdego
	$a\in A$ istnieje kres górny powyższego łańcucha. Definiujemy zatem
	funkcję $f:A\to B$ wzorem
	\[
	f(a)
	:=
	\bigsqcup_{n=1}^{\infty}f_n(a)
	\qquad(a\in A).
	\]
	Jest to poprawnie określona funkcja całkowita, ponieważ
	\[
	f(a)\in B
	\]
	dla każdego $a\in A$.
	
	Pokażemy, że funkcja $f$ jest kresem górnym łańcucha
	$(f_n)_{n\in\mathbb{N}}$.
	
	Najpierw wykażemy, że $f$ jest ograniczeniem górnym. Niech
	$n\in\mathbb{N}$. Z definicji kresu górnego w $\boldsymbol{B}$ mamy
	\[
	f_n(a)
	\sqsubseteq_B
	\bigsqcup_{k=1}^{\infty}f_k(a)
	=
	f(a)
	\qquad(a\in A).
	\]
	Zatem
	\[
	\forall a\in A\quad f_n(a)\sqsubseteq_B f(a),
	\]
	a więc
	\[
	f_n\sqsubseteq f.
	\]
	Ponieważ $n$ było dowolne, otrzymujemy
	\[
	\forall n\in\mathbb{N}\quad f_n\sqsubseteq f.
	\]
	Zatem $f$ jest ograniczeniem górnym rozważanego łańcucha.
	
	Pozostaje wykazać, że $f$ jest najmniejszym ograniczeniem górnym.
	Niech
	\[
	g:A\to B
	\]
	będzie dowolnym ograniczeniem górnym łańcucha $(f_n)_{n\in\mathbb{N}}$.
	Wtedy
	\[
	\forall n\in\mathbb{N}\quad f_n\sqsubseteq g.
	\]
	Z definicji porządku punktowego otrzymujemy
	\[
	\forall n\in\mathbb{N}\;\forall a\in A\quad
	f_n(a)\sqsubseteq_B g(a).
	\]
	
	Ustalmy $a\in A$. Wówczas
	\[
	f_n(a)\sqsubseteq_B g(a)
	\qquad(n\in\mathbb{N}),
	\]
	co oznacza, że $g(a)$ jest ograniczeniem górnym łańcucha
	\[
	f_1(a)\sqsubseteq_B f_2(a)\sqsubseteq_B\ldots
	\]
	w $\boldsymbol{B}$. Ponieważ $f(a)$ jest jego najmniejszym
	ograniczeniem górnym, otrzymujemy
	\[
	f(a)
	=
	\bigsqcup_{n=1}^{\infty}f_n(a)
	\sqsubseteq_B g(a).
	\]
	Ponieważ $a\in A$ było dowolne, mamy
	\[
	\forall a\in A\quad f(a)\sqsubseteq_B g(a).
	\]
	Z definicji porządku punktowego wynika stąd
	\[
	f\sqsubseteq g.
	\]
	Zatem $f$ jest najmniejszym ograniczeniem górnym łańcucha
	$(f_n)_{n\in\mathbb{N}}$.
	
	W konsekwencji
	\[
	\bigsqcup_{n=1}^{\infty}f_n=f,
	\]
	przy czym kres górny jest wyznaczany punktowo:
	\[
	\left(\bigsqcup_{n=1}^{\infty}f_n\right)(a)
	=
	\bigsqcup_{n=1}^{\infty}f_n(a).
	\]
	
	Zatem $\boldsymbol{A\to B}$ posiada element najmniejszy oraz kres
	górny każdego łańcucha, a więc jest zbiorem łańcuchowo zupełnym.
\end{proof}

\begin{wniosek}
	W szczególności, biorąc za $\boldsymbol{B}$ zbiór płaski $\boldsymbol{B} = (B_\bt, \sqsubseteq, \bt)$ otrzymujemy, że zbiór funkcji całkowitych $A\to B_\bt$ jest zbiorem
łańcuchowo zupełnym.
\end{wniosek}

\subsection{Izomorfizm zbiorów łańcuchowo zupełnych}

Okazuje się, że oba spojrzenia opisują ten sam obiekt: funkcje częściowe
$A\rightharpoonup B$ z porządkiem poziomym to dokładnie funkcje całkowite
$A\to B_\bt$ z porządkiem pionowym. Stąd możemy mówić o ich \emph{izomorfizmie}. 

\begin{definicja}
	\emph{Izomorfizmem} zbiorów łańcuchowo zupełnych $\boldsymbol{A}$ i
	$\boldsymbol{B}$ nazywamy bijekcję $\alpha$ zbioru $A$ na $B$
	zachowującą porządek w obie strony oraz element najmniejszy, tzn.
	\begin{itemize}
		\item [(1)] $a_1 \sqsubseteq_A a_2 \iff \alpha(a_1)\sqsubseteq_B\alpha(a_2)$
		dla dowolnych $a_1,a_2\in A$,
		\item [(2)] $\alpha(\bt_A)=\bt_B$.
	\end{itemize}
\end{definicja}

\begin{twierdzenie}\label{tw:izomorfizm_czesciowe_calkowite}
	Niech $A,B$ będą dowolnymi zbiorami, a $B_\bt$ --- płaskim zbiorem
	powstałym z $B$. Wówczas zbiory łańcuchowo
	zupełne
	\[
	(A\rightharpoonup B,\ \subseteq,\ \varnothing)
	\qquad\text{oraz}\qquad
	(A\to B_\bt,\ \sqsubseteq,\ \varphi)
	\]
	są izomorficzne.
\end{twierdzenie}

\begin{proof}
	Zdefiniujmy odwzorowanie
	\[
	\alpha:(A\rightharpoonup B)\longrightarrow(A\to B_\bt)
	\]
	wzorem
	\[
	\alpha(f)(a)=
	\begin{cases}
		f(a),&a\in\dom(f),\\
		\bt,&a\notin\dom(f).
	\end{cases}
	\]
	Odwzorowanie $\alpha$ uzupełnia więc funkcję częściową $f$ wartością
	$\bt$ we wszystkich punktach, w których $f$ nie jest określona.
	
	Zdefiniujmy również odwzorowanie w przeciwną stronę,
	\[
	\beta:(A\to B_\bt)\longrightarrow(A\rightharpoonup B),
	\]
	które z funkcji całkowitej $g:A\to B_\bt$ ,,wycina'' punkty, w których
	przyjmuje ona wartość $\bt$:
	\[
	\beta(g)=\{(a,g(a))\mid a\in A,\ g(a)\neq\bt\}.
	\]
	Sprawdźmy, że $\beta(g)$ jest rzeczywiście funkcją częściową z $A$
	w $B$.
	\begin{itemize}
		\item $\beta(g)\subseteq A\times B$. Istotnie, jeżeli
		$(a,b)\in\beta(g)$, to $b=g(a)$ oraz $g(a)\neq\bt$. Ponieważ
		$g(a)\in B_\bt=B\cup\{\bt\}$, otrzymujemy $b\in B$.
		\item Warunek jednoznaczności. Jeżeli $(a,b)\in\beta(g)$ oraz
		$(a,b')\in\beta(g)$, to z definicji $\beta$ mamy $b=g(a)$ oraz
		$b'=g(a)$, a więc $b=b'$.
	\end{itemize}
	Ponadto z definicji
	\[
	\dom\big(\beta(g)\big)=\{a\in A:\ g(a)\neq\bt\},
	\qquad
	\beta(g)(a)=g(a)\quad\text{dla }a\in\dom\big(\beta(g)\big).
	\]
	
	\emph{Bijektywność.} Pokażemy, że $\beta$ jest odwrotnością $\alpha$,
	tzn.
	\[
	\beta(\alpha(f))=f
	\quad\text{dla każdej } f:A\rightharpoonup B
	\qquad\text{oraz}\qquad
	\alpha(\beta(g))=g
	\quad\text{dla każdej } g:A\to B_\bt.
	\]
	
	\emph{(i) Równość $\beta(\alpha(f))=f$.} Obie strony są podzbiorami
	$A\times B$, więc wystarczy pokazać, że mają te same elementy. Niech
	$(a,b)\in A\times B$. Z definicji $\beta$
	\[
	(a,b)\in\beta(\alpha(f))
	\iff
	\alpha(f)(a)=b
	\ \land\
	\alpha(f)(a)\neq\bt.
	\]
	Ponieważ $b\in B$, mamy $b\neq\bt$, więc drugi warunek wynika
	z pierwszego i
	\[
	(a,b)\in\beta(\alpha(f))
	\iff
	\alpha(f)(a)=b.
	\]
	Jeżeli $a\notin\dom(f)$, to $\alpha(f)(a)=\bt\neq b$, więc równość
	$\alpha(f)(a)=b$ może zachodzić tylko dla $a\in\dom(f)$, a wtedy
	$\alpha(f)(a)=f(a)$. Zatem
	\[
	\alpha(f)(a)=b
	\iff
	a\in\dom(f)\ \land\ f(a)=b
	\iff
	(a,b)\in f.
	\]
	Stąd $\beta(\alpha(f))=f$.
	
	\emph{(ii) Równość $\alpha(\beta(g))=g$.} Obie strony są funkcjami
	całkowitymi z $A$ w $B_\bt$, więc wystarczy porównać ich wartości
	w dowolnym punkcie $a\in A$. Rozważmy dwa przypadki.
	\begin{itemize}
		\item Jeżeli $g(a)\neq\bt$, to $a\in\dom(\beta(g))$ oraz
		$\beta(g)(a)=g(a)$. Z definicji $\alpha$
		\[
		\alpha(\beta(g))(a)=\beta(g)(a)=g(a).
		\]
		\item Jeżeli $g(a)=\bt$, to $a\notin\dom(\beta(g))$, więc
		z definicji $\alpha$
		\[
		\alpha(\beta(g))(a)=\bt=g(a).
		\]
	\end{itemize}
	W obu przypadkach $\alpha(\beta(g))(a)=g(a)$, a ponieważ $a$ było
	dowolne, $\alpha(\beta(g))=g$.
	
	Z (i) wynika, że $\alpha$ jest różnowartościowe: jeżeli
	$\alpha(f)=\alpha(f')$, to
	\[
	f=\beta(\alpha(f))=\beta(\alpha(f'))=f'.
	\]
	Z (ii) wynika, że $\alpha$ jest ,,na'': każda funkcja $g:A\to B_\bt$
	jest postaci $g=\alpha(\beta(g))$. Zatem $\alpha$ jest bijekcją,
	a $\beta$ jest do niej odwrotna.
	
	\emph{Zachowanie porządku.} Pokażemy, że dla dowolnych
	$f,g:A\rightharpoonup B$
	\[
	f\subseteq g
	\iff
	\alpha(f)\sqsubseteq\alpha(g).
	\]
	
	($\Rightarrow$) Załóżmy, że $f\subseteq g$, i ustalmy $a\in A$.
	\begin{itemize}
		\item Jeżeli $a\notin\dom(f)$, to $\alpha(f)(a)=\bt$, a ponieważ
		$\bt$ jest elementem najmniejszym $B_\bt$, mamy
		$\alpha(f)(a)\sqsubseteq\alpha(g)(a)$.
		\item Jeżeli $a\in\dom(f)$, to $(a,f(a))\in f\subseteq g$, więc
		$a\in\dom(g)$ oraz $g(a)=f(a)$. Stąd
		\[
		\alpha(f)(a)=f(a)=g(a)=\alpha(g)(a),
		\]
		a w szczególności $\alpha(f)(a)\sqsubseteq\alpha(g)(a)$.
	\end{itemize}
	Ponieważ $a$ było dowolne, z definicji porządku pionowego
	$\alpha(f)\sqsubseteq\alpha(g)$.
	
	($\Leftarrow$) Załóżmy, że $\alpha(f)\sqsubseteq\alpha(g)$, i niech
	$(a,b)\in f$. Wtedy $a\in\dom(f)$ oraz $f(a)=b$, więc
	$\alpha(f)(a)=b$, przy czym $b\neq\bt$, bo $b\in B$. Z założenia
	\[
	b=\alpha(f)(a)\sqsubseteq\alpha(g)(a).
	\]
	W porządku płaskim element różny od $\bt$ jest mniejszy lub równy
	jedynie samemu sobie (Definicja~\ref{def:zbior_plaski}), zatem
	$\alpha(g)(a)=b\neq\bt$. Oznacza to, że $a\in\dom(g)$ oraz
	$g(a)=\alpha(g)(a)=b$, czyli $(a,b)\in g$. Ponieważ $(a,b)\in f$ było
	dowolne, $f\subseteq g$.
	
	\emph{Element najmniejszy.} Funkcja pusta nie jest nigdzie określona,
	tzn.\ $\dom(\varnothing)=\varnothing$. Zatem dla każdego $a\in A$
	mamy $a\notin\dom(\varnothing)$ i
	\[
	\alpha(\varnothing)(a)=\bt=\varphi(a),
	\]
	czyli $\alpha(\varnothing)=\varphi$.
	
	Wykazaliśmy, że $\alpha$ jest bijekcją, zachowuje porządek w obie
	strony oraz przeprowadza element najmniejszy $\varnothing$ na element
	najmniejszy $\varphi$. Zatem
	\[
	(A\rightharpoonup B,\subseteq,\varnothing)
	\cong
	(A\to B_\bt,\sqsubseteq,\varphi).
	\]
\end{proof}

\begin{wniosek}\label{wn:izomorfizm}
	Każdą funkcję częściową możemy utożsamić z odpowiednią funkcją całkowitą.	
\end{wniosek}

\section{Funkcje w zbiorach łańcuchowo zupełnych}

\subsection{Monotoniczność i ciągłość}

Do przejścia do sedna naszych rozważań będziemy potrzebować jeszcze dwóch kluczowych pojęć.

\begin{definicja}
	Funkcję całkowitą $f : \boldsymbol{A} \to \boldsymbol{B}$ nazywamy \emph{monotoniczną}, jeżeli zachodzi
	\[
	a_1 \sqsubseteq_A a_2
	\implies
	f(a_1) \sqsubseteq_B f(a_2)
	\]
	dla dowolnych $a_1, a_2 \in A$.
\end{definicja}

Analogicznie wprowadzamy definicję ciągłości funkcji w zbiorze łańcuchowo zupełnym.

\begin{definicja}
	Funkcję całkowitą $f : \boldsymbol{A} \to \boldsymbol{B}$ nazywamy \emph{ciągłą}, jeżeli przenosi ona łańcuchy wraz z ich granicami z $\boldsymbol{A}$ do $\boldsymbol{B}$, tzn.
	\begin{itemize}
		\item [(1)] dla każdego łańcucha
		\[
		a_1 \sqsubseteq_A a_2 \sqsubseteq_A a_3 \sqsubseteq_A \ldots
		\]
		w $\boldsymbol{A}$ jego obraz
		\[
		f(a_1), f(a_2), f(a_3), \ldots
		\]
		jest łańcuchem w $\boldsymbol{B}$, tzn:
		\[
			f(a_1) \sqsubseteq_B f(a_2) \sqsubseteq_B f(a_3) \sqsubseteq_B \ldots
		\]
		\item [(2)] funkcja $f$ zachowuje granice łańcuchów, tzn.
		\[
		f\left(\bigsqcup_{i=1}^{\infty} a_i\right)
		=
		\bigsqcup_{i=1}^{\infty} f(a_i).
		\]
	\end{itemize}
\end{definicja}

\begin{twierdzenie}\label{tw:ciagla_monotoniczna}
	Każda funkcja ciągła jest monotoniczna.
\end{twierdzenie}

\begin{proof}
	Niech $f : \boldsymbol{A} \to \boldsymbol{B}$ będzie funkcją ciągłą. Ustalmy dowolne
	$a_1,a_2\in A$ takie, że
	\[
	a_1\sqsubseteq_A a_2.
	\]
	Wówczas ciąg
	\[
	a_1,a_2,a_2,a_2,\ldots
	\]
	jest łańcuchem w $\boldsymbol{A}$, którego granicą jest $a_2$, tzn.
	\[
	\bigsqcup_{i=1}^{\infty} a_i = a_2.
	\]
	Z ciągłości funkcji $f$ wynika, że ciąg
	\[
	f(a_1),f(a_2),f(a_2),f(a_2),\ldots
	\]
	jest łańcuchem w $\boldsymbol{B}$ oraz
	\[
	f\left(\bigsqcup_{i=1}^{\infty} a_i\right)
	=
	\bigsqcup_{i=1}^{\infty} f(a_i).
	\]
	Zatem
	\[
	f(a_2)
	=
	\bigsqcup_{i=1}^{\infty} f(a_i),
	\]
	a więc w szczególności
	\[
	f(a_1)\sqsubseteq_B f(a_2).
	\]
	Stąd
	\[
	a_1\sqsubseteq_A a_2
	\implies
	f(a_1)\sqsubseteq_B f(a_2),
	\]
	co oznacza, że $f$ jest monotoniczna.
\end{proof}

Twierdzenie~\ref{tw:ciagla_monotoniczna} mówi, że z ciągłości wynika
monotoniczność. Zapytajmy teraz w drugą stronę: co wynika z samej
monotoniczności? Okazuje się, że funkcja monotoniczna spełnia
automatycznie warunek~(1) definicji ciągłości oraz jedną z dwóch
nierówności składających się na równość z warunku~(2).

\begin{lemat}\label{lem:monotoniczna_lancuchy}
	Niech $\boldsymbol{A}$ i $\boldsymbol{B}$ będą zbiorami łańcuchowo
	zupełnymi, a $f:\boldsymbol{A}\to\boldsymbol{B}$ funkcją monotoniczną.
	Niech ponadto
	\[
	a_1\sqsubseteq_A a_2\sqsubseteq_A a_3\sqsubseteq_A\ldots
	\]
	będzie dowolnym łańcuchem w $\boldsymbol{A}$. Wówczas:
	\begin{itemize}
		\item[(a)] ciąg $f(a_1),f(a_2),f(a_3),\ldots$ jest łańcuchem
		w $\boldsymbol{B}$, a zatem posiada kres górny;
		\item[(b)] zachodzi nierówność
		\[
		\bigsqcup_{n=1}^{\infty}f(a_n)\sqsubseteq_B f\Big(\bigsqcup_{n=1}^{\infty}a_n\Big).
		\]
	\end{itemize}
\end{lemat}

\begin{proof}
	\emph{(a)} Ustalmy $n\in\mathbb N$. Ponieważ $(a_n)_{n\in\mathbb N}$
	jest łańcuchem, mamy
	\[
	a_n\sqsubseteq_A a_{n+1}.
	\]
	Z monotoniczności funkcji $f$ otrzymujemy
	\[
	f(a_n)\sqsubseteq_B f(a_{n+1}).
	\]
	Ponieważ $n$ było dowolne,
	\[
	f(a_1)\sqsubseteq_B f(a_2)\sqsubseteq_B f(a_3)\sqsubseteq_B\ldots,
	\]
	czyli ciąg $(f(a_n))_{n\in\mathbb N}$ jest łańcuchem w $\boldsymbol{B}$.
	Ponieważ $\boldsymbol{B}$ jest zbiorem łańcuchowo zupełnym, łańcuch ten
	posiada kres górny
	\[
	\bigsqcup_{n=1}^{\infty}f(a_n)\in B.
	\]
	
	\emph{(b)} Oznaczmy kres górny łańcucha $(a_n)_{n\in\mathbb N}$ przez
	\[
	a=\bigsqcup_{n=1}^{\infty}a_n.
	\]
	Z definicji kresu górnego $a$ jest ograniczeniem górnym tego łańcucha,
	tzn.
	\[
	a_n\sqsubseteq_A a
	\qquad(n\in\mathbb N).
	\]
	Z monotoniczności funkcji $f$ wynika zatem
	\[
	f(a_n)\sqsubseteq_B f(a)
	\qquad(n\in\mathbb N).
	\]
	Oznacza to, że $f(a)$ jest ograniczeniem górnym łańcucha
	$(f(a_n))_{n\in\mathbb N}$ z punktu (a). Ponieważ kres górny jest
	najmniejszym ograniczeniem górnym, otrzymujemy
	\[
	\bigsqcup_{n=1}^{\infty}f(a_n)
	\sqsubseteq_B
	f(a)
	=
	f\Big(\bigsqcup_{n=1}^{\infty}a_n\Big).
	\]
\end{proof}

Z lematu wynika, że przy badaniu ciągłości funkcji monotonicznej nie
trzeba sprawdzać całej równości z warunku~(2) --- wystarczy jedna
nierówność, bo druga zachodzi zawsze.

\begin{wniosek}\label{wn:ciaglosc_monotonicznej}
	Niech $\boldsymbol{A}$ i $\boldsymbol{B}$ będą zbiorami łańcuchowo
	zupełnymi. Funkcja monotoniczna $f:\boldsymbol{A}\to\boldsymbol{B}$
	jest ciągła wtedy i tylko wtedy, gdy dla każdego łańcucha
	$(a_n)_{n\in\mathbb N}$ w $\boldsymbol{A}$ zachodzi
	\[
	f\Big(\bigsqcup_{n=1}^{\infty}a_n\Big)\sqsubseteq_B\bigsqcup_{n=1}^{\infty}f(a_n).
	\]
\end{wniosek}

\begin{proof}
	($\Rightarrow$) Jeżeli $f$ jest ciągła, to z warunku~(2) definicji
	ciągłości zachodzi równość
	$f\big(\bigsqcup_{n}a_n\big)=\bigsqcup_{n}f(a_n)$, a więc
	w szczególności żądana nierówność.
	
	($\Leftarrow$) Załóżmy, że $f$ jest monotoniczna i spełnia powyższą
	nierówność dla każdego łańcucha. Sprawdzimy oba warunki definicji
	ciągłości.
	
	Warunek~(1): na mocy Lematu~\ref{lem:monotoniczna_lancuchy} (a) obraz
	dowolnego łańcucha w $\boldsymbol{A}$ jest łańcuchem w $\boldsymbol{B}$.
	
	Warunek~(2): ustalmy łańcuch $(a_n)_{n\in\mathbb N}$. Z założenia
	\[
	f\Big(\bigsqcup_{n=1}^{\infty}a_n\Big)\sqsubseteq_B\bigsqcup_{n=1}^{\infty}f(a_n),
	\]
	a z Lematu~\ref{lem:monotoniczna_lancuchy} (b)
	\[
	\bigsqcup_{n=1}^{\infty}f(a_n)\sqsubseteq_B f\Big(\bigsqcup_{n=1}^{\infty}a_n\Big).
	\]
	Z antysymetrii porządku $\sqsubseteq_B$ otrzymujemy
	\[
	f\Big(\bigsqcup_{n=1}^{\infty}a_n\Big)=\bigsqcup_{n=1}^{\infty}f(a_n).
	\]
	Zatem $f$ spełnia oba warunki definicji ciągłości, czyli jest ciągła.
\end{proof}

\subsection{Przestrzeń funkcji ciągłych}

\begin{definicja}
	Niech $\boldsymbol{A}$ i $\boldsymbol{B}$ będą zbiorami łańcuchowo
	zupełnymi. Zbiór wszystkich funkcji ciągłych z $\boldsymbol{A}$ do
	$\boldsymbol{B}$ oznaczamy przez
	\[
	\mathcal{C}(\boldsymbol{A}\to\boldsymbol{B})
	=
	\{f:\boldsymbol{A}\to\boldsymbol{B}\mid f\text{ jest ciągła}\}.
	\]
	Na zbiorze tym rozważamy porządek pionowy określony przez
	\[
	f\sqsubseteq g
	\iff
	\forall a\in A\quad
	f(a)\sqsubseteq_B g(a).
	\]
\end{definicja}

\begin{twierdzenie}\label{tw:funkcje_ciagle_cpo}
	Niech $\boldsymbol{A}$ i $\boldsymbol{B}$ będą zbiorami łańcuchowo
	zupełnymi. Wówczas
	\[
	\mathcal{C}(\boldsymbol{A}\to\boldsymbol{B})
	\]
	z porządkiem pionowym jest zbiorem łańcuchowo zupełnym.
	Elementem najmniejszym jest funkcja stała
	\[
	\varphi(a)=\bt_B,
	\]
	a kres górny każdego łańcucha funkcji ciągłych wyznaczamy punktowo:
	\[
	\left(\bigsqcup_{n=1}^{\infty}f_n\right)(a)
	=
	\bigsqcup_{n=1}^{\infty}f_n(a).
	\]
\end{twierdzenie}

\begin{proof}
	\emph{Częściowy porządek.} Zbiór
	$\mathcal{C}(\boldsymbol{A}\to\boldsymbol{B})$ jest podzbiorem zbioru
	$A\to B$, a relacja $\sqsubseteq$ na nim jest obcięciem porządku
	pionowego. Zwrotność, przechodniość i antysymetryczność zachodzą dla
	wszystkich funkcji z $A\to B$
	(Twierdzenie~\ref{tw:zbior_funkcji_calkowitych_czesciowo_uporzadkowany}),
	więc tym bardziej dla funkcji ciągłych.
	
	\emph{Kres górny łańcucha.} Niech
	\[
	f_1\sqsubseteq f_2\sqsubseteq f_3\sqsubseteq\ldots
	\]
	będzie łańcuchem w $\mathcal{C}(\boldsymbol{A}\to\boldsymbol{B})$.
	Jest to w szczególności łańcuch w $\boldsymbol{A\to B}$, więc na mocy
	Twierdzenia~\ref{tw:funkcje_calkowite_cpo} posiada on kres górny
	w $\boldsymbol{A\to B}$, równy funkcji
	\[
	f(a)=\bigsqcup_{n=1}^{\infty}f_n(a)
	\qquad(a\in A).
	\]
	Kres ten został jednak wyznaczony w większym zbiorze $A\to B$.
	Pokażemy, że jeżeli $f$ jest ciągła, to jest również kresem górnym
	łańcucha w mniejszym zbiorze $\mathcal{C}(\boldsymbol{A}\to\boldsymbol{B})$.
	\begin{itemize}
		\item $f$ jest ograniczeniem górnym łańcucha, bo $f_n\sqsubseteq f$
		dla każdego $n$ (jest kresem w $A\to B$).
		\item Niech $g\in\mathcal{C}(\boldsymbol{A}\to\boldsymbol{B})$ będzie
		dowolnym ograniczeniem górnym łańcucha. Wtedy $g\in A\to B$ oraz
		$f_n\sqsubseteq g$ dla każdego $n$, więc $g$ jest ograniczeniem
		górnym łańcucha także w $A\to B$. Ponieważ $f$ jest najmniejszym
		ograniczeniem górnym w $A\to B$, otrzymujemy $f\sqsubseteq g$.
	\end{itemize}
	Pozostaje zatem wykazać, że $f$ jest ciągła. Na mocy
	Wniosku~\ref{wn:ciaglosc_monotonicznej} wystarczy pokazać, że
	\begin{itemize}
		\item[(I)] $f$ jest monotoniczna,
		\item[(II)] dla każdego łańcucha $(a_n)_{n\in\mathbb N}$ w
		$\boldsymbol{A}$ zachodzi
		$f\big(\bigsqcup_{n}a_n\big)\sqsubseteq_B\bigsqcup_{n}f(a_n)$.
	\end{itemize}
	
	\emph{(I) Monotoniczność $f$.} Niech $a\sqsubseteq_A a'$ i ustalmy
	$i\in\mathbb N$. Funkcja $f_i$ jest ciągła, a więc na mocy
	Twierdzenia~\ref{tw:ciagla_monotoniczna} monotoniczna, skąd
	\[
	f_i(a)\sqsubseteq_B f_i(a').
	\]
	Ponadto $f_i\sqsubseteq f$ w porządku pionowym, czyli w szczególności
	\[
	f_i(a')\sqsubseteq_B f(a').
	\]
	Z przechodniości otrzymujemy
	\[
	f_i(a)\sqsubseteq_B f(a').
	\]
	Ponieważ $i$ było dowolne, element $f(a')$ jest ograniczeniem górnym
	łańcucha $\big(f_i(a)\big)_{i\in\mathbb N}$. Kresem górnym tego
	łańcucha jest z definicji $f(a)$, a kres jest najmniejszym
	ograniczeniem górnym, więc
	\[
	f(a)\sqsubseteq_B f(a').
	\]
	
	\emph{(II) Nierówność dla kresów.} Niech $(a_n)_{n\in\mathbb N}$ będzie
	łańcuchem w $\boldsymbol{A}$ i niech
	\[
	a=\bigsqcup_{n=1}^{\infty}a_n.
	\]
	Na mocy (I) oraz Lematu~\ref{lem:monotoniczna_lancuchy}(a) ciąg
	$\big(f(a_n)\big)_{n\in\mathbb N}$ jest łańcuchem w $\boldsymbol{B}$,
	więc istnieje jego kres górny
	\[
	s=\bigsqcup_{n=1}^{\infty}f(a_n).
	\]
	Ustalmy $i\in\mathbb N$. Dla każdego $n\in\mathbb N$ mamy
	\[
	f_i(a_n)\sqsubseteq_B f(a_n)\sqsubseteq_B s,
	\]
	gdzie pierwsza nierówność wynika z $f_i\sqsubseteq f$, a druga z tego,
	że $s$ jest ograniczeniem górnym łańcucha $\big(f(a_n)\big)_n$.
	Zatem $s$ jest ograniczeniem górnym łańcucha
	$\big(f_i(a_n)\big)_{n\in\mathbb N}$, skąd, korzystając z ciągłości
	$f_i$,
	\[
	f_i(a)
	=
	f_i\Big(\bigsqcup_{n=1}^{\infty}a_n\Big)
	=
	\bigsqcup_{n=1}^{\infty}f_i(a_n)
	\sqsubseteq_B
	s.
	\]
	Ponieważ $i$ było dowolne, $s$ jest ograniczeniem górnym łańcucha
	$\big(f_i(a)\big)_{i\in\mathbb N}$, którego kresem jest $f(a)$. Stąd
	\[
	f(a)
	=
	\bigsqcup_{i=1}^{\infty}f_i(a)
	\sqsubseteq_B
	s
	=
	\bigsqcup_{n=1}^{\infty}f(a_n).
	\]
	Na mocy Wniosku~\ref{wn:ciaglosc_monotonicznej} funkcja $f$ jest
	ciągła, a więc $f$ jest kresem górnym łańcucha
	$(f_n)_{n\in\mathbb N}$ w $\mathcal{C}(\boldsymbol{A}\to\boldsymbol{B})$.
	
	\emph{Element najmniejszy.} Pokażemy najpierw, że funkcja stała
	$\varphi(a)=\bt_B$ jest ciągła. Jest ona monotoniczna, bo
	$\varphi(a)=\bt_B\sqsubseteq_B\bt_B=\varphi(a')$ dla dowolnych $a,a'$.
	Dla łańcucha $(a_n)_{n\in\mathbb N}$ ciąg $\varphi(a_n)$ jest stały
	i równy $\bt_B$, więc jego kresem jest $\bt_B$ i
	\[
	\varphi\Big(\bigsqcup_{n=1}^{\infty}a_n\Big)
	=
	\bt_B
	=
	\bigsqcup_{n=1}^{\infty}\varphi(a_n).
	\]
	Zatem $\varphi\in\mathcal{C}(\boldsymbol{A}\to\boldsymbol{B})$.
	Ponadto na mocy Twierdzenia~\ref{tw:funkcje_calkowite_cpo} funkcja
	$\varphi$ jest elementem najmniejszym całego zbioru $A\to B$, więc
	$\varphi\sqsubseteq f$ dla każdej funkcji ciągłej $f$. Zatem $\varphi$
	jest elementem najmniejszym zbioru
	$\mathcal{C}(\boldsymbol{A}\to\boldsymbol{B})$.
	
	Wykazaliśmy, że $\mathcal{C}(\boldsymbol{A}\to\boldsymbol{B})$ jest
	częściowo uporządkowany, każdy jego łańcuch ma kres górny wyznaczany
	punktowo i istnieje element najmniejszy $\varphi$. Jest to więc zbiór
	łańcuchowo zupełny.
\end{proof}

\subsection{Funkcje na zbiorach płaskich}

\begin{twierdzenie}\label{tw:monotonna_plaska_ciagla}
	Niech $\boldsymbol{A}$ będzie płaskim zbiorem łańcuchowo zupełnym,
	a $\boldsymbol{B}$ dowolnym zbiorem łańcuchowo zupełnym. Wówczas każda
	monotoniczna funkcja
	\[
	f:\boldsymbol{A}\to\boldsymbol{B}
	\]
	jest ciągła.
\end{twierdzenie}

\begin{proof}
	Niech
	\[
	a_1\sqsubseteq_A a_2\sqsubseteq_A a_3\sqsubseteq_A\ldots
	\]
	będzie dowolnym łańcuchem w $\boldsymbol{A}$. Ponieważ $f$ jest
	monotoniczna, na mocy Lematu~\ref{lem:monotoniczna_lancuchy}(a) ciąg
	$\big(f(a_n)\big)_{n\in\mathbb N}$ jest łańcuchem w $\boldsymbol{B}$,
	a na mocy Wniosku~\ref{wn:ciaglosc_monotonicznej} wystarczy wykazać
	nierówność
	\[
	f\Big(\bigsqcup_{n=1}^{\infty}a_n\Big)
	\sqsubseteq_B
	\bigsqcup_{n=1}^{\infty}f(a_n).
	\]
	
	Pokażemy najpierw, że istnieje indeks $k\in\mathbb N$ taki, że
	\[
	\bigsqcup_{n=1}^{\infty}a_n=a_k.
	\]
	Rozważmy dwa przypadki (por.\ Uwaga~\ref{uw:lancuchy_plaskie}).
	\begin{itemize}
		\item Jeżeli $a_n=\bt_A$ dla wszystkich $n$, to łańcuch jest stały,
		jego kresem jest $\bt_A$ i możemy przyjąć $k=1$.
		\item W przeciwnym razie niech $k$ będzie najmniejszym indeksem
		takim, że $a_k\neq\bt_A$. Dla $n\geq k$ mamy $a_k\sqsubseteq_A a_n$,
		a ponieważ w zbiorze płaskim element różny od $\bt_A$ jest mniejszy
		lub równy tylko samemu sobie, otrzymujemy $a_n=a_k$. Zatem łańcuch
		ma postać
		\[
		\bt_A,\ldots,\bt_A,a_k,a_k,a_k,\ldots
		\]
		Element $a_k$ jest jego ograniczeniem górnym (bo $\bt_A\sqsubseteq_A a_k$
		i $a_k\sqsubseteq_A a_k$), a każde ograniczenie górne $u$ spełnia
		w szczególności $a_k\sqsubseteq_A u$. Zatem
		$\bigsqcup_{n}a_n=a_k$.
	\end{itemize}
	
	W obu przypadkach
	\[
	f\Big(\bigsqcup_{n=1}^{\infty}a_n\Big)=f(a_k).
	\]
	Kres górny $\bigsqcup_{n}f(a_n)$ jest ograniczeniem górnym wszystkich
	wyrazów łańcucha $\big(f(a_n)\big)_{n\in\mathbb N}$, w szczególności
	wyrazu $f(a_k)$. Stąd
	\[
	f\Big(\bigsqcup_{n=1}^{\infty}a_n\Big)
	=
	f(a_k)
	\sqsubseteq_B
	\bigsqcup_{n=1}^{\infty}f(a_n),
	\]
	co na mocy Wniosku~\ref{wn:ciaglosc_monotonicznej} kończy dowód
	ciągłości $f$.
\end{proof}

\begin{wniosek}\label{wn:plaskie_monotoniczne_ciagle}
	Niech $\boldsymbol{A}$ będzie płaskim zbiorem łańcuchowo zupełnym,
	a $\boldsymbol{B}$ dowolnym zbiorem łańcuchowo zupełnym. Jeżeli funkcja
	\[
	f:\boldsymbol{A}\to\boldsymbol{B}
	\]
	zachowuje elementy najmniejsze, tj.
	\[
	f(\bt_A)=\bt_B,
	\]
	to $f$ jest monotoniczna, a w konsekwencji jest ciągła.
\end{wniosek}

\begin{proof}
	Pokażemy, że $f$ jest monotoniczna. Niech $a_1,a_2\in A$ oraz
	\[
	a_1\sqsubseteq_A a_2.
	\]
	Rozważmy dwa przypadki.
	\begin{itemize}
		\item Jeżeli $a_1=\bt_A$, to z założenia $f(a_1)=f(\bt_A)=\bt_B$,
		a ponieważ $\bt_B$ jest elementem najmniejszym $\boldsymbol{B}$,
		\[
		f(a_1)=\bt_B\sqsubseteq_B f(a_2).
		\]
		\item Jeżeli $a_1\neq\bt_A$, to z płaskości $\boldsymbol{A}$
		(Definicja~\ref{def:zbior_plaski}) warunek $a_1\sqsubseteq_A a_2$
		oznacza $a_1=a_2$. Wtedy $f(a_1)=f(a_2)$, a więc z zwrotności
		$f(a_1)\sqsubseteq_B f(a_2)$.
	\end{itemize}
	W obu przypadkach $f(a_1)\sqsubseteq_B f(a_2)$, czyli $f$ jest
	monotoniczna. Ponieważ $\boldsymbol{A}$ jest płaski, na mocy
	Twierdzenia~\ref{tw:monotonna_plaska_ciagla} funkcja $f$ jest ciągła.
\end{proof}

\begin{wniosek}\label{wn:funkcje_czesciowe_jako_ciagle}
	Niech $A$ i $B$ będą dowolnymi zbiorami, a
	\[
	A_\bt=A\cup\{\bt_A\},
	\qquad
	B_\bt=B\cup\{\bt_B\}
	\]
	będą płaskimi zbiorami łańcuchowo zupełnymi otrzymanymi przez dodanie
	elementów najmniejszych. Każda funkcja częściowa
	\[
	f:A\rightharpoonup B
	\]
	wyznacza \emph{kanoniczne rozszerzenie}
	\[
	\widehat{f}:A_\bt\to B_\bt
	\]
	określone wzorem
	\[
	\widehat{f}(a)=
	\begin{cases}
		f(a),&a\in\dom(f),\\
		\bt_B,&a\in A\setminus\dom(f),\\
		\bt_B,&a=\bt_A.
	\end{cases}
	\]
	Funkcja $\widehat{f}$ zachowuje element najmniejszy, a zatem jest
	ciągła.
\end{wniosek}

\begin{proof}
	Najpierw zauważmy, że $\widehat{f}$ jest poprawnie określoną funkcją
	całkowitą z $A_\bt$ w $B_\bt$: każdy element $a\in A_\bt$ spełnia
	dokładnie jeden z warunków $a\in\dom(f)$, $a\in A\setminus\dom(f)$,
	$a=\bt_A$, a w każdym przypadku przypisana wartość należy do
	$B_\bt$ (bo $f(a)\in B$).
	
	Funkcja $\widehat{f}$ jest rozszerzeniem $f$, bo z pierwszego wiersza
	definicji $\widehat{f}(a)=f(a)$ dla każdego $a\in\dom(f)$. Z trzeciego
	wiersza definicji
	\[
	\widehat{f}(\bt_A)=\bt_B,
	\]
	czyli $\widehat{f}$ zachowuje element najmniejszy.
	
	Zbiór $A_\bt$ jest płaskim zbiorem łańcuchowo zupełnym, a $B_\bt$ ---
	zbiorem łańcuchowo zupełnym (Uwaga~\ref{uw:lancuchy_plaskie}). Zatem
	na mocy Wniosku~\ref{wn:plaskie_monotoniczne_ciagle} funkcja
	$\widehat{f}$ jest monotoniczna i ciągła.
	
	Zauważmy jeszcze, że wartości $\widehat{f}$ poza $\dom(f)$ nie są
	wymuszone przez $f$ --- przyjęcie $\bt_B$ wyraża właśnie to, że $f$
	nie jest tam określona. Na $A$ funkcja $\widehat{f}$ pokrywa się
	z funkcją $\alpha(f)$ z Twierdzenia~\ref{tw:izomorfizm_czesciowe_calkowite}.
\end{proof}

%==============================================================================
\section{Równania stałopunktowe w zbiorach łańcuchowo zupełnych}
%==============================================================================

Niniejszy rozdział poświęcimy w pełni zagadnieniu znajdowania punktu stałego
dowolnego odwzorowania ciągłego między zbiorami łańcuchowo zupełnymi. Wrócimy na
chwilę do równania~\eqref{eq:rownanie_stalopunktowe} i postarajmy się odpowiedzieć na
zadane wcześniej pytania, tzn.\ o istnienie rozwiązania równania stałopunktowego oraz
o wybór właściwego rozwiązania, gdy jest ich więcej. Dysponujemy już wszystkimi pojęciami --- zbiorem łańcuchowo zupełnym oraz funkcją ciągłą --- potrzebnymi, aby wykazać istnienie rozwiązania równania~\eqref{eq:rownanie_stalopunktowe} i wskazać wśród rozwiązań jedno wyróżnione: najmniejszy punkt stały. Twierdzenie to nie orzeka natomiast, że rozwiązanie jest jedyne.

Przejdziemy teraz do głównego twierdzenia niniejszego referatu --- twierdzenia Kleenego\footnote{Stephen Cole Kleene (1909--1994) --- amerykański matematyk i~logik, jeden z~twórców teorii rekursji, czyli teorii funkcji obliczalnych. Doktorat uzyskał w~1934~r.\ w~Princeton pod kierunkiem Alonzo Churcha, a~większość kariery naukowej związał z~University of Wisconsin--Madison. Zajmował się m.in.\ funkcjami częściowo rekurencyjnymi, hierarchią arytmetyczną, realizowalnością w~logice intuicjonistycznej oraz wyrażeniami regularnymi --- od jego nazwiska pochodzi nazwa \emph{gwiazdki Kleenego} $A^*$, użytej w~Przykładzie~\ref{prz:alfabet}. Jest autorem klasycznej monografii \emph{Introduction to Metamathematics} (1952)~\cite{kleene1952}. Omawiane twierdzenie nawiązuje do jego \emph{pierwszego twierdzenia o~rekursji}, według którego definicja rekurencyjna wyznacza najmniejszy punkt stały odpowiedniego operatora; sformułowanie w~języku zbiorów łańcuchowo zupełnych pochodzi z~późniejszych prac nad semantyką denotacyjną.} --- łączącego teorię punktu stałego z wprowadzonym pojęciem zbioru łańcuchowo zupełnego.

\begin{twierdzenie}[Kleene]\label{tw:kleene}
	Każda funkcja ciągła $f:\boldsymbol{A}\to\boldsymbol{A}$ ma najmniejszy
	punkt stały. Co więcej, najmniejszy punkt stały jest kresem górnym
	łańcucha kolejnych iteracji tej funkcji na elemencie najmniejszym, tzn.
	\[
	\lfp(f)=\bigsqcup_{i=1}^\infty f^i(\bt),
	\]
	gdzie
	\[
	f^1(\bt)=f(\bt),
	\qquad
	f^{i+1}(\bt)=f\!\left(f^i(\bt)\right)
	\]
	dla $i=1,2,3,\ldots$.
\end{twierdzenie}

\begin{proof}
	Niech $\boldsymbol{A}=(A,\sqsubseteq,\bt)$ oraz niech
	$f:\boldsymbol{A}\to\boldsymbol{A}$ będzie funkcją ciągłą. Dowód przeprowadzimy w kilku kluczowych krokach.
	
	Najpierw pokażemy, że rodzina
	\[
	\left\{f^i(\bt)\right\}_{i=1}^{\infty}
	\]
	jest łańcuchem. Ponieważ $\bt$ jest elementem najmniejszym w $A$, mamy
	\[
	\bt\sqsubseteq f(\bt).
	\]
	Ponieważ $f$ jest ciągła, wobec Twierdzenia~\ref{tw:ciagla_monotoniczna} jest monotoniczna, a więc
	\[
		f(\bt)\sqsubseteq f(f(\bt))=f^2(\bt).
	\]
	Załóżmy indukcyjnie, że dla pewnego $i\in\mathbb{N}$ zachodzi
	\[
	f^i(\bt)\sqsubseteq f^{i+1}(\bt).
	\]
	Wobec Twierdzenia~\ref{tw:ciagla_monotoniczna} ciągłość funkcji $f$
	implikuje jej monotoniczność, zatem
	\[
	f\left(f^i(\bt)\right)
	\sqsubseteq
	f\left(f^{i+1}(\bt)\right).
	\]
	Stąd
	\[
	f^{i+1}(\bt)\sqsubseteq f^{i+2}(\bt).
	\]
	Na mocy zasady indukcji matematycznej
	\[
	f^1(\bt)\sqsubseteq f^2(\bt)\sqsubseteq f^3(\bt)
	\sqsubseteq\ldots,
	\]
	czyli rodzina $\left\{f^i(\bt)\right\}_{i=1}^{\infty}$ jest łańcuchem.
	
	Z łańcuchowej zupełności $\boldsymbol{A}$ wynika istnienie kresu górnego
	łańcucha $\left\{f^i(\bt)\right\}_{i=1}^{\infty}$. Oznaczmy go przez
	\[
	x=\bigsqcup_{i=1}^{\infty}f^i(\bt).
	\]
	
	Ponieważ $f$ jest funkcją ciągłą, zachowuje kresy górne łańcuchów, zatem
	\[
		f(x)
		=f\left(\bigsqcup_{i=1}^{\infty}f^i(\bt)\right)
		=\bigsqcup_{i=1}^{\infty}f\left(f^i(\bt)\right)
		=\bigsqcup_{i=1}^{\infty}f^{i+1}(\bt)
		=\bigsqcup_{i=2}^{\infty}f^i(\bt).
	\]
	
	Ponieważ
	\[
	f^1(\bt)\sqsubseteq f^2(\bt)\sqsubseteq f^3(\bt)
	\sqsubseteq\ldots,
	\]
	element $f^1(\bt)$ jest elementem najmniejszym tego łańcucha.
	Dodanie go do łańcucha
	$\left\{f^i(\bt)\right\}_{i=2}^{\infty}$ nie zmienia jego kresu górnego.
	Stąd
	\[
	\bigsqcup_{i=2}^{\infty}f^i(\bt)
	=
	\bigsqcup_{i=1}^{\infty}f^i(\bt)
	=x.
	\]
	W konsekwencji
	\[
	f(x)=x,
	\]
	czyli $x$ jest punktem stałym odwzorowania $f$.
	
	Pozostaje wykazać, że $x$ jest najmniejszym punktem stałym odwzorowania $f$.
	Weźmy dowolny punkt stały $a\in\Fix(f)$, tj.
	\[
	f(a)=a.
	\]
	Ponieważ $\bt$ jest elementem najmniejszym w $A$, mamy
	\[
	\bt\sqsubseteq a.
	\]
	Z monotoniczności funkcji $f$ otrzymujemy
	\[
	f(\bt)\sqsubseteq f(a)=a,
	\]
	czyli
	\[
	f^1(\bt)\sqsubseteq a.
	\]
	
	Załóżmy indukcyjnie, że dla pewnego $i \in \mathbb{N}$ zachodzi
	\[
	f^i(\bt)\sqsubseteq a.
	\]
	Z monotoniczności funkcji $f$ otrzymujemy
	\[
	f\left(f^i(\bt)\right)\sqsubseteq f(a)=a,
	\]
	a więc
	\[
	f^{i+1}(\bt)\sqsubseteq a.
	\]
	Na mocy zasady indukcji matematycznej
	\[
	f^i(\bt)\sqsubseteq a
	\]
	dla każdego $i \in \mathbb{N}$.
	
	Zatem $a$ jest ograniczeniem górnym łańcucha
	$\left\{f^i(\bt)\right\}_{i=1}^{\infty}$. Ponieważ
	\[
	x=\bigsqcup_{i=1}^{\infty}f^i(\bt)
	\]
	jest jego najmniejszym ograniczeniem górnym, otrzymujemy
	\[
	x\sqsubseteq a.
	\]
	
	Ponieważ $a$ było dowolnym punktem stałym odwzorowania $f$, otrzymujemy,
	że $x$ jest elementem najmniejszym zbioru $\Fix(f)$. Zatem
	\[
	x=\lfp(f).
	\]
\end{proof}

\subsection{Sformułowanie równania stałopunktowego}

Wróćmy teraz do zagadnienia semantyki pętli \kw{while}, rozważanego
w podrozdziale~\ref{subsec:problem_while}. Przypomnijmy, że denotację
instrukcji reprezentujemy jako funkcję częściową
\[
\Cc\sem{c}:\Store\rightharpoonup\Store.
\]
Jak wyjaśniono wcześniej, częściowość tej funkcji pozwala uwzględnić
możliwość niezakończenia wykonania instrukcji. Jeżeli dla pewnego stanu
$\sigma$ nie istnieje stan $\sigma'$ taki, że
\[
(\sigma,\sigma')\in\Cc\sem{c},
\]
oznacza to, że wykonanie instrukcji $c$ rozpoczęte w stanie $\sigma$
nie kończy się.

W szczególności również denotację pętli
\[
\kw{while}\ b\ \kw{do}\ c
\]
będziemy szukać w przestrzeni funkcji częściowych
\[
\Store\rightharpoonup\Store.
\]
Oznaczmy przez
\[
W_{b,c}
=
\Cc\sem{\kw{while}\ b\ \kw{do}\ c}
\]
denotację tej pętli.

W podrozdziale~\ref{subsec:problem_while} otrzymaliśmy równanie
opisujące $W_{b,c}$, w którym denotacja pętli występuje po obu stronach
równania. Wynika to bezpośrednio z charakteru instrukcji \kw{while}:
jeżeli warunek $b$ jest prawdziwy, po wykonaniu ciała $c$ należy ponownie
wykonać całą pętlę.

Aby formalnie rozwiązać to równanie, wprowadzimy funkcjonał
\[
\mathcal{F}:
(\Store\rightharpoonup\Store)
\longrightarrow
(\Store\rightharpoonup\Store).
\]
Argumentem funkcjonału $\mathcal{F}$ będzie dowolna funkcja częściowa
$f:\Store\rightharpoonup\Store$, która reprezentuje proponowaną
denotację dalszego wykonania pętli. Wartość $\mathcal{F}(f)$ opisuje
jedno rozwinięcie pętli: sprawdzenie warunku $b$, wykonanie ciała $c$
w przypadku prawdziwego warunku oraz dalsze wykonanie opisane przez
$f$.

Funkcjonał $\mathcal{F}$ definiujemy następująco:
\[
\begin{aligned}
	\mathcal{F}(f)
	={}&
	\left\{
	(\sigma,\sigma)
	\;\middle|\;
	(\sigma,\kw{false})\in\Bb\sem{b}
	\right\}\\
	&{}\cup
	\left\{
	(\sigma,\sigma')
	\;\middle|\;
	\begin{array}{c}
		(\sigma,\kw{true})\in\Bb\sem{b},\\
		\exists\sigma''\in\Store\;
		\big(
		(\sigma,\sigma'')\in\Cc\sem{c}
		\land
		(\sigma'',\sigma')\in f
		\big)
	\end{array}
	\right\}.
\end{aligned}
\]

Pierwszy składnik opisuje sytuację, w której warunek $b$ jest fałszywy.
Wówczas wykonanie pętli kończy się bez wykonania jej ciała, dlatego stan
końcowy jest równy stanowi początkowemu:
\[
(\sigma,\sigma)\in\mathcal{F}(f).
\]

Drugi składnik opisuje sytuację, w której warunek $b$ jest prawdziwy.
Najpierw wykonywana jest instrukcja $c$, prowadząc ze stanu $\sigma$
do stanu pośredniego $\sigma''$. Następnie dalsze wykonanie pętli jest
reprezentowane przez funkcję $f$, która prowadzi ze stanu $\sigma''$ do
stanu końcowego $\sigma'$. Wówczas
\[
(\sigma,\sigma')\in\mathcal{F}(f).
\]

Jeżeli teraz jako argument funkcjonału przyjmiemy właściwą denotację
pętli, czyli $f=W_{b,c}$, otrzymujemy
\[
\begin{aligned}
	\mathcal{F}(W_{b,c})
	={}&
	\left\{
	(\sigma,\sigma)
	\;\middle|\;
	(\sigma,\kw{false})\in\Bb\sem{b}
	\right\}\\
	&{}\cup
	\left\{
	(\sigma,\sigma')
	\;\middle|\;
	\begin{array}{c}
		(\sigma,\kw{true})\in\Bb\sem{b},\\
		\exists\sigma''\in\Store\;
		\big(
		(\sigma,\sigma'')\in\Cc\sem{c}
		\land
		(\sigma'',\sigma')\in W_{b,c}
		\big)
	\end{array}
	\right\}.
\end{aligned}
\]

Zgodnie z równaniem opisującym działanie pętli otrzymujemy zatem
\[
W_{b,c}=\mathcal{F}(W_{b,c}).
\]
Denotacja $W_{b,c}$ jest więc punktem stałym funkcjonału $\mathcal F$.

Nie każdy punkt stały musi jednak odpowiadać właściwej denotacji pętli.
Rozważmy na przykład pętlę
\[
\kw{while}\ \kw{true}\ \kw{do}\ \kw{skip},
\]
czyli $b=\kw{true}$ oraz $c=\kw{skip}$. Obliczmy funkcjonał
$\mathcal F$ dla tej pętli. Z definicji denotacji stałej logicznej
\[
\Bb\sem{\kw{true}}=\{(\sigma,\kw{true})\mid\sigma\in\Store\},
\]
więc nie istnieje stan $\sigma$, dla którego
$(\sigma,\kw{false})\in\Bb\sem{\kw{true}}$. Pierwszy składnik w
definicji $\mathcal F(f)$ jest zatem zbiorem pustym:
\[
\left\{
(\sigma,\sigma)
\;\middle|\;
(\sigma,\kw{false})\in\Bb\sem{\kw{true}}
\right\}
=\varnothing.
\]
W drugim składniku warunek $(\sigma,\kw{true})\in\Bb\sem{\kw{true}}$
jest spełniony dla każdego $\sigma$. Ponadto
\[
\Cc\sem{\kw{skip}}=\{(\sigma,\sigma)\mid\sigma\in\Store\},
\]
więc warunek $(\sigma,\sigma'')\in\Cc\sem{\kw{skip}}$ oznacza dokładnie
$\sigma''=\sigma$. Wobec tego
\[
\exists\sigma''\in\Store\;
\big(
(\sigma,\sigma'')\in\Cc\sem{\kw{skip}}
\land
(\sigma'',\sigma')\in f
\big)
\iff
(\sigma,\sigma')\in f.
\]
Otrzymujemy
\[
\mathcal F(f)
=
\varnothing
\cup
\left\{(\sigma,\sigma')\;\middle|\;(\sigma,\sigma')\in f\right\}
=
f
\]
dla każdej funkcji częściowej $f:\Store\rightharpoonup\Store$. Zatem
\emph{każda} funkcja częściowa jest punktem stałym funkcjonału
$\mathcal F$ --- na przykład funkcja pusta $\varnothing$, ale także
identyczność $\{(\sigma,\sigma)\mid\sigma\in\Store\}$.

Które z tych rozwiązań jest właściwe? Pętla
$\kw{while}\ \kw{true}\ \kw{do}\ \kw{skip}$ nie kończy się dla żadnego
stanu początkowego, więc jej denotacja nie powinna zawierać żadnej pary
$(\sigma,\sigma')$ --- powinna to być funkcja pusta $\varnothing$.
Identyczność byłaby błędna: twierdziłaby, że pętla kończy się
natychmiast, nie zmieniając stanu. Funkcja pusta jest przy tym
\emph{najmniejszym} punktem stałym, bo $\varnothing\subseteq f$ dla
każdej $f$. Ogólnie, najmniejszy punkt stały zawiera tylko te pary,
których istnienie wymusza równanie, a więc nie przypisuje pętli wyników,
których obliczenie w rzeczywistości nie osiąga.

Dlatego za denotację pętli przyjmujemy najmniejszy punkt stały względem
porządku inkluzji. Z twierdzenia Kleenego będzie on dany przez
\[
\lfp(\mathcal F)
=
\bigsqcup_{n=1}^{\infty}\mathcal F^n(\varnothing),
\]
o ile wykażemy, że $\mathcal F$ jest funkcjonałem ciągłym na
łańcuchowo zupełnym zbiorze
$(\Store\rightharpoonup\Store,\subseteq,\varnothing)$.

Aby zastosować twierdzenie Kleenego, musimy zatem wykazać, że przestrzeń
\[
\Store\rightharpoonup\Store
\]
jest zbiorem łańcuchowo zupełnym oraz że funkcjonał $\mathcal F$ jest
ciągły. Pierwszy z tych warunków wynika z
Twierdzenia~\ref{tw:funkcje_czesciowe_cpo}. Stosując to twierdzenie do
$A=B=\Store$, otrzymujemy, że
\[
\boldsymbol{\Store\rightharpoonup\Store}
=
(\Store\rightharpoonup\Store,\subseteq,\varnothing)
\]
jest zbiorem łańcuchowo zupełnym. Jego elementem najmniejszym jest
funkcja pusta $\varnothing$, a kres górny każdego łańcucha jest równy
sumie jego elementów:
\[
\bigsqcup_{n=1}^{\infty}f_n
=
\bigcup_{n=1}^{\infty}f_n.
\]

\subsection{Własności funkcjonału}

Pozostaje zatem zbadać własności funkcjonału $\mathcal{F}$, w
szczególności jego monotoniczność i ciągłość. 

Zaczniemy jednak od sprawdzenia jego poprawności. 

\begin{lemat}\label{lemat:F_dobrze_okreslony}
	Funkcjonał
	\[
	\mathcal{F}:
	(\Store\rightharpoonup\Store)
	\longrightarrow
	(\Store\rightharpoonup\Store)
	\]
	jest dobrze określony.
\end{lemat}

\begin{proof}
	Niech
	\[
	f:\Store\rightharpoonup\Store
	\]
	będzie dowolną funkcją częściową. Musimy wykazać, że
	\[
	\mathcal{F}(f)
	\]
	jest również funkcją częściową z $\Store$ do $\Store$.
	
	Z definicji
	\[
	\begin{aligned}
		\mathcal{F}(f)
		={}&
		\left\{
		(\sigma,\sigma)
		\;\middle|\;
		(\sigma,\kw{false})\in\Bb\sem{b}
		\right\}\\
		&{}\cup
		\left\{
		(\sigma,\sigma')
		\;\middle|\;
		\begin{array}{c}
			(\sigma,\kw{true})\in\Bb\sem{b},\\
			\exists\sigma''\in\Store\;
			\big(
			(\sigma,\sigma'')\in\Cc\sem{c}
			\land
			(\sigma'',\sigma')\in f
			\big)
		\end{array}
		\right\}.
	\end{aligned}
	\]
	
	Najpierw zauważmy, że
	\[
	\mathcal{F}(f)\subseteq\Store\times\Store,
	\]
	ponieważ wszystkie stany występujące w definicji $\mathcal{F}(f)$
	należą do $\Store$. Pozostaje zatem wykazać warunek jednoznaczności,
	tzn. że dla każdego $\sigma\in\Store$ istnieje co najwyżej jeden
	$\sigma'\in\Store$ taki, że
	\[
	(\sigma,\sigma')\in\mathcal{F}(f).
	\]
	
	Niech więc
	\[
	(\sigma,\sigma_1)\in\mathcal{F}(f)
	\quad\land\quad
	(\sigma,\sigma_2)\in\mathcal{F}(f).
	\]
	Pokażemy, że
	\[
	\sigma_1=\sigma_2.
	\]
	
	Wartość warunku $b$ dla ustalonego stanu $\sigma$ jest jednoznacznie
	określona. Rozpatrujemy zatem dwa przypadki.
	
	\emph{Przypadek 1.}
	Załóżmy, że
	\[
	(\sigma,\kw{false})\in\Bb\sem{b}.
	\]
	Wówczas obie pary należą do pierwszego składnika definicji
	$\mathcal{F}(f)$, ponieważ drugi składnik wymaga spełnienia warunku
	\[
	(\sigma,\kw{true})\in\Bb\sem{b}.
	\]
	Zatem
	\[
	\sigma_1=\sigma
	\quad\land\quad
	\sigma_2=\sigma,
	\]
	a więc
	\[
	\sigma_1=\sigma_2.
	\]
	
	\emph{Przypadek 2.}
	Załóżmy, że
	\[
	(\sigma,\kw{true})\in\Bb\sem{b}.
	\]
	Wówczas, z definicji $\mathcal{F}(f)$, istnieją
	\[
	\sigma_1'',\sigma_2''\in\Store
	\]
	takie, że
	\[
	(\sigma,\sigma_1'')\in\Cc\sem{c}
	\quad\land\quad
	(\sigma_1'',\sigma_1)\in f
	\]
	oraz
	\[
	(\sigma,\sigma_2'')\in\Cc\sem{c}
	\quad\land\quad
	(\sigma_2'',\sigma_2)\in f.
	\]
	
	Ponieważ $\Cc\sem{c}$ jest funkcją częściową, dla ustalonego stanu
	$\sigma$ istnieje co najwyżej jeden stan końcowy. Stąd
	\[
	(\sigma,\sigma_1'')\in\Cc\sem{c}
	\quad\land\quad
	(\sigma,\sigma_2'')\in\Cc\sem{c}
	\implies
	\sigma_1''=\sigma_2''.
	\]
	Oznaczmy wspólny stan przez $\sigma''$. Otrzymujemy wtedy
	\[
	(\sigma'',\sigma_1)\in f
	\quad\land\quad
	(\sigma'',\sigma_2)\in f.
	\]
	Ponieważ $f$ jest funkcją częściową, dla ustalonego stanu $\sigma''$
	istnieje co najwyżej jeden stan końcowy. Zatem
	\[
	\sigma_1=\sigma_2.
	\]
	
	W obu przypadkach otrzymujemy
	\[
	(\sigma,\sigma_1)\in\mathcal{F}(f)
	\quad\land\quad
	(\sigma,\sigma_2)\in\mathcal{F}(f)
	\implies
	\sigma_1=\sigma_2.
	\]
	Zatem $\mathcal{F}(f)$ spełnia warunek jednoznaczności funkcji
	częściowej, a więc
	\[
	\mathcal{F}(f):\Store\rightharpoonup\Store.
	\]
	Ponieważ $f:\Store\rightharpoonup\Store$ było dowolne, funkcjonał
	$\mathcal{F}$ jest dobrze określony.
\end{proof}

Aby móc zastosować twierdzenie Kleenego (Twierdzenie \ref{tw:kleene}), oprócz dobrego określenia
funkcjonału $\mathcal{F}$ musimy zbadać jego zachowanie względem
porządku na przestrzeni
\[
\Store\rightharpoonup\Store.
\]
Pierwszą z potrzebnych własności jest monotoniczność. Intuicyjnie,
jeżeli funkcja $g$ zawiera więcej informacji niż funkcja $f$, tzn.
\[
f\subseteq g,
\]
to zastosowanie funkcjonału $\mathcal{F}$ do $g$ nie powinno prowadzić
do utraty żadnej z par uzyskanych dla $f$. Pokażemy zatem, że
\[
f\subseteq g
\implies
\mathcal{F}(f)\subseteq\mathcal{F}(g).
\]

\begin{lemat}\label{lemat:F_monotoniczny}
	Funkcjonał $\mathcal{F}$ jest monotoniczny względem inkluzji wykresów,
	tzn. dla dowolnych funkcji częściowych
	\[
	f,g:\Store\rightharpoonup\Store
	\]
	zachodzi
	\[
	f\subseteq g
	\implies
	\mathcal{F}(f)\subseteq\mathcal{F}(g).
	\]
\end{lemat}

\begin{proof}
	Załóżmy, że
	\[
	f\subseteq g.
	\]
	Weźmy dowolną parę
	\[
	(\sigma,\sigma')\in\mathcal{F}(f).
	\]
	Pokażemy, że
	\[
	(\sigma,\sigma')\in\mathcal{F}(g).
	\]
	
	Z definicji funkcjonału $\mathcal{F}$ rozpatrujemy dwa przypadki
	zależnie od wartości warunku $b$ w stanie $\sigma$.
	
	\emph{Przypadek 1.}
	Załóżmy, że
	\[
	(\sigma,\kw{false})\in\Bb\sem{b}.
	\]
	Wówczas z definicji $\mathcal{F}(f)$ otrzymujemy
	\[
	\sigma'=\sigma.
	\]
	Zatem
	\[
	(\sigma,\sigma')
	=
	(\sigma,\sigma)
	\in\mathcal{F}(g),
	\]
	ponieważ pierwszy składnik definicji $\mathcal{F}(g)$ nie zależy od
	jej argumentu.
	
	\emph{Przypadek 2.}
	Załóżmy, że
	\[
	(\sigma,\kw{true})\in\Bb\sem{b}.
	\]
	Wówczas z definicji $\mathcal{F}(f)$ istnieje
	\[
	\sigma''\in\Store
	\]
	takie, że
	\[
	(\sigma,\sigma'')\in\Cc\sem{c}
	\]
	oraz
	\[
	(\sigma'',\sigma')\in f.
	\]
	Z założenia
	\[
	f\subseteq g
	\]
	wynika
	\[
	(\sigma'',\sigma')\in g.
	\]
	Wobec tego
	\[
	(\sigma,\sigma'')\in\Cc\sem{c}
	\quad\land\quad
	(\sigma'',\sigma')\in g,
	\]
	a zatem, na mocy definicji funkcjonału $\mathcal{F}$,
	\[
	(\sigma,\sigma')\in\mathcal{F}(g).
	\]
	
	W obu przypadkach otrzymujemy
	\[
	(\sigma,\sigma')\in\mathcal{F}(f)
	\implies
	(\sigma,\sigma')\in\mathcal{F}(g).
	\]
	Ponieważ $(\sigma,\sigma')$ było dowolnym elementem
	$\mathcal{F}(f)$, wynika stąd
	\[
	\mathcal{F}(f)\subseteq\mathcal{F}(g).
	\]
	Zatem $\mathcal{F}$ jest monotoniczny.
\end{proof}

Aby zastosować twierdzenie Kleenego do funkcjonału $\mathcal{F}$, musimy
wykazać, że jest on ciągły. Ciągłość oznacza, że $\mathcal{F}$ zachowuje
kresy górne łańcuchów. W rozważanej przestrzeni funkcji częściowych kres
górny łańcucha jest równy sumie jego elementów, dlatego dla dowolnego
łańcucha
\[
f_1\subseteq f_2\subseteq f_3\subseteq\ldots
\]
wystarczy wykazać równość
\[
\mathcal{F}
\left(
\bigcup_{n=1}^{\infty}f_n
\right)
=
\bigcup_{n=1}^{\infty}\mathcal{F}(f_n).
\]
Pokażemy ją, dowodząc obu inkluzji.

\begin{twierdzenie}\label{tw:F_ciagly}
	Funkcjonał
	\[
	\mathcal{F}:
	(\Store\rightharpoonup\Store)
	\longrightarrow
	(\Store\rightharpoonup\Store)
	\]
	jest ciągły względem porządku inkluzji wykresów.
\end{twierdzenie}

\begin{proof}
	Niech
	\[
	f_1\subseteq f_2\subseteq f_3\subseteq\ldots
	\]
	będzie dowolnym łańcuchem funkcji częściowych
	\[
	f_n:\Store\rightharpoonup\Store.
	\]
	Z Twierdzenia~\ref{tw:funkcje_czesciowe_cpo} mamy
	\[
	\bigsqcup_{n=1}^{\infty}f_n
	=
	\bigcup_{n=1}^{\infty}f_n.
	\]
	
	Z Lematu~\ref{lemat:F_monotoniczny} funkcjonał $\mathcal{F}$ jest
	monotoniczny, zatem
	\[
	\mathcal{F}(f_1)
	\subseteq
	\mathcal{F}(f_2)
	\subseteq
	\mathcal{F}(f_3)
	\subseteq\ldots.
	\]
	Wobec tego
	\[
	\bigsqcup_{n=1}^{\infty}\mathcal{F}(f_n)
	=
	\bigcup_{n=1}^{\infty}\mathcal{F}(f_n).
	\]
	
	Pokażemy więc, że
	\[
	\mathcal{F}
	\left(
	\bigcup_{n=1}^{\infty}f_n
	\right)
	=
	\bigcup_{n=1}^{\infty}\mathcal{F}(f_n).
	\]
	
	Najpierw wykażemy inkluzję
	\[
	\mathcal{F}
	\left(
	\bigcup_{n=1}^{\infty}f_n
	\right)
	\subseteq
	\bigcup_{n=1}^{\infty}\mathcal{F}(f_n).
	\]
	
	Niech
	\[
	(\sigma,\sigma')
	\in
	\mathcal{F}
	\left(
	\bigcup_{n=1}^{\infty}f_n
	\right).
	\]
	
	\emph{Przypadek 1.}
	Jeżeli
	\[
	(\sigma,\kw{false})\in\Bb\sem{b},
	\]
	to z definicji $\mathcal{F}$ mamy
	\[
	\sigma'=\sigma.
	\]
	Ten warunek nie zależy od argumentu funkcjonału, a więc
	\[
	(\sigma,\sigma')
	\in
	\mathcal{F}(f_n)
	\]
	dla każdego $n\in\mathbb{N}$. W szczególności
	\[
	(\sigma,\sigma')
	\in
	\bigcup_{n=1}^{\infty}\mathcal{F}(f_n).
	\]
	
	\emph{Przypadek 2.}
	Jeżeli
	\[
	(\sigma,\kw{true})\in\Bb\sem{b},
	\]
	to z definicji $\mathcal{F}$ istnieje
	$\sigma''\in\Store$ takie, że
	\[
	(\sigma,\sigma'')\in\Cc\sem{c}
	\]
	oraz
	\[
	(\sigma'',\sigma')
	\in
	\bigcup_{n=1}^{\infty}f_n.
	\]
	Z definicji sumy zbiorów istnieje $k\in\mathbb{N}$ takie, że
	\[
	(\sigma'',\sigma')\in f_k.
	\]
	Wówczas
	\[
	(\sigma,\sigma'')\in\Cc\sem{c}
	\quad\land\quad
	(\sigma'',\sigma')\in f_k,
	\]
	a więc
	\[
	(\sigma,\sigma')\in\mathcal{F}(f_k).
	\]
	Stąd
	\[
	(\sigma,\sigma')
	\in
	\bigcup_{n=1}^{\infty}\mathcal{F}(f_n).
	\]
	
	Zatem
	\[
	\mathcal{F}
	\left(
	\bigcup_{n=1}^{\infty}f_n
	\right)
	\subseteq
	\bigcup_{n=1}^{\infty}\mathcal{F}(f_n).
	\]
	
	Inkluzja przeciwna
	\[
	\bigcup_{n=1}^{\infty}\mathcal{F}(f_n)
	\subseteq
	\mathcal{F}
	\left(
	\bigcup_{n=1}^{\infty}f_n
	\right)
	\]
	wynika z monotoniczności $\mathcal{F}$. Ustalmy $k\in\mathbb N$.
	Z definicji sumy zbiorów
	\[
	f_k\subseteq\bigcup_{n=1}^{\infty}f_n,
	\]
	więc na mocy Lematu~\ref{lemat:F_monotoniczny}
	\[
	\mathcal{F}(f_k)
	\subseteq
	\mathcal{F}\left(\bigcup_{n=1}^{\infty}f_n\right).
	\]
	Ponieważ $k$ było dowolne, każdy ze zbiorów $\mathcal{F}(f_k)$ jest
	zawarty w $\mathcal{F}\left(\bigcup_{n}f_n\right)$, a więc również ich
	suma:
	\[
	\bigcup_{k=1}^{\infty}\mathcal{F}(f_k)
	\subseteq
	\mathcal{F}\left(\bigcup_{n=1}^{\infty}f_n\right).
	\]
	(Jest to dokładnie nierówność z Lematu~\ref{lem:monotoniczna_lancuchy}(b)
	zapisana w porządku inkluzji.)
	
	Z obu inkluzji wynika
	\[
	\mathcal{F}
	\left(
	\bigcup_{n=1}^{\infty}f_n
	\right)
	=
	\bigcup_{n=1}^{\infty}\mathcal{F}(f_n).
	\]
	Zatem
	\[
	\mathcal{F}
	\left(
	\bigsqcup_{n=1}^{\infty}f_n
	\right)
	=
	\bigsqcup_{n=1}^{\infty}\mathcal{F}(f_n).
	\]
	Wobec tego $\mathcal{F}$ jest ciągły.
\end{proof}

%==============================================================================
\subsection{Wyznaczenie denotacji pętli \kw{while}}
%==============================================================================

Możemy teraz zastosować twierdzenie Kleenego do równania
\[
W_{b,c}=\mathcal{F}\left(W_{b,c}\right).
\]
W tym celu sprawdziliśmy wszystkie wymagane założenia. Z
Twierdzenia~\ref{tw:funkcje_czesciowe_cpo} wynika, że przestrzeń
\[
\boldsymbol{\Store\rightharpoonup\Store}
=
(\Store\rightharpoonup\Store,\subseteq,\varnothing)
\]
jest zbiorem łańcuchowo zupełnym z elementem najmniejszym $\varnothing$.
Z Lematu~\ref{lemat:F_dobrze_okreslony} wynika, że funkcjonał
\[
\mathcal{F}:
(\Store\rightharpoonup\Store)
\longrightarrow
(\Store\rightharpoonup\Store)
\]
jest dobrze określony, natomiast na mocy
Lematu~\ref{lemat:F_monotoniczny} oraz
Twierdzenia~\ref{tw:F_ciagly} jest on ciągły.

Wszystkie założenia twierdzenia Kleenego są zatem spełnione. Wynika stąd,
że $\mathcal{F}$ posiada najmniejszy punkt stały, który jest równy kresowi
górnemu ciągu kolejnych przybliżeń od elementu najmniejszego:
\[
\lfp(\mathcal{F})
=
\bigsqcup_{i=1}^{\infty}\mathcal{F}^i(\varnothing).
\]

Zgodnie z wcześniejszymi rozważaniami \emph{definiujemy} denotację pętli
jako ten najmniejszy punkt stały:
\[
	W_{b,c}
	:=
	\lfp(\mathcal{F})
	=
	\bigsqcup_{i=1}^{\infty}\mathcal{F}^i(\varnothing).
\]

Na mocy Twierdzenia~\ref{tw:funkcje_czesciowe_cpo} kres górny łańcucha
funkcji częściowych jest równy sumie jego elementów, dlatego
\[
	W_{b,c}
	=
	\bigcup_{i=1}^{\infty}\mathcal{F}^i(\varnothing).
\]

Możemy teraz jawnie zobaczyć, jak powstają kolejne przybliżenia.

Pierwsze przybliżenie otrzymujemy przez zastosowanie funkcjonału
$\mathcal{F}$ do funkcji pustej:
\[
\begin{aligned}
	\mathcal{F}^1(\varnothing)
	=
	\mathcal{F}(\varnothing)
	=
	\left\{
	(\sigma,\sigma)
	\;\middle|\;
	(\sigma,\kw{false})\in\Bb\sem{b}
	\right\}.
\end{aligned}
\]
Opisuje ono wszystkie wykonania pętli, które kończą się bez wykonania
ciała $c$.

Drugie przybliżenie ma postać
\[
\begin{aligned}
	\mathcal{F}^2(\varnothing)
	={}&
	\left\{
	(\sigma,\sigma)
	\;\middle|\;
	(\sigma,\kw{false})\in\Bb\sem{b}
	\right\}
	\\
	&{}\cup
	\left\{
	(\sigma,\sigma')
	\;\middle|\;
	\begin{array}{c}
		(\sigma,\kw{true})\in\Bb\sem{b},\\
		(\sigma,\sigma')\in\Cc\sem{c},\\
		(\sigma',\kw{false})\in\Bb\sem{b}
	\end{array}
	\right\}.
\end{aligned}
\]
Zawiera ono wszystkie wykonania, które kończą się bez wykonania ciała
albo po wykonaniu ciała $c$ dokładnie jeden raz.

Trzecie przybliżenie jest równe
\[
\begin{aligned}
	\mathcal{F}^3(\varnothing)
	={}&
	\left\{
	(\sigma,\sigma)
	\;\middle|\;
	(\sigma,\kw{false})\in\Bb\sem{b}
	\right\}
	\\
	&{}\cup
	\left\{
	(\sigma,\sigma')
	\;\middle|\;
	\begin{array}{c}
		(\sigma,\kw{true})\in\Bb\sem{b},\\
		(\sigma,\sigma')\in\Cc\sem{c},\\
		(\sigma',\kw{false})\in\Bb\sem{b}
	\end{array}
	\right\}
	\\
	&{}\cup
	\left\{
	(\sigma,\sigma'')
	\;\middle|\;
	\begin{array}{c}
		(\sigma,\kw{true})\in\Bb\sem{b},\\
		\exists\sigma'\in\Store\;
		\big(
		(\sigma,\sigma')\in\Cc\sem{c}
		\land
		(\sigma',\kw{true})\in\Bb\sem{b}
		\\
		\qquad\qquad\qquad
		\land
		(\sigma',\sigma'')\in\Cc\sem{c}
		\land
		(\sigma'',\kw{false})\in\Bb\sem{b}
		\big)
	\end{array}
	\right\}.
\end{aligned}
\]

Widzimy zatem, że kolejne przybliżenia tworzą łańcuch
\[
\mathcal{F}^1(\varnothing)
\subseteq
\mathcal{F}^2(\varnothing)
\subseteq
\mathcal{F}^3(\varnothing)
\subseteq\ldots.
\]
Przybliżenie $\mathcal{F}^i(\varnothing)$ zawiera wszystkie wykonania
pętli, które mogą zakończyć się po co najwyżej $i-1$ wykonaniach jej
ciała. Każde kolejne zastosowanie funkcjonału może zatem dodać wykonania
wymagające jeszcze jednego przejścia przez ciało pętli.

W konsekwencji semantyka pętli jest dana przez sumę wszystkich tych
przybliżeń:
\[
	W_{b,c}
	=
	\bigsqcup_{i=1}^{\infty}\mathcal{F}^i(\varnothing)
	=
	\bigcup_{i=1}^{\infty}\mathcal{F}^i(\varnothing).
\]

Ostatecznie otrzymujemy jawną postać denotacji:
\[
	\Cc\sem{\kw{while}\ b\ \kw{do}\ c}
	=
	\bigsqcup_{i=1}^{\infty}\mathcal{F}^i(\varnothing)
	=
	\bigcup_{i=1}^{\infty}\mathcal{F}^i(\varnothing).
\]



\begin{przyklad}
	Rozważmy przykład
	\[
	\kw{while}\ \texttt{foo}<5\ \kw{do}\ 
	\texttt{foo}:=\texttt{foo}+1.
	\]
	Mamy zatem
	\[
	b=\texttt{foo}<5,
	\qquad
	c=\texttt{foo}:=\texttt{foo}+1.
	\]
	
	Dla pierwszego przybliżenia otrzymujemy
	\[
	\begin{aligned}
		\mathcal{F}(\varnothing)
		={}&
		\left\{
		(\sigma,\sigma)
		\;\middle|\;
		(\sigma,\kw{false})\in\Bb\sem{b}
		\right\}\\
		&{}\cup
		\left\{
		(\sigma,\sigma')
		\;\middle|\;
		\begin{array}{c}
			(\sigma,\kw{true})\in\Bb\sem{b},\\
			\exists\sigma''\;
			\big(
			(\sigma,\sigma'')\in\Cc\sem{c}
			\land
			(\sigma'',\sigma')\in\varnothing
			\big)
		\end{array}
		\right\}\\
		={}&
		\left\{
		(\sigma,\sigma)
		\;\middle|\;
		\sigma(\texttt{foo})\geq5
		\right\}.
	\end{aligned}
	\]
	
	Drugie przybliżenie ma postać
	\[
	\begin{aligned}
		\mathcal{F}^2(\varnothing)
		={}&
		\left\{
		(\sigma,\sigma)
		\;\middle|\;
		(\sigma,\kw{false})\in\Bb\sem{b}
		\right\}\\
		&{}\cup
		\left\{
		(\sigma,\sigma')
		\;\middle|\;
		\begin{array}{c}
			(\sigma,\kw{true})\in\Bb\sem{b},\\
			\exists\sigma''\;
			\big(
			(\sigma,\sigma'')\in\Cc\sem{c}
			\land
			(\sigma'',\sigma')\in\mathcal{F}(\varnothing)
			\big)
		\end{array}
		\right\}\\
		={}&
		\left\{
		(\sigma,\sigma)
		\;\middle|\;
		\sigma(\texttt{foo})\geq5
		\right\}\\
		&{}\cup
		\left\{
		\left(
		\sigma,
		\sigma[\texttt{foo}\mapsto\sigma(\texttt{foo})+1]
		\right)
		\;\middle|\;
		\begin{array}{c}
			\sigma(\texttt{foo})<5,\\
			\sigma(\texttt{foo})+1\geq5
		\end{array}
		\right\}.
	\end{aligned}
	\]
	
	Ponieważ
	\[
	\sigma(\texttt{foo})<5
	\quad\land\quad
	\sigma(\texttt{foo})+1\geq5
	\]
	implikuje
	\[
	\sigma(\texttt{foo})+1=5,
	\]
	otrzymujemy
	\[
	\mathcal{F}^2(\varnothing)
	=
	\left\{
	(\sigma,\sigma)
	\;\middle|\;
	\sigma(\texttt{foo})\geq5
	\right\}
	\cup
	\left\{
	\left(
	\sigma,
	\sigma[\texttt{foo}\mapsto\sigma(\texttt{foo})+1]
	\right)
	\;\middle|\;
	\sigma(\texttt{foo})+1=5
	\right\}.
	\]
	
	Trzecie przybliżenie otrzymujemy analogicznie:
	\[
	\begin{aligned}
		\mathcal{F}^3(\varnothing)
		={}&
		\left\{
		(\sigma,\sigma)
		\;\middle|\;
		\sigma(\texttt{foo})\geq5
		\right\}\\
		&{}\cup
		\left\{
		\left(
		\sigma,
		\sigma[\texttt{foo}\mapsto\sigma(\texttt{foo})+1]
		\right)
		\;\middle|\;
		\sigma(\texttt{foo})+1=5
		\right\}\\
		&{}\cup
		\left\{
		\left(
		\sigma,
		\sigma[\texttt{foo}\mapsto\sigma(\texttt{foo})+2]
		\right)
		\;\middle|\;
		\sigma(\texttt{foo})+2=5
		\right\}.
	\end{aligned}
	\]
	
	Czwarte przybliżenie ma postać
	\[
	\begin{aligned}
		\mathcal{F}^4(\varnothing)
		={}&
		\left\{
		(\sigma,\sigma)
		\;\middle|\;
		\sigma(\texttt{foo})\geq5
		\right\}\\
		&{}\cup
		\left\{
		\left(
		\sigma,
		\sigma[\texttt{foo}\mapsto\sigma(\texttt{foo})+1]
		\right)
		\;\middle|\;
		\sigma(\texttt{foo})+1=5
		\right\}\\
		&{}\cup
		\left\{
		\left(
		\sigma,
		\sigma[\texttt{foo}\mapsto\sigma(\texttt{foo})+2]
		\right)
		\;\middle|\;
		\sigma(\texttt{foo})+2=5
		\right\}\\
		&{}\cup
		\left\{
		\left(
		\sigma,
		\sigma[\texttt{foo}\mapsto\sigma(\texttt{foo})+3]
		\right)
		\;\middle|\;
		\sigma(\texttt{foo})+3=5
		\right\}.
	\end{aligned}
	\]
	
	Biorąc sumę wszystkich przybliżeń, otrzymujemy
	\[
	\begin{aligned}
		W_{b,c}
		={}&
		\bigcup_{i=1}^{\infty}\mathcal{F}^i(\varnothing)\\
		={}&
		\left\{
		(\sigma,\sigma)
		\;\middle|\;
		\sigma(\texttt{foo})\geq5
		\right\}\\
		&{}\cup
		\left\{
		\left(
		\sigma,
		\sigma[\texttt{foo}\mapsto\sigma(\texttt{foo})+n]
		\right)
		\;\middle|\;
		\sigma(\texttt{foo})+n=5
		\land
		n\geq1
		\right\}.
	\end{aligned}
	\]
	
	Ostatnią równość możemy zapisać w postaci zamkniętej:
	\[
		W_{b,c}
		=
		\left\{
		(\sigma,\sigma)
		\;\middle|\;
		\sigma(\texttt{foo})\geq5
		\right\}
		\cup
		\left\{
		\left(
		\sigma,
		\sigma[\texttt{foo}\mapsto\sigma(\texttt{foo})+n]
		\right)
		\;\middle|\;
		n=5-\sigma(\texttt{foo})
		\land
		n\geq1
		\right\}.
	\]
	
	W szczególności, jeżeli początkowo
	\[
	\sigma(\texttt{foo})=1,
	\]
	to
	\[
	n=5-1=4,
	\]
	a więc
	\[
	\Cc\sem{\kw{while}\ \texttt{foo}<5\ \kw{do}\
		\texttt{foo}:=\texttt{foo}+1}(\sigma)
	=
	\sigma[\texttt{foo}\mapsto5].
	\]
\end{przyklad}

%==============================================================================
\subsection{Rekurencyjna definicja silni}
%==============================================================================

Konstrukcję najmniejszego punktu stałego można zastosować również do
rekurencyjnych definicji funkcji. 

Rozważmy klasyczną definicję silni:
\[
\operatorname{silnia}(0)=1,
\qquad
\operatorname{silnia}(n+1)
=
(n+1)\cdot\operatorname{silnia}(n).
\]

Dziedziną funkcji silnia jest zbiór
\[
\mathbb{N}_0=\{0,1,2,\ldots\}.
\]
Oznaczmy wartość silni dla argumentu $n$ przez
\[
S_n:=n!.
\]
Wówczas
\[
S_0=1,
\qquad
S_{n+1}=(n+1) \cdot S_n.
\]

Niech
\[
D:=\mathbb{N}_{0,\bt}
=
\mathbb{N}_0\cup\{\bt\}
\]
będzie płaskim zbiorem łańcuchowo zupełnym otrzymanym przez dodanie
elementu najmniejszego $\bt$.

Rozważamy przestrzeń wszystkich funkcji całkowitych
\[
D\to D
\]
uporządkowaną porządkiem pionowym
\[
f\sqsubseteq g
\iff
\forall x\in D\quad
f(x)\sqsubseteq_D g(x).
\]
Na mocy Twierdzenia~\ref{tw:funkcje_calkowite_cpo} przestrzeń
\[
\boldsymbol{D\to D}
=
(D\to D,\sqsubseteq,\varphi)
\]
jest zbiorem łańcuchowo zupełnym, gdzie elementem najmniejszym jest
funkcja stała
\[
\varphi(x)=\bt \quad (x \in D).
\]

Niech
\[
S:D\to D
\]
oznacza poszukiwaną denotację funkcji $\operatorname{silnia}$. Z jej
definicji rekurencyjnej oczekujemy
\[
S(0)=1
\]
oraz
\[
S(n+1)=(n+1)\cdot S(n).
\]
Możemy zatem zapisać tę definicję jako równanie stałopunktowe
\[
S=\mathcal{F}(S),
\]
gdzie
\[
\mathcal{F}:
(D\to D)
\longrightarrow
(D\to D)
\]
jest funkcjonałem określonym wzorem
\[
\mathcal{F}(f)(x)
=
\begin{cases}
	\bt,&x=\bt,\\
	1,&x=0,\\
	(n+1)\cdot f(n),&x=n+1,\quad n\in\mathbb{N}_0.
\end{cases}
\]
W powyższym wzorze przyjmujemy dodatkowo
\[
(n+1)\cdot\bt=\bt
\qquad(n\in\mathbb{N}_0).
\]
Dzięki temu dla każdej funkcji
\[
f:D\to D
\]
funkcja $\mathcal{F}(f)$ jest również funkcją całkowitą
\[
\mathcal{F}(f):D\to D.
\]

Aby zastosować twierdzenie Kleenego, musimy wykazać, że funkcjonał
$\mathcal{F}$ jest ciągły względem porządku pionowego.

\begin{twierdzenie}\label{tw:silnia_funkcjonal_ciagly}
	Funkcjonał
	\[
	\mathcal{F}:(D\to D)\longrightarrow(D\to D)
	\]
	określony powyższym wzorem jest ciągły.
\end{twierdzenie}

\begin{proof}
	\emph{Krok 1: pomocnicza funkcja $h_n$.} Ustalmy $n\in\mathbb N_0$
	i rozważmy funkcję $h_n:D\to D$ ,,mnożenia przez $n+1$'':
	\[
	h_n(d)=
	\begin{cases}
		\bt,&d=\bt,\\
		(n+1)\cdot d,&d\in\mathbb N_0.
	\end{cases}
	\]
	Zgodnie z przyjętą konwencją $(n+1)\cdot\bt=\bt$ możemy pisać krótko
	$h_n(d)=(n+1)\cdot d$ dla wszystkich $d\in D$. W szczególności
	\[
	h_n(\bt)=\bt,
	\]
	czyli $h_n$ zachowuje element najmniejszy. Ponieważ $D$ jest płaskim
	zbiorem łańcuchowo zupełnym, na mocy
	Wniosku~\ref{wn:plaskie_monotoniczne_ciagle} funkcja $h_n$ jest
	monotoniczna i ciągła.
	
	Za pomocą funkcji $h_n$ wzór na $\mathcal{F}$ możemy zapisać jako
	\[
	\mathcal{F}(f)(\bt)=\bt,
	\qquad
	\mathcal{F}(f)(0)=1,
	\qquad
	\mathcal{F}(f)(n+1)=(n+1)\cdot f(n)=h_n\big(f(n)\big)
	\quad(n\in\mathbb N_0).
	\]
	
	\emph{Krok 2: monotoniczność $\mathcal{F}$.} Niech
	$g,g':D\to D$ oraz $g\sqsubseteq g'$, tzn.\
	$g(x)\sqsubseteq_D g'(x)$ dla każdego $x\in D$. Porównajmy wartości
	$\mathcal{F}(g)$ i $\mathcal{F}(g')$ w dowolnym punkcie $x\in D$.
	\begin{itemize}
		\item Dla $x=\bt$: $\mathcal{F}(g)(\bt)=\bt=\mathcal{F}(g')(\bt)$.
		\item Dla $x=0$: $\mathcal{F}(g)(0)=1=\mathcal{F}(g')(0)$.
		\item Dla $x=n+1$: z założenia $g(n)\sqsubseteq_D g'(n)$, a z
		monotoniczności $h_n$
		\[
		\mathcal{F}(g)(n+1)
		=
		h_n\big(g(n)\big)
		\sqsubseteq_D
		h_n\big(g'(n)\big)
		=
		\mathcal{F}(g')(n+1).
		\]
	\end{itemize}
	W każdym przypadku $\mathcal{F}(g)(x)\sqsubseteq_D\mathcal{F}(g')(x)$,
	więc $\mathcal{F}(g)\sqsubseteq\mathcal{F}(g')$. Funkcjonał
	$\mathcal{F}$ jest zatem monotoniczny.
	
	\emph{Krok 3: obraz łańcucha jest łańcuchem.} Niech
	\[
	f_1\sqsubseteq f_2\sqsubseteq f_3\sqsubseteq\ldots
	\]
	będzie łańcuchem w $\boldsymbol{D\to D}$. Z kroku 2 i
	Lematu~\ref{lem:monotoniczna_lancuchy}(a) ciąg
	\[
	\mathcal{F}(f_1)\sqsubseteq\mathcal{F}(f_2)\sqsubseteq\mathcal{F}(f_3)\sqsubseteq\ldots
	\]
	jest łańcuchem, a więc ma kres górny. Jest to warunek~(1) definicji
	ciągłości.
	
	\emph{Krok 4: zachowanie kresów górnych.} Oznaczmy
	\[
	f=\bigsqcup_{i=1}^{\infty}f_i.
	\]
	Na mocy Twierdzenia~\ref{tw:funkcje_calkowite_cpo} kresy górne w
	$\boldsymbol{D\to D}$ wyznaczamy punktowo, więc
	\[
	f(x)=\bigsqcup_{i=1}^{\infty}f_i(x)
	\qquad\text{oraz}\qquad
	\Big(\bigsqcup_{i=1}^{\infty}\mathcal{F}(f_i)\Big)(x)
	=
	\bigsqcup_{i=1}^{\infty}\mathcal{F}(f_i)(x)
	\qquad(x\in D).
	\]
	Musimy wykazać, że
	$\mathcal{F}(f)(x)=\bigsqcup_{i}\mathcal{F}(f_i)(x)$ dla każdego
	$x\in D$.
	\begin{itemize}
		\item Dla $x=\bt$: $\mathcal{F}(f)(\bt)=\bt$ oraz
		$\mathcal{F}(f_i)(\bt)=\bt$ dla każdego $i$. Kresem górnym łańcucha
		stałego $\bt,\bt,\ldots$ jest $\bt$, więc obie strony są równe $\bt$.
		\item Dla $x=0$: analogicznie obie strony są równe $1$.
		\item Dla $x=n+1$: ciąg $\big(f_i(n)\big)_{i\in\mathbb N}$ jest
		łańcuchem w $D$ (bo $f_i\sqsubseteq f_{i+1}$ w porządku pionowym),
		a jego kresem jest $f(n)$. Korzystając z ciągłości $h_n$, otrzymujemy
		\[
		\mathcal{F}(f)(n+1)
		=
		h_n\big(f(n)\big)
		=
		h_n\Big(\bigsqcup_{i=1}^{\infty}f_i(n)\Big)
		=
		\bigsqcup_{i=1}^{\infty}h_n\big(f_i(n)\big)
		=
		\bigsqcup_{i=1}^{\infty}\mathcal{F}(f_i)(n+1).
		\]
	\end{itemize}
	Wartości funkcji $\mathcal{F}(f)$ i $\bigsqcup_{i}\mathcal{F}(f_i)$ są
	równe w każdym punkcie, więc
	\[
	\mathcal{F}\Big(\bigsqcup_{i=1}^{\infty}f_i\Big)
	=
	\bigsqcup_{i=1}^{\infty}\mathcal{F}(f_i).
	\]
	Jest to warunek~(2) definicji ciągłości. Zatem funkcjonał
	$\mathcal{F}$ jest ciągły.
\end{proof}

Możemy teraz zastosować twierdzenie Kleenego. Otrzymujemy
\[
\lfp(\mathcal{F})
=
\bigsqcup_{i=1}^{\infty}\mathcal{F}^i(\varphi).
\]

Oznaczmy kolejne przybliżenia przez
\[
f_i=\mathcal{F}^i(\varphi).
\]
Otrzymujemy wówczas następującą tabelę:
\[
\begin{array}{c|ccccccc}
	x
	&\bt&0&1&2&3&4&\cdots\\
	\hline
	\varphi
	&\bt&\bt&\bt&\bt&\bt&\bt&\cdots\\
	f_1
	&\bt&1&\bt&\bt&\bt&\bt&\cdots\\
	f_2
	&\bt&1&1&\bt&\bt&\bt&\cdots\\
	f_3
	&\bt&1&1&2&\bt&\bt&\cdots\\
	f_4
	&\bt&1&1&2&6&\bt&\cdots\\
	f_5
	&\bt&1&1&2&6&24&\cdots\\
	\vdots
	&\vdots&\vdots&\vdots&\vdots&\vdots&\vdots&\ddots
\end{array}
\]

Wprowadzona wcześniej notacja $S_n=n!$ pozwala zapisać tę tabelę bardziej
symbolicznie:
\[
\begin{array}{c|ccccccc}
	x
	&\bt&0&1&2&3&4&\cdots\\
	\hline
	\varphi
	&\bt&\bt&\bt&\bt&\bt&\bt&\cdots\\
	f_1
	&\bt&S_0&\bt&\bt&\bt&\bt&\cdots\\
	f_2
	&\bt&S_0&S_1&\bt&\bt&\bt&\cdots\\
	f_3
	&\bt&S_0&S_1&S_2&\bt&\bt&\cdots\\
	f_4
	&\bt&S_0&S_1&S_2&S_3&\bt&\cdots\\
	f_5
	&\bt&S_0&S_1&S_2&S_3&S_4&\cdots\\
	\vdots
	&\vdots&\vdots&\vdots&\vdots&\vdots&\vdots&\ddots
\end{array}
\]

Pierwsze przybliżenie ma postać
\[
f_1=\mathcal{F}(\varphi).
\]
Dla argumentu $0$ otrzymujemy
\[
f_1(0)=1=S_0,
\]
natomiast dla $n\geq1$
\[
f_1(n)
=
n\cdot\varphi(n-1)
=
\bt.
\]

Dla drugiego przybliżenia mamy
\[
f_2(0)=1=S_0,
\]
oraz
\[
f_2(1)
=
1\cdot f_1(0)
=
1
=
S_1.
\]
Natomiast
\[
f_2(2)
=
2\cdot f_1(1)
=
\bt.
\]

Analogicznie
\[
f_3(2)
=
2\cdot f_2(1)
=
2
=
S_2,
\]
a następnie
\[
f_4(3)
=
3\cdot f_3(2)
=
6
=
S_3
\]
oraz
\[
f_5(4)
=
4\cdot f_4(3)
=
24
=
S_4.
\]

Ponieważ $\mathcal{F}$ jest ciągły, na mocy
Twierdzenia~\ref{tw:ciagla_monotoniczna} jest również monotoniczny.
Ponieważ $\varphi$ jest elementem najmniejszym, mamy
\[
\varphi\sqsubseteq\mathcal{F}(\varphi)=f_1.
\]
Stosując monotoniczność $\mathcal{F}$ indukcyjnie, otrzymujemy łańcuch
\[
\varphi
\sqsubseteq
f_1
\sqsubseteq
f_2
\sqsubseteq
f_3
\sqsubseteq
\ldots.
\]

W granicy otrzymujemy
\[
\lfp(\mathcal{F})
=
\bigsqcup_{i=1}^{\infty}f_i.
\]
Ponieważ kres górny w przestrzeni funkcji całkowitych jest wyznaczany
punktowo, dla każdego $n\in\mathbb{N}_0$ mamy
\[
\lfp(\mathcal{F})(n)=S_n=n!,
\]
a ponadto
\[
\lfp(\mathcal{F})(\bt)=\bt.
\]

Zatem najmniejszy punkt stały funkcjonału $\mathcal{F}$ jest funkcją
\[
S:D\to D
\]
określoną przez
\[
S(x)
=
\begin{cases}
	\bt,&x=\bt,\\
	x!,&x\in\mathbb{N}_0.
\end{cases}
\]

W szczególności
\[
S(\bt)=\bt.
\]
Ponieważ $D$ jest płaskim zbiorem łańcuchowo zupełnym, z
Wniosku~\ref{wn:plaskie_monotoniczne_ciagle} wynika, że funkcja $S$
jest ciągła.

Otrzymaliśmy zatem funkcję ciągłą, która dla każdego
$n\in\mathbb{N}_0$ przyjmuje wartość
\[
S(n)=n!,
\]
a więc stanowi rozszerzenie zwykłej funkcji silnia na zbiorze łańcuchowo zupełnym.

%==============================================================================
\section{Podejście alternatywne}
%==============================================================================

W poprzednich częściach punkty stałe rozważaliśmy z perspektywy porządku
informacyjnego. Zbiór łańcuchowo zupełny dostarczał przestrzeni, w której
można było porównywać kolejne przybliżenia, natomiast ciągłość funkcji
gwarantowała istnienie najmniejszego punktu stałego, który otrzymywaliśmy
jako kres górny ciągu iteracji od elementu najmniejszego.

To samo zagadnienie można rozpatrywać z innej perspektywy, zastępując
strukturę porządkową strukturą opartą na pojęciu odległości. Takie podejście
nie prowadzi jedynie do innego sposobu zapisywania tych samych konstrukcji.
Pozwala ono również wykorzystać narzędzia analizy klasycznej, w szczególności
pojęcie ciągu Cauchy'ego oraz twierdzenie o odwzorowaniu zwężającym.

W pracy \emph{Partial Metric Topology}~\cite{matthews1994}
S.~G.~Matthews wprowadza w tym celu pojęcie \emph{metryki częściowej}
(ang.\ \emph{partial metric}). Jej istotną cechą jest to, że obiekt może mieć
niezerową odległość od samego siebie. Wartość $p(x,x)$ może dzięki temu
opisywać stopień określenia obiektu: obiekty o wartości $p(x,x)=0$ są
całkowicie określone, natomiast dodatnia wartość wskazuje na obecność
nieokreśloności.

Podejście to jest szczególnie interesujące w kontekście teorii przedstawionej
wcześniej, ponieważ metryka częściowa nie jest niezależna od pojęcia
porządku. Zostaje z niej wyprowadzony porządek informacyjny, a własności
związane z granicami łańcuchów znajdują swój odpowiednik w pojęciach
zbieżności i zupełności metrycznej.

W dalszej części przedstawimy podstawowe pojęcia związane z metrykami
częściowymi, pokażemy sposób odzyskiwania z nich porządku informacyjnego,
a następnie przedstawimy twierdzenie Matthewsa o punkcie stałym kontrakcji.
Dzięki temu otrzymamy drugie, metryczne ujęcie problemu punktów stałych,
które można porównać z wcześniejszym podejściem porządkowym.

%==============================================================================
\subsection{Podstawy metryczne}
%==============================================================================

Zanim wprowadzimy metrykę częściową, przypomnijmy klasyczne pojęcie
metryki. Będzie ono stanowiło punkt odniesienia dla konstrukcji
Matthewsa.

\begin{definicja}\label{def:metryka}
	Metryką na zbiorze $U$ nazywamy funkcję
	\[
	d:U\times U\to [0, \infty)
	\]
	spełniającą dla dowolnych $x,y,z\in U$ następujące warunki:
	\begin{itemize}[leftmargin=1.5cm]
		\item[(M1)] $d(x,y)=0 \iff x=y$;
		\item[(M2)] $d(x,y)=d(y,x)$;
		\item[(M3)] $d(x,z)\leq d(x,y)+d(y,z)$.
	\end{itemize}
\end{definicja}

Z warunku (M1) w szczególności wynika, że
\[
d(x,x)=0
\]
dla każdego $x\in U$. W zwykłej metryce własna odległość obiektu nie
niesie więc żadnej dodatkowej informacji. W podejściu Matthewsa własność
ta zostanie osłabiona. Odległość obiektu od samego siebie może być dodatnia
i może służyć do opisu stopnia jego określenia.

Prowadzi to do pojęcia metryki częściowej, wprowadzonego przez
S.~G.~Matthewsa~\cite{matthews1994}.

%==============================================================================
\subsection{Metryka częściowa}
%==============================================================================

\begin{definicja}\label{def:metryka_czesciowa}
	\emph{Metryką częściową} na zbiorze $U$ nazywamy funkcję
	\[
	p:U\times U\to [0, \infty)
	\]
	spełniającą dla dowolnych $x,y,z\in U$ następujące warunki:
	\begin{itemize}[leftmargin=1.5cm]
		\item[(P1)] $x=y \iff p(x,x)=p(x,y)=p(y,y)$;
		\item[(P2)] $p(x,x)\leq p(x,y)$;
		\item[(P3)] $p(x,y)=p(y,x)$;
		\item[(P4)] $p(x,z)\leq p(x,y)+p(y,z)-p(y,y)$.
	\end{itemize}
\end{definicja}

Warunek (P1) zastępuje pierwszy aksjomat zwykłej metryki. W szczególności
nie wymaga on, aby
\[
p(x,x)=0.
\]
Możliwe jest więc, że
\[
p(x,x)>0,
\]
a wartość ta może być interpretowana jako miara stopnia nieokreśloności
obiektu $x$. Obiekt o własnej odległości równej zero jest natomiast
całkowicie określony.

Warunek (P2) oznacza, że własna odległość $p(x,x)$ jest najmniejszą
odległością, jaką obiekt $x$ może mieć od dowolnego elementu:
\[
p(x,x)\leq p(x,y).
\]
Warunek (P3) zachowuje symetrię zwykłej metryki. W warunku (P4)
nierówność trójkąta otrzymuje natomiast postać
\[
p(x,z)\leq p(x,y)+p(y,z)-p(y,y),
\]
w której uwzględniona zostaje własna odległość punktu pośredniego.

\begin{uwaga}
	Jeżeli dla każdego $x\in U$ zachodzi
	\[
	p(x,x)=0,
	\]
	to warunki (P1)--(P4) redukują się do warunków metryki z
	Definicji~\ref{def:metryka}. Zwykłą metrykę można więc traktować jako
	szczególny przypadek metryki częściowej.
\end{uwaga}

\begin{przyklad}\label{prz:plaska_pmetryka}
	Niech $D$ będzie dowolnym zbiorem oraz
	\[
	D_\bt=D\cup\{\bt\},
	\]
	gdzie $\bt\notin D$. Na zbiorze $D_\bt$ definiujemy funkcję
	\[
	p:D_\bt\times D_\bt\to\{0,1\}
	\]
	wzorem
	\[
	p(x,y)=
	\begin{cases}
		0,&x=y\in D,\\
		1,&\text{w przeciwnym przypadku}.
	\end{cases}
	\]
	Pokażemy, że $p$ jest metryką częściową.
	
	\emph{(P1).} Załóżmy najpierw, że
	\[
	x=y.
	\]
	Jeżeli $x=y\in D$, to
	\[
	p(x,x)=p(x,y)=p(y,y)=0.
	\]
	Jeżeli natomiast
	\[
	x=y=\bt,
	\]
	to
	\[
	p(x,x)=p(x,y)=p(y,y)=1.
	\]
	Zatem
	\[
	x=y
	\implies
	p(x,x)=p(x,y)=p(y,y).
	\]
	
	Załóżmy teraz, że
	\[
	p(x,x)=p(x,y)=p(y,y).
	\]
	Jeżeli $x\neq y$, to z definicji $p$ mamy
	\[
	p(x,y)=1.
	\]
	Jeżeli jednocześnie $x,y\in D$, to
	\[
	p(x,x)=p(y,y)=0,
	\]
	co daje sprzeczność. Jeżeli natomiast jeden z elementów jest równy
	$\bt$, a drugi należy do $D$, to ich własne odległości wynoszą odpowiednio
	$1$ i $0$, co również daje sprzeczność. Jedynym pozostałym przypadkiem
	jest $x=y=\bt$, sprzeczny z założeniem $x\neq y$. Zatem
	\[
	p(x,x)=p(x,y)=p(y,y)
	\implies
	x=y.
	\]
	W konsekwencji zachodzi
	\[
	x=y
	\iff
	p(x,x)=p(x,y)=p(y,y).
	\]
	
	\emph{(P2).} Jeżeli $x\in D$, to
	\[
	p(x,x)=0\leq p(x,y).
	\]
	Jeżeli $x=\bt$, to
	\[
	p(x,x)=p(\bt,\bt)=1
	\]
	oraz
	\[
	p(\bt,y)=1
	\]
	dla każdego $y\in D_\bt$. W obu przypadkach
	\[
	p(x,x)\leq p(x,y).
	\]
	
	\emph{(P3).} Z definicji funkcji $p$ wynika bezpośrednio
	\[
	p(x,y)=p(y,x).
	\]
	
	\emph{(P4).} Rozważmy dwa przypadki.
	
	Jeżeli
	\[
	y=\bt,
	\]
	to
	\[
	p(x,y)=p(y,z)=p(\bt,\bt)=1,
	\]
	a więc
	\[
	p(x,y)+p(y,z)-p(y,y)
	=
	1+1-1
	=
	1.
	\]
	Ponieważ $p(x,z)\in\{0,1\}$, otrzymujemy
	\[
	p(x,z)\leq
	p(x,y)+p(y,z)-p(y,y).
	\]
	
	Jeżeli natomiast
	\[
	y\in D,
	\]
	to
	\[
	p(y,y)=0.
	\]
	Jeżeli $x=z$, to
	\[
	p(x,z)=0,
	\]
	a nierówność jest oczywista. Jeżeli $x\neq z$, to
	\[
	p(x,z)=1.
	\]
	Nie mogą jednocześnie zachodzić
	\[
	p(x,y)=0
	\qquad\text{oraz}\qquad
	p(y,z)=0,
	\]
	ponieważ oznaczałoby to
	\[
	x=y
	\qquad\text{oraz}\qquad
	y=z,
	\]
	a więc $x=z$, wbrew założeniu. Zatem co najmniej jedna z liczb
	$p(x,y)$ i $p(y,z)$ jest równa $1$, a więc
	\[
	p(x,y)+p(y,z)-p(y,y)\geq1=p(x,z).
	\]
	W obu przypadkach zachodzi (P4).
	
	Zatem funkcja $p$ spełnia warunki (P1)--(P4), a więc jest metryką
	częściową na $D_\bt$.
\end{przyklad}

\begin{uwaga}
	Warto zauważyć, że
	\[
	p(\bt,\bt)=1,
	\qquad
	p(d,d)=0
	\quad(d\in D).
	\]
	Element $\bt$ ma więc dodatni rozmiar, podczas gdy elementy zbioru $D$
	są całkowicie określone.
\end{uwaga}

%==============================================================================
\subsection{Porządek indukowany przez metrykę częściową}
%==============================================================================

Metryka częściowa niesie ze sobą porządek. 

\begin{definicja}\label{def:porzadek_indukowany_pmetryka}
	Niech $p$ będzie metryką częściową na zbiorze $U$. Relację
	$\sqsubseteq_p \, \subseteq U \times U$ definiujemy wzorem
	\[
	x\sqsubseteq_p y
	\iff
	p(x,x)=p(x,y).
	\]
	Relację tę nazywamy \emph{porządkiem indukowanym przez metrykę częściową}.
\end{definicja}

\begin{twierdzenie}\label{tw:porzadek_pmetryka}
	Dla każdej metryki częściowej $p$ relacja $\sqsubseteq_p$ jest
	porządkiem częściowym.
\end{twierdzenie}

\begin{proof}
	Musimy wykazać zwrotność, antysymetryczność oraz przechodniość.
	
	\emph{Zwrotność.}
	Dla każdego $x\in U$ zachodzi
	\[
	p(x,x)=p(x,x),
	\]
	a zatem, zgodnie z definicją relacji $\sqsubseteq_p$,
	\[
	x\sqsubseteq_p x.
	\]
	
	\emph{Antysymetryczność.}
	Załóżmy, że
	\[
	x\sqsubseteq_p y
	\qquad\text{oraz}\qquad
	y\sqsubseteq_p x.
	\]
	Z definicji $\sqsubseteq_p$ otrzymujemy
	\[
	p(x,x)=p(x,y)
	\]
	oraz
	\[
	p(y,y)=p(y,x).
	\]
	Z warunku (P3), czyli symetrii pmetryki,
	\[
	p(y,x)=p(x,y).
	\]
	Stąd
	\[
	p(x,x)=p(x,y)=p(y,y).
	\]
	Na mocy warunku (P1) wynika z tego
	\[
	x=y.
	\]
	
	\emph{Przechodniość.}
	Załóżmy, że
	\[
	x\sqsubseteq_p y
	\qquad\text{oraz}\qquad
	y\sqsubseteq_p z.
	\]
	Wówczas
	\[
	p(x,x)=p(x,y)
	\]
	oraz
	\[
	p(y,y)=p(y,z).
	\]
	Z warunku (P4) otrzymujemy
	\[
	\begin{aligned}
		p(x,z)
		&\leq p(x,y)+p(y,z)-p(y,y)\\
		&=p(x,y)\\
		&=p(x,x).
	\end{aligned}
	\]
	Z drugiej strony, z warunku (P2) mamy
	\[
	p(x,x)\leq p(x,z).
	\]
	Zatem
	\[
	p(x,z)=p(x,x),
	\]
	czyli, z definicji $\sqsubseteq_p$,
	\[
	x\sqsubseteq_p z.
	\]
	
	Relacja $\sqsubseteq_p$ jest zatem zwrotna, antysymetryczna
	i przechodnia, a więc jest porządkiem częściowym.
\end{proof}

\begin{przyklad}\label{prz:porzadek_pmetryka_plaska}
	Dla pmetryki z Przykładu~\ref{prz:plaska_pmetryka} porządek
	indukowany przez pmetrykę jest dokładnie porządkiem płaskim
	z Definicji~\ref{def:zbior_plaski}.
	
	Istotnie, dla każdego $x\in D_\bt$ mamy
	\[
	p(\bt,\bt)=p(\bt,x)=1,
	\]
	a więc
	\[
	\bt\sqsubseteq_p x.
	\]
	Z kolei dla $x,y\in D$, $x\neq y$, zachodzi
	\[
	p(x,x)=0
	\neq
	1=p(x,y),
	\]
	więc
	\[
	x\not\sqsubseteq_p y.
	\]
	Podobnie dla $d\in D$ mamy
	\[
	p(d,d)=0
	\neq
	1=p(d,\bt),
	\]
	a więc $d\not\sqsubseteq_p\bt$. Zatem jedynym elementem mniejszym od
	pozostałych jest $\bt$, a dwa różne elementy zbioru $D$ są
	nieporównywalne --- relacja $\sqsubseteq_p$ spełnia warunek
	$d_1\sqsubseteq_p d_2\iff d_1=\bt\ \text{lub}\ d_1=d_2$.
\end{przyklad}

%==============================================================================
\subsection{Ciągi Cauchy'ego i zupełność}
%==============================================================================

W podejściu porządkowym podstawowym pojęciem związanym z istnieniem granic
był łańcuch oraz jego kres górny. To właśnie istnienie kresów górnych
łańcuchów stanowiło istotę zupełności zbioru łańcuchowo zupełnego.

W podejściu metrycznym rolę przybliżeń pełni natomiast ciąg elementów,
których wzajemne odległości stają się coraz mniejsze. Prowadzi to do
pojęcia ciągu Cauchy'ego. Zupełność metryki częściowej będzie oznaczała,
że każdy taki ciąg posiada granicę należącą do rozważanego zbioru.

Istotne jest przy tym, że porządek indukowany przez metrykę częściową
pozwala połączyć oba sposoby opisu. Łańcuch względem tego porządku okazuje
się szczególnym przypadkiem ciągu Cauchy'ego. W przypadku zupełnej
metryki częściowej jego granica odpowiada następnie kresowi górnemu
łańcucha.

Wprowadzimy najpierw pojęcie ciągu Cauchy'ego i zupełności metryki
częściowej, a następnie pokażemy bezpośrednio, w jaki sposób wcześniejsze
łańcuchy pojawiają się w tym metrycznym ujęciu.

\begin{definicja}\label{def:zbieznosc_pmetryka}
	Niech $p$ będzie metryką częściową na zbiorze $U$. Ciąg
	$(x_n)_{n\in\mathbb{N}}$ nazywamy \emph{zbieżnym} do elementu
	$x\in U$, jeżeli
	\[
	\lim_{n\to\infty}p(x_n,x)=p(x,x).
	\]
	Wówczas piszemy
	\[
	x_n\longrightarrow x.
	\]
\end{definicja}

Wprowadźmy teraz odpowiednik klasycznego pojęcia ciągu Cauchy'ego.

\begin{definicja}\label{def:cauchy_pmetryka}
	Niech $p$ będzie metryką częściową na zbiorze $U$. Ciąg
	$(x_n)_{n\in\mathbb{N}}$ nazywamy \emph{ciągiem Cauchy'ego}, jeżeli
	istnieje granica
	\[
	\lim_{n,m\to\infty}p(x_n,x_m).
	\]
\end{definicja}

W zwykłej przestrzeni metrycznej zbieżność ciągu Cauchy'ego prowadzi do
granicy, której odległość od samej siebie wynosi zero. W przypadku
metryki częściowej granica może mieć jednak dodatni rozmiar. Dlatego
wymagamy dodatkowo, aby rozmiar punktu granicznego był zgodny z granicą
wzajemnych odległości ciągu.

\begin{definicja}\label{def:zupelnosc_pmetryka}
	Metrykę częściową $p$ na zbiorze $U$ nazywamy \emph{zupełną}, jeżeli
	każdy ciąg Cauchy'ego $(x_n)_{n\in\mathbb{N}}$ jest zbieżny do pewnego
	$x\in U$ oraz
	\[
	p(x,x)
	=
	\lim_{n,m\to\infty}p(x_n,x_m).
	\]
\end{definicja}

Mamy już wszystkie pojęcia potrzebne do porównania zupełności
metrycznej z zupełnością porządkową. Wcześniej rozważaliśmy łańcuchy
postaci
\[
x_1\sqsubseteq x_2\sqsubseteq x_3\sqsubseteq\ldots
\]
i wymagaliśmy istnienia ich kresu górnego. Teraz pokażemy, że w obecnym
ujęciu łańcuch względem porządku $\sqsubseteq_p$ jest jednocześnie
ciągiem Cauchy'ego względem metryki częściowej $p$.

Jest to pierwszy bezpośredni związek pomiędzy obiema konstrukcjami:
porządek informacyjny może być odczytywany z pmetryki, a rosnący ciąg
informacyjnie coraz bardziej określonych przybliżeń staje się ciągiem
Cauchy'ego.

\begin{twierdzenie}\label{tw:lanuch_jest_cauchy}
	Niech $p$ będzie metryką częściową na zbiorze $U$, a
	\[
	x_1\sqsubseteq_p x_2\sqsubseteq_p x_3\sqsubseteq_p\ldots
	\]
	będzie łańcuchem względem porządku indukowanego przez $p$. Wówczas
	ciąg $(x_n)_{n\in\mathbb N}$ jest ciągiem Cauchy'ego względem $p$.
\end{twierdzenie}

\begin{proof}
	Z definicji porządku indukowanego przez pmetrykę, dla $n\leq m$ mamy
	\[
	x_n\sqsubseteq_p x_m
	\iff
	p(x_n,x_n)=p(x_n,x_m).
	\]
	Zatem
	\[
	p(x_n,x_m)=p(x_n,x_n).
	\]
	
	Z warunku (P2) oraz symetrii (P3) otrzymujemy dla $n\leq m$
	\[
	p(x_m,x_m)
	\leq
	p(x_m,x_n)
	=
	p(x_n,x_m)
	=
	p(x_n,x_n).
	\]
	A więc ciąg
	\[
	p(x_1,x_1),p(x_2,x_2),p(x_3,x_3),\ldots
	\]
	jest nierosnący. Ponieważ wartości pmetryki są nieujemne, ciąg ten
	jest ograniczony z dołu przez $0$. Istnieje zatem granica
	\[
	L:=\lim_{n\to\infty}p(x_n,x_n).
	\]
	
	Pokażemy teraz, że
	\[
	\lim_{n,m\to\infty}p(x_n,x_m)=L.
	\]
	Niech $\varepsilon>0$. Z definicji granicy ciągu
	$\bigl(p(x_n,x_n)\bigr)_{n\in\mathbb N}$ istnieje $N\in\mathbb N$
	takie, że dla każdego $k\geq N$ zachodzi
	\[
	|p(x_k,x_k)-L|<\varepsilon.
	\]
	
	Jeżeli $n,m\geq N$, to rozważamy mniejszy z indeksów
	\[
	k=\min\{n,m\}.
	\]
	Mamy $k\geq N$. Jeżeli $n\leq m$, to
	\[
	p(x_n,x_m)=p(x_n,x_n)=p(x_k,x_k),
	\]
	a jeżeli $m\leq n$, to z symetrii pmetryki
	\[
	p(x_n,x_m)=p(x_m,x_n)=p(x_m,x_m)=p(x_k,x_k).
	\]
	W obu przypadkach
	\[
	p(x_n,x_m)=p(x_k,x_k).
	\]
	Zatem
	\[
	|p(x_n,x_m)-L|
	=
	|p(x_k,x_k)-L|
	<\varepsilon.
	\]
	
	Otrzymaliśmy więc
	\[
	\forall\varepsilon>0\;\exists N\in\mathbb N\;
	\forall n,m\geq N\quad
	|p(x_n,x_m)-L|<\varepsilon,
	\]
	czyli
	\[
	\lim_{n,m\to\infty}p(x_n,x_m)=L.
	\]
	Z definicji ciągu Cauchy'ego wynika, że
	$(x_n)_{n\in\mathbb N}$ jest ciągiem Cauchy'ego.
\end{proof}

Dotychczas pokazaliśmy, że porządek indukowany przez metrykę częściową
pozwala interpretować łańcuchy jako ciągi Cauchy'ego. W przypadku
zupełnej metryki częściowej każdy łańcuch ma więc granicę. Pozostaje
zbadać, jak ta granica odnosi się do porządku $\sqsubseteq_p$, tzn.\ czy
jest ograniczeniem górnym łańcucha i czy jest najmniejszym spośród
takich ograniczeń. Okazuje się, że obie własności zachodzą.

\begin{twierdzenie}\label{tw:granica_kresem_gornym}
	Niech $p$ będzie zupełną metryką częściową na zbiorze $U$. Wówczas
	każdy łańcuch
	\[
	x_1\sqsubseteq_p x_2\sqsubseteq_p x_3\sqsubseteq_p\ldots
	\]
	posiada kres górny względem porządku $\sqsubseteq_p$, równy granicy
	tego łańcucha w sensie Definicji~\ref{def:zbieznosc_pmetryka}.
\end{twierdzenie}

\begin{proof}
	Niech
	\[
	x_1\sqsubseteq_p x_2\sqsubseteq_p x_3\sqsubseteq_p\ldots
	\]
	będzie łańcuchem względem $\sqsubseteq_p$.
	
	\emph{Krok 1: istnienie granicy.} Na mocy
	Twierdzenia~\ref{tw:lanuch_jest_cauchy} ciąg $(x_n)_{n\in\mathbb N}$
	jest ciągiem Cauchy'ego. W dowodzie tego twierdzenia pokazaliśmy
	ponadto, że ciąg $\big(p(x_n,x_n)\big)_{n\in\mathbb N}$ jest
	nierosnący i ograniczony z dołu przez $0$, więc ma granicę
	\[
	L=\lim_{n\to\infty}p(x_n,x_n),
	\]
	oraz że
	\[
	\lim_{n,m\to\infty}p(x_n,x_m)=L.
	\]
	Ponieważ metryka częściowa $p$ jest zupełna
	(Definicja~\ref{def:zupelnosc_pmetryka}), istnieje $x\in U$ takie, że
	\[
	x_n\longrightarrow x,
	\quad\text{tzn.}\quad
	\lim_{n\to\infty}p(x_n,x)=p(x,x),
	\qquad\text{oraz}\qquad
	p(x,x)=\lim_{n,m\to\infty}p(x_n,x_m)=L.
	\]
	
	\emph{Krok 2: $x$ jest ograniczeniem górnym łańcucha.} Ustalmy
	$n\in\mathbb N$ i weźmy dowolne $m\geq n$. Ponieważ
	$x_n\sqsubseteq_p x_m$, z definicji porządku indukowanego
	\[
	p(x_n,x_m)=p(x_n,x_n).
	\]
	Warunek (P4) zastosowany do elementów $x_n$, $x_m$ oraz $x$ daje
	\[
	p(x_n,x)
	\leq
	p(x_n,x_m)+p(x_m,x)-p(x_m,x_m)
	=
	p(x_n,x_n)+p(x_m,x)-p(x_m,x_m).
	\]
	Nierówność ta zachodzi dla każdego $m\geq n$, a jej lewa strona nie
	zależy od $m$. Przejdźmy z $m$ do nieskończoności po prawej stronie:
	\[
	p(x_m,x)\xrightarrow[m\to\infty]{}p(x,x)=L,
	\qquad
	p(x_m,x_m)\xrightarrow[m\to\infty]{}L.
	\]
	Granica prawej strony wynosi więc
	\[
	p(x_n,x_n)+L-L=p(x_n,x_n),
	\]
	a ponieważ nierówności nieostre zachowują się przy przejściu do
	granicy,
	\[
	p(x_n,x)\leq p(x_n,x_n).
	\]
	Z drugiej strony warunek (P2) daje
	\[
	p(x_n,x_n)\leq p(x_n,x).
	\]
	Obie nierówności razem dają
	\[
	p(x_n,x_n)=p(x_n,x),
	\]
	czyli z definicji porządku indukowanego $x_n\sqsubseteq_p x$.
	Ponieważ $n$ było dowolne, $x$ jest ograniczeniem górnym łańcucha.
	
	\emph{Krok 3: $x$ jest najmniejszym ograniczeniem górnym.} Niech
	$y\in U$ będzie dowolnym ograniczeniem górnym łańcucha, tzn.\
	$x_n\sqsubseteq_p y$ dla każdego $n$. Z definicji porządku
	indukowanego
	\[
	p(x_n,y)=p(x_n,x_n)
	\qquad(n\in\mathbb N).
	\]
	Warunek (P4) zastosowany do elementów $x$, $x_n$ oraz $y$ daje
	\[
	p(x,y)
	\leq
	p(x,x_n)+p(x_n,y)-p(x_n,x_n)
	=
	p(x,x_n)
	\qquad(n\in\mathbb N).
	\]
	Z symetrii (P3) oraz zbieżności $x_n\to x$ mamy
	\[
	p(x,x_n)=p(x_n,x)\xrightarrow[n\to\infty]{}p(x,x).
	\]
	Przechodząc w nierówności $p(x,y)\leq p(x,x_n)$ do granicy przy
	$n\to\infty$ (lewa strona nie zależy od $n$), otrzymujemy
	\[
	p(x,y)\leq p(x,x).
	\]
	Z warunku (P2) zachodzi nierówność przeciwna $p(x,x)\leq p(x,y)$,
	więc
	\[
	p(x,x)=p(x,y),
	\]
	czyli $x\sqsubseteq_p y$.
	
	Element $x$ jest więc ograniczeniem górnym łańcucha i jest mniejszy
	lub równy każdemu innemu ograniczeniu górnemu. Zatem
	\[
	x=\bigsqcup_{n=1}^{\infty}x_n,
	\]
	czyli kres górny łańcucha istnieje i jest równy jego granicy.
\end{proof}

\begin{uwaga}
	Powyższe twierdzenie pokazuje, w jakim sensie zupełność metryki
	częściowej odtwarza zupełność porządkową: w zbiorze $(U,\sqsubseteq_p)$
	każdy łańcuch ma kres górny, a jest nim granica łańcucha.
	
	Nie oznacza to jednak, że $(U,\sqsubseteq_p)$ jest zbiorem łańcuchowo
	zupełnym w sensie rozdziału~\ref{sec:cpo},
	ponieważ porządek $\sqsubseteq_p$ nie musi mieć elementu najmniejszego.
	Przykładowo, dla zwykłej metryki $d$ (tzn.\ gdy $d(x,x)=0$ dla
	wszystkich $x$) mamy
	\[
	x\sqsubseteq_d y
	\iff
	d(x,y)=0
	\iff
	x=y,
	\]
	więc porządek indukowany jest dyskretny i przy co najmniej dwóch
	elementach zbioru $U$ nie ma elementu najmniejszego.
\end{uwaga}

\subsection{Kontrakcje i twierdzenie o punkcie stałym}

W poprzednich podrozdziałach wprowadziliśmy pojęcie metryki częściowej,
porządek indukowany przez pmetrykę oraz pojęcia zbieżności, ciągu
Cauchy'ego i zupełności. Możemy teraz przejść do właściwego zastosowania
tego aparatu, czyli wyznaczania punktów stałych.

W klasycznej przestrzeni metrycznej kluczową rolę odgrywa twierdzenie
Banacha o punkcie stałym. Jego założeniem jest istnienie kontrakcji, czyli
funkcji, która zmniejsza odległości między elementami co najmniej przez
stały czynnik mniejszy od $1$. Analogiczne pojęcie można wprowadzić dla
metryki częściowej.

W dalszej części pokażemy, że zupełność metryki częściowej pozwala
zastosować iterację funkcji
\[
x,\quad f(x),\quad f^2(x),\quad f^3(x),\ldots
\]
do konstrukcji punktu stałego. Otrzymamy w ten sposób odpowiednik
klasycznej metody iteracji Banacha, sformułowany dla metryk częściowych.

\begin{definicja}\label{def:kontrakcja_pmetryka}
	Niech $p$ będzie metryką częściową na zbiorze $U$. Funkcję
	$f:U\to U$ nazywamy \emph{kontrakcją} względem $p$, jeżeli istnieje
	stała $c\in[0,1)$ taka, że
	\[
	p(f(x),f(y))\leq c \cdot p(x,y)
	\]
	dla wszystkich $x,y\in U$.
\end{definicja}

Kontrakcja pozwala kontrolować odległości pomiędzy kolejnymi iteracjami
funkcji. Ustalmy zatem dowolny element początkowy $x_0\in U$ i rozważmy
ciąg określony rekurencyjnie przez
\[
x_{n+1}=f(x_n),
\qquad n\in\mathbb N_0.
\]
Równoważnie,
\[
x_n=f^n(x_0).
\]

W pierwszej kolejności pokażemy, że tak otrzymany ciąg jest ciągiem
Cauchy'ego.

\begin{lemat}\label{lem:iteracje_kontrakcji_cauchy}
	Niech $(U,p)$ będzie przestrzenią częściowo metryczną, a
	$f:U\to U$ kontrakcją względem $p$ ze współczynnikiem
	$c\in[0,1)$. Dla ustalonego $x_0\in U$ niech
	\[
	x_{n+1}=f(x_n),
	\qquad n\in\mathbb N_0.
	\]
	Wówczas ciąg $(x_n)_{n\in\mathbb N_0}$ jest ciągiem Cauchy'ego.
\end{lemat}

\begin{proof}
	Z warunku kontrakcji dla $x_n$ i $x_{n-1}$ otrzymujemy dla każdego
	$n\geq 1$
	\[
	\begin{aligned}
		p(x_{n+1},x_n)
		&=
		p(f(x_n),f(x_{n-1}))\\
		&\leq
		c\,p(x_n,x_{n-1}).
	\end{aligned}
	\]
	Stosując tę nierówność indukcyjnie, otrzymujemy
	\[
	p(x_{n+1},x_n)
	\leq
	c^n p(x_1,x_0)
	\]
	dla każdego $n\in\mathbb N_0$.
	
	Niech teraz $m>n$. Wielokrotne zastosowanie warunku (P4) daje
	\[
	\begin{aligned}
		p(x_n,x_m)
		&\leq
		p(x_n,x_{n+1})
		+p(x_{n+1},x_m)
		-p(x_{n+1},x_{n+1})\\
		&\leq
		p(x_n,x_{n+1})+p(x_{n+1},x_m),
	\end{aligned}
	\]
	gdzie w drugim kroku wykorzystaliśmy nieujemność
	$p(x_{n+1},x_{n+1})$.
	
	Kontynuując to rozumowanie, otrzymujemy
	\[
	p(x_n,x_m)
	\leq
	\sum_{k=n}^{m-1}p(x_k,x_{k+1}).
	\]
	Z uzyskanego wcześniej oszacowania wynika więc
	\[
	\begin{aligned}
		p(x_n,x_m)
		&\leq
		p(x_0,x_1)\sum_{k=n}^{m-1}c^k\\
		&\leq
		p(x_0,x_1)\sum_{k=n}^{\infty}c^k\\
		&=
		p(x_0,x_1)\frac{c^n}{1-c}.
	\end{aligned}
	\]
	
	Ponieważ $0\leq c<1$, zachodzi
	\[
	\lim_{n\to\infty}
	p(x_0,x_1)\frac{c^n}{1-c}=0.
	\]
	Zatem dla każdego $\varepsilon>0$ istnieje $N\in\mathbb N$ takie,
	że dla każdego $n\geq N$
	\[
	p(x_0,x_1)\frac{c^n}{1-c}<\varepsilon.
	\]
	Jeżeli $n,m\geq N$, to bez straty ogólności możemy założyć, że
	$n\leq m$. Wówczas
	\[
	p(x_n,x_m)
	\leq
	p(x_0,x_1)\frac{c^n}{1-c}
	<\varepsilon.
	\]
	Z symetrii pmetryki ten sam wniosek zachodzi również dla $m<n$.
	Otrzymujemy zatem
	\[
	\forall\varepsilon>0\;\exists N\in\mathbb N\;
	\forall n,m\geq N\quad
	p(x_n,x_m)<\varepsilon.
	\]
	W konsekwencji
	\[
	\lim_{n,m\to\infty}p(x_n,x_m)=0.
	\]
	Granica wzajemnych odległości istnieje, a więc z definicji
	$(x_n)_{n\in\mathbb N_0}$ jest ciągiem Cauchy'ego.
\end{proof}

Powyższy lemat pokazuje, że iteracje kontrakcji tworzą ciąg Cauchy'ego.
Jeżeli przestrzeń jest zupełna, ciąg ten posiada zatem granicę należącą
do $U$. Co więcej, z definicji zupełności własna odległość tej granicy
jest równa granicy wzajemnych odległości wyrazów ciągu.

W naszym przypadku granica ta jest równa zeru, co wynika bezpośrednio
z oszacowania uzyskanego w dowodzie lematu. Otrzymujemy w ten sposób
punkt, o którym następnie wykażemy, że jest punktem stałym funkcji
$f$.

\begin{twierdzenie}[Matthews, \cite{matthews1994}]
	\label{tw:punkt_staly_matthews}
	Niech $(U,p)$ będzie zupełną przestrzenią metryczną częściową, a
	$f:U\to U$ będzie kontrakcją względem $p$, tzn. istnieje
	$c\in[0,1)$ takie, że
	\[
	p(f(x),f(y))\leq c\,p(x,y)
	\]
	dla wszystkich $x,y\in U$.
	
	Wówczas istnieje dokładnie jeden punkt $x^*\in U$ spełniający
	\[
	f(x^*)=x^*.
	\]
	Ponadto
	\[
	p(x^*,x^*)=0.
	\]
\end{twierdzenie}

\begin{proof}
	Wybierzmy dowolny element $x_0\in U$ i zdefiniujmy ciąg
	$(x_n)_{n\in\mathbb N_0}$ przez
	\[
	x_{n+1}=f(x_n).
	\]
	
	Z Lematu~\ref{lem:iteracje_kontrakcji_cauchy} wynika, że
	$(x_n)_{n\in\mathbb N_0}$ jest ciągiem Cauchy'ego. Ponieważ $p$ jest
	zupełna, istnieje $x^*\in U$ takie, że
	\[
	x_n\longrightarrow x^*
	\]
	oraz
	\[
	p(x^*,x^*)
	=
	\lim_{n,m\to\infty}p(x_n,x_m).
	\]
	Z Lematu~\ref{lem:iteracje_kontrakcji_cauchy} mamy
	\[
	\lim_{n,m\to\infty}p(x_n,x_m)=0,
	\]
	a więc
	\[
	p(x^*,x^*)=0.
	\]
	
	Pokażemy teraz, że $x^*$ jest punktem stałym funkcji $f$. Z warunku
	(P4) otrzymujemy
	\[
	p(f(x^*),x^*)
	\leq
	p(f(x^*),x_{n+1})
	+
	p(x_{n+1},x^*)
	-
	p(x_{n+1},x_{n+1}).
	\]
	Stąd, korzystając z nieujemności pmetryki oraz z równości
	$x_{n+1}=f(x_n)$,
	\[
	\begin{aligned}
		p(f(x^*),x^*)
		&\leq
		p(f(x^*),f(x_n))+p(x_{n+1},x^*)\\
		&\leq
		c\,p(x^*,x_n)+p(x_{n+1},x^*).
	\end{aligned}
	\]
	Przechodząc z $n$ do nieskończoności, otrzymujemy
	\[
	p(f(x^*),x^*)\leq 0.
	\]
	Ponieważ pmetryka przyjmuje wartości nieujemne,
	\[
	p(f(x^*),x^*)=0.
	\]
	Z warunku (P2) wynika
	\[
	p(f(x^*),f(x^*))\leq p(f(x^*),x^*)=0,
	\]
	a zatem
	\[
	p(f(x^*),f(x^*))=0.
	\]
	Wraz z
	\[
	p(x^*,x^*)=0
	\]
	otrzymujemy
	\[
	p(f(x^*),f(x^*))
	=
	p(f(x^*),x^*)
	=
	p(x^*,x^*)
	=
	0.
	\]
	Na mocy (P1)
	\[
	f(x^*)=x^*.
	\]
	
	Pozostaje wykazać jednoznaczność. Niech $y\in U$ będzie drugim
	punktem stałym funkcji $f$. Wówczas
	\[
	\begin{aligned}
		p(x^*,y)
		&=
		p(f(x^*),f(y))\\
		&\leq
		c\,p(x^*,y).
	\end{aligned}
	\]
	Ponieważ $c<1$ oraz $p(x^*,y)\geq0$, otrzymujemy
	\[
	p(x^*,y)=0.
	\]
	Z warunku (P2) wynika
	\[
	p(x^*,x^*)=0,
	\qquad
	p(y,y)=0,
	\]
	a więc z (P1)
	\[
	x^*=y.
	\]
	Punkt stały jest zatem jedyny.
\end{proof}

\subsection{Porównanie z podejściem porządkowym}

\begin{comment}
W tej części przedstawimy związek pomiędzy strukturą porządkową
wykorzystywaną wcześniej a podejściem metrycznym rozwijanym w teorii
dziedzin ilościowych. Podstawowe pojęcia dotyczące zbiorów skierowanych,
zupełnych zbiorów częściowo uporządkowanych, relacji $\ll$, dziedzin
ciągłych oraz topologii Scotta należą do klasycznej teorii dziedzin.
W dalszej części korzystamy z ich ujęcia przedstawionego między innymi
w pracach Waszkiewicza
\cite{Waszkiewicz2001,Waszkiewicz2003}.

Szczególne znaczenie dla dalszych rozważań ma praca
\cite{Waszkiewicz2001}, w której badany jest związek pomiędzy pomiarami
na dziedzinach, topologią Scotta oraz metrykami częściowymi. Przedstawiona
tam konstrukcja pozwala, przy odpowiednich założeniach dotyczących
dziedziny i pomiaru, otrzymać metrykę częściową zgodną ze strukturą
porządkową. Zagadnienie to zostało następnie rozwinięte w
\cite{Waszkiewicz2003}, gdzie analizowany jest związek pomiędzy
pomiarami a metrykami częściowymi w ramach ilościowej teorii dziedzin.

W kolejnych podrozdziałach najpierw przypomnimy potrzebne pojęcia
porządkowe, a następnie przejdziemy od relacji $\ll$ i topologii Scotta
do pojęcia pomiaru. Dopiero na tej podstawie skonstruujemy odpowiednią
metrykę częściową. Dzięki takiej kolejności konstrukcja metryczna nie
jest wprowadzana niezależnie od wcześniej omawianej struktury
porządkowej, lecz wynika z informacji zawartej w tej strukturze.

\subsubsection{Zbiory skierowane i dziedziny}

W rozdziałach 5--7 rozważaliśmy zbiory łańcuchowo zupełne, których struktura porządkowa była wystarczająca do sformułowania twierdzenia Kleenego i wyznaczenia najmniejszego punktu stałego odpowiedniego funkcjonału. W dalszej części potrzebna będzie szersza struktura porządkowa, w której zamiast łańcuchów rozważa się zbiory skierowane.

\begin{definicja}\label{def:zbior_skierowany}
	Niech $(P,\sqsubseteq)$ będzie zbiorem częściowo uporządkowanym. Niepusty podzbiór $D\subseteq P$ nazywamy \emph{zbiorem skierowanym}, jeżeli dla każdych $x,y\in D$ istnieje element $z\in D$ taki, że
	\[
	x\sqsubseteq z
	\qquad\text{oraz}\qquad
	y\sqsubseteq z.
	\]
\end{definicja}

Warunek ten oznacza, że dowolne dwa elementy zbioru skierowanego mają wspólne ograniczenie górne należące do tego samego zbioru. Nie wymagamy przy tym, aby elementy zbioru były ze sobą porównywalne. Z tego względu pojęcie zbioru skierowanego jest szersze niż pojęcie łańcucha.

\begin{przyklad}\label{prz:zbior_skierowany_nie_lancuch}
	Rozważmy zbiór
	\[
	P=\{\bot,a,b,c\}
	\]
	z porządkiem określonym przez relacje
	\[
	\bot\sqsubseteq a\sqsubseteq c,
	\qquad
	\bot\sqsubseteq b\sqsubseteq c,
	\]
	przy czym elementy $a$ i $b$ są nieporównywalne:
	\[
	a\not\sqsubseteq b,
	\qquad
	b\not\sqsubseteq a.
	\]
	
	Strukturę tego porządku można przedstawić za pomocą diagramu Hassego:
	
	\[
	\begin{tikzpicture}[
		every node/.style={circle, fill=white, inner sep=2pt},
		>=stealth
		]
		\node (bottom) at (0,0) {$\bot$};
		\node (a) at (-1.5,1.5) {$a$};
		\node (b) at (1.5,1.5) {$b$};
		\node (c) at (0,3) {$c$};
		
		\draw (bottom) -- (a);
		\draw (bottom) -- (b);
		\draw (a) -- (c);
		\draw (b) -- (c);
	\end{tikzpicture}
	\]
	
	Zbiór
	\[
	D=\{a,b,c\}
	\]
	nie jest łańcuchem, ponieważ $a$ i $b$ nie są porównywalne. Jest jednak zbiorem skierowanym. Dla pary $a,b$ elementem będącym ich wspólnym ograniczeniem górnym i należącym do $D$ jest $c$. Pozostałe pary elementów są już porównywalne. Zatem $D$ spełnia warunek z definicji~\ref{def:zbior_skierowany}.
\end{przyklad}

Każdy łańcuch jest zbiorem skierowanym.

\begin{twierdzenie}\label{tw:lancuch_skierowany}
	Niech $(P,\sqsubseteq)$ będzie zbiorem częściowo uporządkowanym. Każdy łańcuch w $P$ jest zbiorem skierowanym.
\end{twierdzenie}

\begin{proof}
	Niech $C\subseteq P$ będzie łańcuchem. Z definicji łańcucha dla dowolnych $x,y\in C$ zachodzi
	\[
	x\sqsubseteq y
	\qquad\text{lub}\qquad
	y\sqsubseteq x.
	\]
	W pierwszym przypadku możemy przyjąć $z=y$, natomiast w drugim przypadku $z=x$. W obu przypadkach $z\in C$ oraz
	\[
	x\sqsubseteq z
	\qquad\text{oraz}\qquad
	y\sqsubseteq z.
	\]
	Zatem $C$ spełnia warunek z definicji~\ref{def:zbior_skierowany}, a więc jest zbiorem skierowanym.
\end{proof}

Pojęcie zbioru skierowanego pozwala rozszerzyć pojęcie zupełności, które w przypadku zbiorów łańcuchowo zupełnych dotyczyło jedynie łańcuchów. Otrzymujemy w ten sposób następującą definicję.

\begin{definicja}\label{def:zupełny_zbior_czesciowo_uporzadkowany}
	Zbiór częściowo uporządkowany $(P,\sqsubseteq)$ nazywamy \emph{zupełnym zbiorem częściowo uporządkowanym}, jeżeli każdy niepusty zbiór skierowany $D\subseteq P$ posiada supremum w $P$, oznaczane przez
	\[
	\bigsqcup D.
	\]
\end{definicja}

Każdy zupełny zbiór częściowo uporządkowany jest zatem także łańcuchowo zupełny.

\begin{twierdzenie}\label{tw:zupelnosc_skierowana_lancuchowa}
	Jeżeli $(P,\sqsubseteq)$ jest zupełnym zbiorem częściowo uporządkowanym, to każdy łańcuch $C\subseteq P$ posiada supremum w $P$.
\end{twierdzenie}

\begin{proof}
	Na mocy twierdzenia~\ref{tw:lancuch_skierowany} każdy łańcuch $C\subseteq P$ jest zbiorem skierowanym. Z definicji zupełnego zbioru częściowo uporządkowanego~\ref{def:zupełny_zbior_czesciowo_uporzadkowany} każdy niepusty zbiór skierowany posiada supremum w $P$. W szczególności istnieje
	\[
	\bigsqcup C\in P.
	\]
	Zatem każdy łańcuch w $P$ posiada supremum, czyli $(P,\sqsubseteq)$ jest łańcuchowo zupełny.
\end{proof}

Wynika stąd, że zupełność względem zbiorów skierowanych jest warunkiem silniejszym niż zupełność względem łańcuchów. W teorii dziedzin struktura ta jest podstawowa, ponieważ pozwala traktować nie tylko liniowo uporządkowane ciągi coraz dokładniejszych przybliżeń, lecz również rodziny przybliżeń, których elementy nie muszą być wzajemnie porównywalne.

\subsubsection{Relacja way-below}

Sama zupełność względem zbiorów skierowanych nie pozwala jeszcze
scharakteryzować sposobu, w jaki elementy dziedziny mogą być
przybliżane przez elementy prostsze. Do tego celu służy relacja
\emph{way-below}, oznaczana przez $\ll$. Relacja ta jest jednym
z podstawowych narzędzi teorii dziedzin i pozwala formalnie opisać
pojęcie przybliżenia elementu za pomocą elementów zawartych w
skierowanych rodzinach przybliżeń
\cite{Waszkiewicz2001}.

\begin{definicja}\label{def:relacja_way_below}
	Niech $(P,\sqsubseteq)$ będzie zupełnym zbiorem częściowo uporządkowanym.
	Dla $x,y\in P$ mówimy, że $x$ jest \emph{znacznie poniżej} $y$,
	co oznaczamy przez
	\[
	x\ll y,
	\]
	jeżeli dla każdego niepustego zbioru skierowanego $D\subseteq P$,
	jeżeli
	\[
	y\sqsubseteq\bigsqcup D,
	\]
	to istnieje element $d\in D$ taki, że
	\[
	x\sqsubseteq d.
	\]
\end{definicja}

Definicja~\ref{def:relacja_way_below} mówi, że element $x$ jest
przybliżeniem elementu $y$ na tyle silnym, że nie można otrzymać
$y$ jako kresu górnego zbioru skierowanego bez pojawienia się
w tym zbiorze elementu, który zawiera w sobie informację zawartą
w $x$. W szczególności relacja $\ll$ jest silniejsza od zwykłej
relacji porządku $\sqsubseteq$.

\begin{twierdzenie}\label{tw:way_below_porządek}
	Niech $(P,\sqsubseteq)$ będzie zupełnym zbiorem częściowo uporządkowanym.
	Jeżeli
	\[
	x\ll y,
	\]
	to
	\[
	x\sqsubseteq y.
	\]
	Ponadto, jeżeli
	\[
	x'\sqsubseteq x\ll y\sqsubseteq y',
	\]
	to
	\[
	x'\ll y'.
	\]
\end{twierdzenie}

\begin{proof}
	Załóżmy najpierw, że
	\[
	x\ll y.
	\]
	Rozważmy zbiór jednoelementowy
	\[
	D=\{y\}.
	\]
	Jest on skierowany, a ponadto
	\[
	\bigsqcup D=y.
	\]
	Zatem
	\[
	y\sqsubseteq\bigsqcup D.
	\]
	Na mocy definicji relacji $\ll$ istnieje $d\in D$ takie, że
	\[
	x\sqsubseteq d.
	\]
	Jedynym elementem zbioru $D$ jest $y$, więc
	\[
	x\sqsubseteq y.
	\]
	
	Załóżmy teraz, że
	\[
	x'\sqsubseteq x\ll y\sqsubseteq y'.
	\]
	Chcemy wykazać, że
	\[
	x'\ll y'.
	\]
	Niech $D\subseteq P$ będzie dowolnym niepustym zbiorem skierowanym
	spełniającym
	\[
	y'\sqsubseteq\bigsqcup D.
	\]
	Ponieważ
	\[
	y\sqsubseteq y'\sqsubseteq\bigsqcup D,
	\]
	otrzymujemy
	\[
	y\sqsubseteq\bigsqcup D.
	\]
	Z założenia $x\ll y$ istnieje zatem $d\in D$ takie, że
	\[
	x\sqsubseteq d.
	\]
	Ponieważ
	\[
	x'\sqsubseteq x,
	\]
	otrzymujemy
	\[
	x'\sqsubseteq d.
	\]
	Spełniony został warunek z definicji~\ref{def:relacja_way_below},
	a więc
	\[
	x'\ll y'.
	\]
\end{proof}

Relację $\ll$ można dobrze zobrazować na przykładzie zbioru
potęgowego.

\begin{przyklad}\label{prz:way_below_zbior_potegowy}
	Niech $X$ będzie dowolnym zbiorem. Rozważmy zbiór potęgowy
	$\mathcal{P}(X)$ wyposażony w porządek zawierania:
	\[
	A\sqsubseteq B
	\quad\Longleftrightarrow\quad
	A\subseteq B.
	\]
	Wtedy dla dowolnych $A,B\subseteq X$ zachodzi
	\[
	A\ll B
	\quad\Longleftrightarrow\quad
	A\subseteq B
	\text{ oraz }
	A\text{ jest skończony}.
	\]
\end{przyklad}

\begin{proof}
	Załóżmy najpierw, że $A$ jest skończonym podzbiorem $B$. Niech
	$\mathcal{D}\subseteq\mathcal{P}(X)$ będzie niepustą rodziną
	skierowaną względem zawierania oraz załóżmy, że
	\[
	B\subseteq\bigsqcup\mathcal{D}.
	\]
	W porządku zawierania supremum rodziny skierowanej jest jej sumą:
	\[
	\bigsqcup\mathcal{D}
	=
	\bigcup_{D\in\mathcal{D}}D.
	\]
	Zatem
	\[
	A\subseteq\bigcup_{D\in\mathcal{D}}D.
	\]
	Ponieważ $A$ jest skończony, możemy zapisać
	\[
	A=\{a_1,\ldots,a_n\}.
	\]
	Dla każdego $i\in\{1,\ldots,n\}$ istnieje $D_i\in\mathcal{D}$ takie,
	że
	\[
	a_i\in D_i.
	\]
	
	Pokażemy indukcyjnie, że istnieje $D\in\mathcal{D}$ takie, że
	\[
	D_i\subseteq D
	\]
	dla wszystkich $i\in\{1,\ldots,n\}$. Dla $n=1$ wystarczy przyjąć
	$D=D_1$. Załóżmy, że dla pewnego $k<n$ istnieje
	$D'\in\mathcal{D}$ takie, że
	\[
	D_i\subseteq D'
	\]
	dla wszystkich $i\in\{1,\ldots,k\}$. Ponieważ $\mathcal{D}$ jest
	skierowana, istnieje $D\in\mathcal{D}$ takie, że
	\[
	D'\subseteq D
	\qquad\text{oraz}\qquad
	D_{k+1}\subseteq D.
	\]
	Wtedy
	\[
	D_i\subseteq D
	\]
	dla wszystkich $i\in\{1,\ldots,k+1\}$. Zatem dla pewnego
	$D\in\mathcal{D}$ zachodzi
	\[
	A\subseteq D.
	\]
	Otrzymujemy więc
	\[
	A\sqsubseteq D,
	\]
	co na mocy definicji~\ref{def:relacja_way_below} oznacza
	\[
	A\ll B.
	\]
	
	W drugą stronę załóżmy, że
	\[
	A\ll B.
	\]
	Na mocy twierdzenia~\ref{tw:way_below_porządek} mamy
	\[
	A\subseteq B.
	\]
	Przypuśćmy, że $A$ jest nieskończony. Rozważmy rodzinę
	\[
	\mathcal{D}
	=
	\{F\subseteq B\mid F\text{ jest skończony}\}.
	\]
	Rodzina $\mathcal{D}$ jest skierowana względem zawierania. Rzeczywiście,
	dla dowolnych $F_1,F_2\in\mathcal{D}$ zbiór
	\[
	F_1\cup F_2
	\]
	jest skończony, należy do $\mathcal{D}$ oraz spełnia
	\[
	F_1\subseteq F_1\cup F_2,
	\qquad
	F_2\subseteq F_1\cup F_2.
	\]
	Ponadto
	\[
	\bigsqcup\mathcal{D}
	=
	\bigcup_{F\in\mathcal{D}}F
	=
	B.
	\]
	Zatem
	\[
	B\subseteq\bigsqcup\mathcal{D}.
	\]
	Z założenia $A\ll B$ wynika więc, że istnieje $F\in\mathcal{D}$
	takie, że
	\[
	A\subseteq F.
	\]
	Jest to niemożliwe, ponieważ $F$ jest skończony, natomiast $A$ jest
	nieskończony. Otrzymujemy sprzeczność. Zatem $A$ musi być skończony.
\end{proof}

Powyższy przykład pokazuje podstawową interpretację relacji $\ll$.
Element $B$ może być nieskończony, ale jego informacja może być
przybliżana przez skończone zbiory $A$. W tym sensie elementy
spełniające relację $A\ll B$ reprezentują skończone przybliżenia
elementu $B$.

Relacja $\ll$ posiada również własność interpolacji, która będzie
istotna przy przejściu od struktury porządkowej do topologii Scotta.

\begin{twierdzenie}\label{tw:interpolacja_way_below}
	Niech $(P,\sqsubseteq)$ będzie dziedziną ciągłą. Jeżeli
	\[
	x\ll z,
	\]
	to istnieje element $y\in P$ taki, że
	\[
	x\ll y\ll z.
	\]
\end{twierdzenie}

Twierdzenie~\ref{tw:interpolacja_way_below} jest nazywane własnością
interpolacji relacji $\ll$. Oznacza ono, że pomiędzy dowolnym
przybliżeniem $x$ a elementem $z$, który jest przez nie przybliżany,
można znaleźć pośrednie przybliżenie $y$. W szczególności relacja
$\ll$ nie opisuje jedynie pojedynczego poziomu przybliżenia, lecz
pozwala budować ich hierarchię. Własność ta jest standardową
własnością dziedzin ciągłych i będzie wykorzystywana przy opisie
topologii Scotta
\cite{Waszkiewicz2001}.

\begin{definicja}\label{def:dziedzina_ciagla}
	Niech $(P,\sqsubseteq)$ będzie zupełnym zbiorem częściowo uporządkowanym.
	Mówimy, że $P$ jest \emph{dziedziną ciągłą}, jeżeli dla każdego
	$x\in P$ zbiór
	\[
	\mathord{\Downarrow}x
	=
	\{y\in P\mid y\ll x\}
	\]
	jest skierowany oraz
	\[
	x=\bigsqcup\mathord{\Downarrow}x.
	\]
\end{definicja}

Warunek z definicji~\ref{def:dziedzina_ciagla} oznacza, że każdy
element dziedziny może zostać otrzymany jako supremum swoich
przybliżeń. Struktura porządkowa nie opisuje więc jedynie relacji
pomiędzy gotowymi elementami, lecz dostarcza również sposobu ich
stopniowego przybliżania.

\begin{definicja}\label{def:baza_dziedziny}
	Niech $(P,\sqsubseteq)$ będzie dziedziną ciągłą. Podzbiór
	$B\subseteq P$ nazywamy \emph{bazą} dziedziny $P$, jeżeli dla każdego
	$x\in P$ zbiór
	\[
	\{b\in B\mid b\ll x\}
	\]
	jest skierowany oraz
	\[
	x=
	\bigsqcup\{b\in B\mid b\ll x\}.
	\]
\end{definicja}

Jeżeli baza jest przeliczalna, otrzymujemy klasę dziedzin, dla których
struktura przybliżeń może zostać opisana za pomocą przeliczalnego
zbioru elementów.

\begin{definicja}\label{def:dziedzina_omega_ciagla}
	Dziedzinę ciągłą $P$ nazywamy \emph{dziedziną $\omega$-ciągłą}, jeżeli
	posiada ona przeliczalną bazę.
\end{definicja}


\subsubsection{Topologia Scotta}

Struktura porządkowa dziedziny pozwala nie tylko określać sposób
przybliżania elementów za pomocą relacji $\ll$, lecz również
wprowadzić topologię opisującą zachowanie tych przybliżeń. Naturalną
topologią związaną z częściowym porządkiem jest topologia Scotta
\cite{Waszkiewicz2001}.

\begin{definicja}\label{def:zbior_scotta_otwarty}
	Niech $(P,\sqsubseteq)$ będzie zupełnym zbiorem częściowo uporządkowanym.
	Podzbiór $U\subseteq P$ nazywamy \emph{Scott-otwartym}, jeżeli spełnia
	następujące dwa warunki:
	
	\begin{enumerate}
		\item $U$ jest górnym zbiorem, tzn. dla każdych $x,y\in P$ zachodzi
		\[
		x\in U\ \text{oraz}\ x\sqsubseteq y
		\quad\Longrightarrow\quad
		y\in U;
		\]
		
		\item dla każdego niepustego zbioru skierowanego $D\subseteq P$,
		jeżeli
		\[
		\bigsqcup D\in U,
		\]
		to
		\[
		D\cap U\neq\varnothing.
		\]
	\end{enumerate}
\end{definicja}

Pierwszy warunek oznacza, że jeżeli element należy do zbioru
otwartego, to każdy element zawierający co najmniej tyle samo
informacji również do niego należy. Drugi warunek opisuje zachowanie
zbioru otwartego względem granic skierowanych przybliżeń. Jeżeli
supremum zbioru skierowanego należy do $U$, to już pewien element tego
zbioru musi należeć do $U$. Oznacza to, że przynależność do zbioru
otwartego nie może pojawić się dopiero w granicy skierowanego procesu
przybliżania.

\begin{definicja}\label{def:topologia_scotta}
	Niech $(P,\sqsubseteq)$ będzie zupełnym zbiorem częściowo uporządkowanym.
	\emph{Topologią Scotta} nazywamy rodzinę
	\[
	\sigma(P)
	=
	\{U\subseteq P\mid U\text{ jest Scott-otwarty}\}.
	\]
\end{definicja}

Rodzina $\sigma(P)$ jest topologią na $P$. Zbiór pusty oraz cały
zbiór $P$ są Scott-otwarte. Dowolna suma zbiorów Scott-otwartych jest
Scott-otwarta, ponieważ własność bycia górnym zbiorem zachowuje się
przy sumach, a jeżeli supremum zbioru skierowanego należy do sumy,
to należy do jednego ze składników tej sumy. Również skończony
przekrój zbiorów Scott-otwartych jest Scott-otwarty.

Istotny przykład topologii Scotta otrzymujemy dla zbioru potęgowego.

\begin{przyklad}\label{prz:scott_potega}
	Rozważmy zupełny zbiór częściowo uporządkowany
	\[
	(\mathcal{P}(X),\subseteq).
	\]
	Dla dowolnego skończonego zbioru $A\subseteq X$ zdefiniujmy
	\[
	U_A
	=
	\{B\subseteq X\mid A\subseteq B\}.
	\]
	Zbiór $U_A$ jest Scott-otwarty.
	
	Istotnie, jeżeli $B\in U_A$ oraz $B\subseteq C$, to
	\[
	A\subseteq B\subseteq C,
	\]
	a więc $C\in U_A$. Zatem $U_A$ jest górnym zbiorem.
	
	Pozostaje sprawdzić drugi warunek z
	definicji~\ref{def:zbior_scotta_otwarty}. Niech $\mathcal{D}$ będzie
	niepustą rodziną skierowaną pod względem zawierania oraz załóżmy, że
	\[
	\bigsqcup\mathcal{D}
	=
	\bigcup_{D\in\mathcal{D}}D
	\in U_A.
	\]
	Oznacza to, że
	\[
	A\subseteq\bigcup_{D\in\mathcal{D}}D.
	\]
	Ponieważ $A$ jest skończony, dla każdego $a\in A$ istnieje
	$D_a\in\mathcal{D}$ takie, że
	\[
	a\in D_a.
	\]
	Skierowanie rodziny $\mathcal{D}$ pozwala znaleźć jeden element
	$D\in\mathcal{D}$ będący ograniczeniem górnym wszystkich skończenie
	wielu zbiorów $D_a$. W konsekwencji
	\[
	A\subseteq D,
	\]
	czyli
	\[
	D\in U_A.
	\]
	Zatem
	\[
	\mathcal{D}\cap U_A\neq\varnothing,
	\]
	a więc $U_A$ jest Scott-otwarty.
\end{przyklad}

Przykład~\ref{prz:scott_potega} pokazuje, że topologia Scotta
bezpośrednio odzwierciedla pojęcie skończonego przybliżenia. Aby
zdecydować, czy zbiór $B$ należy do otoczenia $U_A$, wystarczy
sprawdzić, czy zawiera on skończony fragment informacji określony
przez $A$.

Związek pomiędzy topologią Scotta a relacją $\ll$ staje się
szczególnie wyraźny w dziedzinach ciągłych.

\begin{twierdzenie}\label{tw:uparrow_way_below_scott}
	Niech $(P,\sqsubseteq)$ będzie dziedziną ciągłą. Dla każdego
	$x\in P$ zbiór
	\[
	\mathord{\Uparrow}x
	=
	\{y\in P\mid x\ll y\}
	\]
	jest Scott-otwarty.
\end{twierdzenie}

\begin{proof}
	Najpierw pokażemy, że $\mathord{\Uparrow}x$ jest górnym zbiorem.
	Niech $y\in\mathord{\Uparrow}x$ oraz
	\[
	y\sqsubseteq z.
	\]
	Zatem
	\[
	x\ll y.
	\]
	Na mocy twierdzenia~\ref{tw:way_below_porządek}, z relacji
	\[
	x\ll y\sqsubseteq z
	\]
	wynika
	\[
	x\ll z.
	\]
	Stąd
	\[
	z\in\mathord{\Uparrow}x.
	\]
	
	Pozostaje sprawdzić drugi warunek z
	definicji~\ref{def:zbior_scotta_otwarty}. Niech $D\subseteq P$ będzie
	niepustym zbiorem skierowanym oraz załóżmy, że
	\[
	\bigsqcup D\in\mathord{\Uparrow}x.
	\]
	Oznacza to, że
	\[
	x\ll\bigsqcup D.
	\]
	Na mocy definicji~\ref{def:relacja_way_below} istnieje $d\in D$
	takie, że
	\[
	x\sqsubseteq d.
	\]
	Nie wystarcza to jednak do stwierdzenia, że $d\in\mathord{\Uparrow}x$.
	Korzystamy zatem z własności interpolacji
	twierdzenia~\ref{tw:interpolacja_way_below}. Istnieje $y\in P$ takie,
	że
	\[
	x\ll y\ll\bigsqcup D.
	\]
	Ponieważ
	\[
	y\ll\bigsqcup D,
	\]
	z definicji relacji $\ll$ istnieje $d\in D$ takie, że
	\[
	y\sqsubseteq d.
	\]
	Z twierdzenia~\ref{tw:way_below_porządek} otrzymujemy
	\[
	x\ll y\sqsubseteq d,
	\]
	a więc
	\[
	x\ll d.
	\]
	Stąd
	\[
	d\in D\cap\mathord{\Uparrow}x.
	\]
	Ostatecznie
	\[
	D\cap\mathord{\Uparrow}x\neq\varnothing,
	\]
	co kończy dowód.
\end{proof}

Topologia Scotta pozwala zatem traktować strukturę porządkową dziedziny
jako przestrzeń topologiczną. W szczególności w dziedzinie ciągłej
relacja $\ll$ dostarcza naturalnych zbiorów otwartych, które opisują
otoczenia elementów w sensie przybliżeń porządkowych.

W dalszej części potrzebne będzie przypisanie elementom dziedziny
wartości liczbowych opisujących stopień ich informacji. Taką rolę
pełni \emph{pomiar}. W przypadku dziedzin z przeliczalną bazą pomiar
można skonstruować na podstawie przeliczalnej rodziny zbiorów
Scott-otwartych \cite{Waszkiewicz2001,Waszkiewicz2003}. Prowadzi to
do następnego etapu konstrukcji, w którym struktura topologiczno-
porządkowa zostanie wyposażona w ilościowy opis.

\subsubsection{Pomiary}

Dotychczasowa konstrukcja opierała się wyłącznie na strukturze
porządkowej dziedziny. Relacja $\ll$ pozwala opisywać, kiedy jeden
element stanowi przybliżenie drugiego, natomiast topologia Scotta
opisuje sposób, w jaki te przybliżenia zachowują się względem
skierowanych supremów. Kolejnym krokiem jest przypisanie elementom
dziedziny wartości liczbowych, które pozwolą ilościowo opisać stopień
ich przybliżenia.

W tym celu wykorzystuje się pojęcie \emph{pomiaru}. Idea ta pochodzi
z ilościowej teorii dziedzin rozwijanej przez Martina i została
wykorzystana przez Waszkiewicza w konstrukcji odległości związanej
z pomiarem \cite{Waszkiewicz2001,Waszkiewicz2003}.

Wprowadzając pomiar, wygodnie jest odwrócić naturalny porządek na
zbiorze wartości liczbowych. Wynika to z interpretacji elementów
dziedziny jako informacji: jeżeli
\[
x\sqsubseteq y,
\]
to $y$ zawiera co najmniej tyle informacji co $x$, a zatem jego
wartość pomiaru powinna być nie większa. W szczególności pomiar będzie
funkcją monotoniczną względem porządku odwrotnego na
$[0,\infty)$.

\begin{definicja}\label{def:przestrzen_pomiarowa}
	Niech $P$ będzie zbiorem częściowo uporządkowanym. Dla
	$x\in P$ i $\varepsilon>0$ oraz monotonicznej funkcji
	\[
	\mu\colon P\to[0,\infty)^{\mathrm{op}}
	\]
	definiujemy zbiór elementów $\varepsilon$-bliskich $x$ przez
	\[
	\mu_\varepsilon(x)
	=
	\{y\in P\mid y\sqsubseteq x,\ \mu(y)<\mu(x)+\varepsilon\}.
	\]
\end{definicja}

Zauważmy, że monotoniczność funkcji
\[
\mu\colon P\to[0,\infty)^{\mathrm{op}}
\]
oznacza, że
\[
x\sqsubseteq y
\quad\Longrightarrow\quad
\mu(x)\geq\mu(y)
\]
w zwykłym porządku liczb rzeczywistych. Wartość $\mu(x)$ można więc
interpretować jako ilość informacji, która pozostaje do określenia,
aby osiągnąć element $x$. Im bardziej określony jest element, tym
mniejsza może być jego wartość pomiaru.

Zbiór $\mu_\varepsilon(x)$ zawiera elementy znajdujące się poniżej
$x$, których wartość pomiaru różni się od $\mu(x)$ o mniej niż
$\varepsilon$. Jeżeli $\mu$ dobrze odzwierciedla strukturę dziedziny,
zbiory te powinny dostarczać lokalnego opisu topologii Scotta.

\begin{definicja}\label{def:pomiar_indukujacy_topologie}
	Niech $P$ będzie zupełnym zbiorem częściowo uporządkowanym,
	a $\sigma(P)$ jego topologią Scotta. Monotoniczną funkcję
	\[
	\mu\colon P\to[0,\infty)^{\mathrm{op}}
	\]
	nazywamy \emph{indukującą topologię Scotta} na podzbiorze
	$X\subseteq P$, jeżeli dla każdego $U\in\sigma(P)$ i każdego
	$x\in U\cap X$ istnieje $\varepsilon>0$ takie, że
	\[
	\mu_\varepsilon(x)\subseteq U.
	\]
\end{definicja}

Warunek z definicji~\ref{def:pomiar_indukujacy_topologie} oznacza,
że wartości pomiaru pozwalają rozpoznać otoczenia Scotta. Jeżeli
$x$ należy do zbioru Scott-otwartego $U$, to istnieje odpowiednio
małe $\varepsilon>0$ takie, że wszystkie elementy będące poniżej
$x$ i mające wartość pomiaru dostatecznie bliską $\mu(x)$ również
należą do $U$.

W teorii dziedzin interesują nas pomiary, które są zgodne nie tylko
z porządkiem, lecz także z jego zupełnością. Prowadzi to do pojęcia
ciągłości pomiaru.

\begin{definicja}\label{def:pomiar}
	Niech $P$ będzie dziedziną ciągłą. Funkcję
	\[
	\mu\colon P\to[0,\infty)^{\mathrm{op}}
	\]
	nazywamy \emph{pomiarem}, jeżeli jest Scott-ciągła oraz indukuje
	topologię Scotta na swoim jądrze
	\[
	\ker\mu
	=
	\{x\in P\mid\mu(x)=0\}.
	\]
\end{definicja}

Scott-ciągłość pomiaru oznacza w szczególności, że dla każdego
niepustego zbioru skierowanego $D\subseteq P$ zachodzi
\[
\mu\left(\bigsqcup D\right)
=
\bigsqcup_{x\in D}^{[0,\infty)^{\mathrm{op}}}\mu(x).
\]
Ponieważ porządek w $[0,\infty)^{\mathrm{op}}$ jest odwrócony,
powyższa równość odpowiada w zwykłym porządku liczbowym zależności
\[
\mu\left(\bigsqcup D\right)
=
\inf_{x\in D}\mu(x).
\]

W szczególności wartość pomiaru elementu będącego supremum
skierowanej rodziny jest graniczną wartością pomiarów elementów tej
rodziny. Jest to zgodne z interpretacją pomiaru jako ilości
informacji brakującej do pełnego określenia elementu.

Szczególnie istotny jest przypadek dziedzin posiadających
przeliczalną bazę. Wówczas pomiar można skonstruować bezpośrednio
z przeliczalnej rodziny zbiorów Scott-otwartych.

\begin{twierdzenie}\label{tw:pomiar_z_bazy_scotta}
	Niech $P$ będzie dziedziną ciągłą posiadającą przeliczalną bazę
	topologii Scotta
	\[
	\{U_n\mid n\in\mathbb{N}\}.
	\]
	Wtedy funkcja
	\[
	\mu\colon P\to[0,1]
	\]
	określona wzorem
	\[
	\mu(x)
	=
	1-
	\sum_{\substack{n\in\mathbb{N}\\x\in U_n}}
	2^{-(n+1)}
	\]
	jest pomiarem indukującym topologię Scotta.
\end{twierdzenie}

\begin{proof}
	Dla każdego $x\in P$ zbiór indeksów
	\[
	I_x=\{n\in\mathbb{N}\mid x\in U_n\}
	\]
	jest podzbiorem $\mathbb{N}$. Ponieważ
	\[
	\sum_{n=0}^{\infty}2^{-(n+1)}=1,
	\]
	otrzymujemy
	\[
	0\leq
	\sum_{n\in I_x}2^{-(n+1)}
	\leq 1,
	\]
	a zatem
	\[
	0\leq\mu(x)\leq1.
	\]
	Funkcja $\mu$ jest więc dobrze określona.
	
	Załóżmy, że
	\[
	x\sqsubseteq y.
	\]
	Ponieważ każdy zbiór $U_n$ jest górnym zbiorem, z
	$x\in U_n$ wynika
	\[
	y\in U_n.
	\]
	Stąd
	\[
	I_x\subseteq I_y,
	\]
	a więc
	\[
	\sum_{n\in I_x}2^{-(n+1)}
	\leq
	\sum_{n\in I_y}2^{-(n+1)}.
	\]
	W konsekwencji
	\[
	\mu(x)\geq\mu(y).
	\]
	Oznacza to, że $\mu$ jest monotoniczna jako funkcja
	\[
	P\to[0,\infty)^{\mathrm{op}}.
	\]
	
	Pozostaje wykazać, że $\mu$ indukuje topologię Scotta. Niech
	$U\in\sigma(P)$ oraz niech
	\[
	x\in U.
	\]
	Ponieważ $\{U_n\mid n\in\mathbb{N}\}$ jest bazą topologii Scotta,
	istnieje indeks $k\in\mathbb{N}$ taki, że
	\[
	x\in U_k\subseteq U.
	\]
	
	Wybierzmy $\varepsilon>0$ takie, że
	\[
	\varepsilon<2^{-(k+1)}.
	\]
	Pokażemy, że
	\[
	\mu_\varepsilon(x)\subseteq U.
	\]
	Niech $y\in\mu_\varepsilon(x)$. Wtedy
	\[
	y\sqsubseteq x
	\]
	oraz
	\[
	\mu(y)<\mu(x)+\varepsilon.
	\]
	Załóżmy, dla dowodu nie wprost, że
	\[
	y\notin U_k.
	\]
	Ponieważ $x\in U_k$, indeks $k$ należy do $I_x$, natomiast nie
	należy do $I_y$. Zatem
	\[
	\sum_{n\in I_x}2^{-(n+1)}
	-
	\sum_{n\in I_y}2^{-(n+1)}
	\geq
	2^{-(k+1)}.
	\]
	Wynika stąd
	\[
	\mu(y)-\mu(x)
	\geq
	2^{-(k+1)}
	>
	\varepsilon,
	\]
	co przeczy nierówności
	\[
	\mu(y)<\mu(x)+\varepsilon.
	\]
	Zatem
	\[
	y\in U_k.
	\]
	Ponieważ
	\[
	U_k\subseteq U,
	\]
	otrzymujemy
	\[
	y\in U.
	\]
	Stąd
	\[
	\mu_\varepsilon(x)\subseteq U.
	\]
	Wobec dowolności $U$ i $x$ funkcja $\mu$ indukuje topologię Scotta.
\end{proof}

Powyższa konstrukcja pokazuje, że przeliczalna struktura topologiczna
dziedziny może zostać zakodowana za pomocą jednej funkcji o wartościach
rzeczywistych. Każdy zbiór bazowy $U_n$ otrzymuje w konstrukcji
odpowiednią wagę $2^{-(n+1)}$, a wartość $\mu(x)$ zależy od tego,
które z tych zbiorów zawierają element $x$.

Warto zwrócić uwagę na szczególną rolę jądra pomiaru. Jeżeli
\[
x\in\ker\mu,
\]
to
\[
\mu(x)=0.
\]
W konstrukcji metrycznej przedstawionej w następnym podrozdziale
elementy jądra będą zachowywać się jak punkty o zerowej odległości
własnej. Właśnie ta możliwość odróżnia metrykę częściową od zwykłej
metryki, dla której odległość punktu od samego siebie zawsze wynosi
zero.

Pomiar stanowi zatem ilościowy opis struktury dziedziny. Z jednej
strony pozostaje związany z relacją $\ll$ i topologią Scotta, z drugiej
zaś przyjmuje wartości w zbiorze liczbowym. Pozwala to przejść od
czysto porządkowego opisu dziedziny do konstrukcji funkcji
przypominającej odległość. W następnym podrozdziale wykorzystamy
pomiar $\mu$ do zdefiniowania odległości pomiędzy dwoma elementami
dziedziny:
\[
d_\mu(x,y)
=
\inf\{\mu(z)\mid z\ll x,\ z\ll y\}.
\]
Konstrukcja ta jest podstawowym krokiem w przejściu od teorii
pomiarów do metryk częściowych
\cite{Waszkiewicz2001,Waszkiewicz2003}.

\end{comment}

%%%%%%%%%%%%%%%%%%%%%%%%%%%%%%%%%%%%%%%%%%%%%%%%%%%%%%%%%%%%%%%%%%%%%%%%%%%%%%%%%%%%%

	Wprowadzenie metryk częściowych pozwoliło otrzymać alternatywny sposób
	konstrukcji punktów stałych. Wcześniej punkt stały otrzymywaliśmy za
	pomocą struktury porządkowej: ciąg kolejnych przybliżeń był rosnącym
	łańcuchem, a jego kres górny był punktem stałym funkcjonału. W podejściu
	Matthewsa podstawową rolę odgrywa natomiast kontrakcja oraz zupełność
	metryki częściowej.
	
	Oba podejścia wykorzystują zatem iterację funkcji, lecz wymagają
	odmiennych założeń. Warto wskazać tę różnicę, ponieważ pokazuje ona, że
	twierdzenie Matthewsa nie jest po prostu innym sformułowaniem
	twierdzenia Kleenego.
	
	\begin{table}[h]
		\centering
		\begin{tabular}{p{0.42\textwidth}p{0.42\textwidth}}
			\hline
			\textbf{Podejście porządkowe} &
			\textbf{Podejście metryczne} \\
			\hline
			zbiór łańcuchowo zupełny z porządkiem $\sqsubseteq$ &
			zupełna metryka częściowa $p$ \\[2mm]
			
			ciąg rosnący względem $\sqsubseteq$ &
			ciąg Cauchy'ego względem $p$ \\[2mm]
			
			istnienie kresu górnego &
			istnienie granicy \\[2mm]
			
			ciąg iteracji $f^n(\bt)$, start w $\bt$ &
			ciąg iteracji $f^n(x_0)$, dowolne $x_0$ \\[2mm]
			
			ciągłość funkcjonału &
			kontrakcja \\[2mm]
			
			najmniejszy punkt stały (może nie być jedyny) &
			jedyny punkt stały $x^*$, przy czym $p(x^*,x^*)=0$ \\[2mm]
			
			twierdzenie Kleenego &
			twierdzenie Matthewsa \\[2mm]
			\hline
		\end{tabular}
		\caption{Porównanie dwóch podejść do konstrukcji punktu stałego.}
		\label{tab:porownanie}
	\end{table}
	
	Podobieństwo obu konstrukcji (zob.\ Tabela~\ref{tab:porownanie}) jest
	widoczne przede wszystkim w sposobie wyznaczania punktu stałego. W obu
	przypadkach wielokrotnie stosujemy badaną funkcję do pewnego elementu
	początkowego i otrzymujemy ciąg przybliżeń
	\[
	x_0,\quad f(x_0),\quad f^2(x_0),\quad\ldots,
	\]
	którego element graniczny okazuje się punktem stałym. Różnica pojawia
	się już na tym etapie: w twierdzeniu Kleenego iterację trzeba rozpocząć
	od elementu najmniejszego $\bt$ (tylko wtedy otrzymujemy łańcuch
	i najmniejszy punkt stały), natomiast w twierdzeniu Matthewsa punkt
	początkowy $x_0$ może być dowolny.
	
	Różny jest również mechanizm gwarantujący istnienie elementu
	granicznego. W podejściu porządkowym wykorzystujemy ciągłość funkcji
	oraz istnienie kresu górnego rosnącego łańcucha. W podejściu metrycznym
	wykorzystujemy kontrakcję, która zapewnia własność Cauchy'ego ciągu
	iteracji, oraz zupełność metryki częściowej, która gwarantuje istnienie
	jego granicy.
	
	Najistotniejsza jest jednak różnica w otrzymanym wyniku. Twierdzenie
	Kleenego wskazuje \emph{najmniejszy} punkt stały, przy czym punktów
	stałych może być wiele --- jak w przypadku pętli
	$\kw{while}\ \kw{true}\ \kw{do}\ \kw{skip}$, dla której każda funkcja
	częściowa jest punktem stałym, a właściwą denotacją jest
	$\varnothing$. Twierdzenie Matthewsa daje natomiast punkt stały
	\emph{jedyny}, i to taki, że $p(x^*,x^*)=0$, czyli element całkowicie
	określony. W szczególności metodą tą nie otrzymamy denotacji, która
	jest gdzieś nieokreślona, np.\ funkcji pustej opisującej pętlę
	nieskończoną.
	
	Warunki ciągłości i kontrakcji nie są przy tym porównywalne wprost,
	ponieważ odnoszą się do różnych struktur: pierwszy do porządku,
	drugi do metryki częściowej. Twierdzenie Matthewsa nie zastępuje więc
	twierdzenia Kleenego w ogólnym przypadku. Stanowi natomiast inne
	narzędzie konstrukcji punktu stałego, szczególnie użyteczne wtedy,
	gdy działanie funkcji można kontrolować ilościowo za pomocą metryki
	częściowej, a poszukiwany obiekt ma być całkowicie określony.

%========================= Bibliografia =========================
\phantomsection
\addcontentsline{toc}{section}{\refname}
\begin{thebibliography}{99}

\bibitem{gordon1979}
  M.~J.~C.~Gordon, \emph{The Denotational Description of Programming Languages: An Introduction}. Springer-Verlag, 1979.
  
\bibitem{matthews1994}
  S.\,G.~Matthews, \emph{Partial Metric Topology}. W: \emph{Papers on General Topology and Applications} (Eighth Summer Conference at Queens College), Annals of the New York Academy of Sciences, vol.~728, 1994, s.~183--197.

\bibitem{cs152denot}
  \emph{Denotational Semantics}, CS~152: Programming Languages, wykład~6, Harvard School of Engineering and Applied Sciences, 9~lutego 2023.

\bibitem{cs4110denot}
  \emph{Denotational Semantics}, CS~4110: Programming Languages and Logics, wykład~8, Cornell University, jesień 2012.
  
% Pozycje cytowane wyłącznie w wykomentowanym fragmencie o teorii dziedzin:
% \bibitem{Waszkiewicz2001}
% P.~Waszkiewicz,
% \emph{Distance and Measurement in Domain Theory},
% Electronic Notes in Theoretical Computer Science,
% 45 (2001), 448--462.
% DOI: 10.1016/S1571-0661(04)80975-7.

% \bibitem{Waszkiewicz2003}
% P.~Waszkiewicz,
% \emph{Quantitative Continuous Domains},
% Applied Categorical Structures,
% 11(1) (2003), 41--67.
% DOI: 10.1023/A:1023012924892.

\bibitem{naur1960}
P. Naur (red.),
\emph{Report on the Algorithmic Language ALGOL 60},
Communications of the ACM, 3(5), 1960, s.~299--314.

\bibitem{kleene1952}
S.\,C.~Kleene,
\emph{Introduction to Metamathematics},
North-Holland, Amsterdam; D.~Van Nostrand, New York, 1952.

\end{thebibliography}

\end{document}
